\documentclass[11pt]{amsart}
\usepackage{amsmath,amssymb}
\setlength{\emergencystretch}{3em}
\tolerance=600
\usepackage{tikz-cd}
\usepackage{tikz}
\usepackage[hidelinks]{hyperref}

\newtheorem{thm}{Theorem}[section]
\newtheorem{lem}[thm]{Lemma}
\newtheorem{cor}[thm]{Corollary}
\newtheorem{prop}[thm]{Proposition}
\theoremstyle{definition}
\newtheorem{defi}[thm]{Definition}
\newtheorem{rem}[thm]{Remark}
\newtheorem{ex}[thm]{Example}
\newtheorem{que}[thm]{Question}

\newcommand{\Set}{\mathbf{Set}}
\newcommand{\Fin}{\mathbf{FinSet}}
\newcommand{\Sub}{\mathrm{Sub}}
\newcommand{\Mod}{\mathrm{Mod}}
\newcommand{\Sent}{\mathrm{Sent}}
\newcommand{\Lan}{\mathrm{Lan}}
\newcommand{\Aut}{\mathrm{Aut}}
\newcommand{\Max}{\mathrm{Max}}
\newcommand{\PP}{\mathsf P}
\newcommand{\HH}{\mathsf H}
\newcommand{\cX}{\mathcal X}
\newcommand{\cC}{\mathcal C}
\newcommand{\cE}{\mathcal E}
\newcommand{\cH}{\mathcal H}
\newcommand{\cM}{\mathcal M}
\newcommand{\cS}{\mathcal S}
\newcommand{\bT}{\mathbb T}
\newcommand{\id}{\mathrm{id}}

\begin{document}

\title[Finite sorts and the logics proving linearity]{Intermediate logics
and the imaginaries of a theory with finite sorts: the logics proving
linearity}
\author{Ryota Shibusawa}
\thanks{Corresponding author. E-mail: \texttt{r-shibusawa@daiichi-koudai.ac.jp}}
\address{Daiichi Institute of Technology}
\email{r-shibusawa@daiichi-koudai.ac.jp}
\keywords{intermediate logic, G\"odel--Dummett logic, weak excluded middle, Heyting category, pretopos completion, imaginaries, Kripke model, EI-category, profinite group, coherent theory}
\subjclass[2020]{03G30 (primary), 03B55, 03C13, 18F10, 20E18 (secondary)}

\begin{abstract}
For an intermediate logic $L$ and a coherent theory, reading the
theory in first-order logic over $L$ and adding the imaginaries need
not commute. For theories whose models are finite on every sort the
question was open for the logics proving linearity with a
three-element chain frame, unless the theory has finitely many
sentences. We settle it for all these logics, including
G\"odel--Dummett logic: the two operations commute if and only if
every endomorphism of a model is invertible and no leg tree is an
obstruction at a chain of non-invertible homomorphisms of models, no
longer than the chain frames of the logic. A leg tree is a finite tree of such chains for a finite
language; pulled back to the larger finite languages, with the
homomorphisms which become invertible after restriction as new
branches, it defines two closed subgroups of the automorphism group of
a model, and it is an obstruction if they differ. The condition passes
to logics with shorter chain frames, but for a limit
of a tower of finite categories it holds for the logics of all finite
chains and fails for G\"odel--Dummett logic. The two subgroups are
intersections over the solutions of finite level of a system of
implications between equivalence relations;
if the chain frames are bounded, they are given by a recursion with
open equivalence relations, and the criterion can be stated without
leg trees.
\end{abstract}

\maketitle

\section{Introduction}

For an intermediate logic $L$ and a small coherent category $\cC$ let
$\HH_L\cC$ be the syntactic category of the internal coherent theory of
$\cC$ in first-order logic over $L$, and $\PP$ the pretopos completion,
which adds the imaginaries. There is a full and faithful comparison
$c^L_\cC\colon\HH_L\PP\cC\to\PP\HH_L\cC$, and it is an equivalence if
and only if $\HH_L\PP\cC$ is a pretopos \cite[Lemma~2.2]{Interm}. We
say that \emph{the translation into $L$ commutes with the imaginaries
of $\cC$} if $c^L_\cC$ is an equivalence.

A \emph{model} of $\cC$ is a coherent functor $M\colon\cC\to\Set$, a
\emph{homomorphism} is a natural transformation, and $\cC$ \emph{has
finite sorts} if the set $M(A)$ is finite for every model $M$ and
every object $A$. For these $\cC$ the comparison was determined in
\cite{Finsorts} in terms of the category $\Mod(\cC)$ of models, for
all $L$ except those proving the linearity axiom $(p\to q)\vee(q\to
p)$ with $n_L\ge3$; here $n_L\in\{1,2,\dots,\infty\}$ is the supremum
of the numbers of elements of the finite chains which are frames of
$L$. In that case it is necessary that $\Mod(\cC)$ is an
\emph{EI-category} (every endomorphism is invertible), and this is not
sufficient \cite{Strings}. A \emph{sentence} is a subobject of the
terminal object of $\cC$. For theories with finitely many sentences
the case $n_L\ge3$ was settled in \cite{Proftrees}. This paper treats
all theories with finite sorts and all logics proving linearity.

The models of the pretopos completion are those of $\cC$, so let $\cC$
be a pretopos with finite sorts. A \emph{finite language}
$\Lambda=(S,R)$ consists of a finite set $S$ of objects of $\cC$ and a
finite set $R$ of subobjects of finite products of objects in $S$; the
theory $\bT_\Lambda$ in this language has as axioms the coherent
sequents valid in $\cC$. Its category of models is equivalent to a
finite category $\cX_\Lambda$, and for $\Lambda\subseteq\Lambda'$ the
restriction of structures is a functor
$\cX_{\Lambda'}\to\cX_\Lambda$. A \emph{string} is a sequence
$M_0\to\dots\to M_r$ of non-invertible homomorphisms of models with at
most $n_L$ models. If $\Mod(\cC)$ is an EI-category, then so is
$\cX_\Lambda$ for cofinally many $\Lambda$ \cite{Finsorts}, and we
consider these $\Lambda$ only. A string $s'$ of $\cX_{\Lambda'}$
restricts to a string $s'|_\Lambda$ of $\cX_\Lambda$, after the
homomorphisms whose restrictions are invertible have been deleted; in
the same way a string $\tilde s$ of models of $\cC$ has restrictions
$\tilde s|_\Lambda$, and $(\tilde s|_{\Lambda'})|_\Lambda=\tilde
s|_\Lambda$ for $\Lambda\subseteq\Lambda'$. The automorphism group
$G_{\tilde s}$ of the first model of $\tilde s$ is profinite, and it
contains the closed subgroup $H_{\tilde s}$ of the automorphisms of
the first model which extend to an automorphism of the string $\tilde
s$. A \emph{leg tree}
$T$ of $\cX_\Lambda$ at $\tilde s|_\Lambda$ \cite[Def.~6.14]{Strings}
is a finite tree of strings of $\cX_\Lambda$ which extend $\tilde
s|_\Lambda$ to the left. It has a \emph{pullback} to every
$\Lambda'\supseteq\Lambda$ (Section~\ref{sec:pull}), a leg tree at
$\tilde s|_{\Lambda'}$: a branch is replaced by all branches of
$\cX_{\Lambda'}$ over it, and new branches are added at every node for
the arrows of $\cX_{\Lambda'}$ into the first model of its string
which are not invertible but have an invertible restriction to
$\Lambda$. We call these the \emph{collapsing children}. Forming joins
of equivalence relations in two ways, on all homomorphisms and on
those which are homomorphisms of strings, at every $\Lambda'$, and
passing to the limit, we obtain closed subgroups
\[
R_{(\Lambda,T)}\le H_{\tilde s}\cap S_{(\Lambda,T)}
\]
of $G_{\tilde s}$ (Section~\ref{sec:inverse}); the pair $(\Lambda,T)$
is an \emph{obstruction at $\tilde s$} if they are different.

\begin{thm}\label{thm:A}
Let $\cC$ be a small coherent category with finite sorts and $L$ a
non-classical intermediate logic proving linearity. Then $c^L_\cC$ is
an equivalence if and only if $\Mod(\cC)$ is an EI-category and no leg
tree of a finite language of $\PP\cC$ is an obstruction at a string of
models.
\end{thm}

The strings have at most $n_L$ models; for $n_L=\infty$ they have
arbitrary finite length, and the depth of the pulled back trees is not
bounded. The models of $\HH_L\PP\cC$ are described by more general
objects. A \emph{thread} is a family $\sigma=(s_\Lambda)_\Lambda$ of
strings $s_\Lambda$ of $\cX_\Lambda$ with
$s_{\Lambda'}|_\Lambda=s_\Lambda$ for $\Lambda\subseteq\Lambda'$. The
restrictions of a string of models form a thread; if $n_L$ is finite,
or if the chains of prime filters of sentences are bounded, every
thread arises in this way (Remark~\ref{rem:strings}), but in general
the lengths of the $s_\Lambda$ need not be bounded. If
$\Mod(\cC)$ is an EI-category, every model of $\HH_L\PP\cC$ is
isomorphic to the model $\hat G_\sigma$ of a thread $\sigma$, and its
automorphism group is the group $H_\sigma$ of the automorphisms of the
first model whose restriction to every $\Lambda$ extends to an
automorphism of $s_\Lambda$ (Proposition~\ref{prop:models}).
Obstructions are defined at threads in the same way, and the criterion
is first proved for threads (Theorem~\ref{thm:general}); an
obstruction at a thread gives one at a string of models
(Theorem~\ref{thm:bounded}). For theories with
finitely many sentences Theorem~\ref{thm:A} is
\cite[Thm.~1.1]{Proftrees}. There the theory is approximated by its
restrictions to finite sets of sorts, which keep all sentences; no
homomorphism becomes invertible and there are no collapsing children.
A finite language has finitely many sentences up to equivalence, and
non-isomorphic models can have isomorphic restrictions; this was the
difficulty named in \cite[Question~10.1]{Proftrees}, which
Theorem~\ref{thm:A} answers. With \cite{Finsorts} for the other
logics, it is now determined for every intermediate logic and every
coherent theory with finite sorts whether the translation commutes
with the imaginaries; for the logics of Theorem~\ref{thm:A} the
condition is stated through the finite languages
(Section~\ref{sec:questions}). If $n_L$ is finite, the models of
$\HH_L\cC$ have the following description. A model of $\HH_L\cC$ is a
coherent functor to sets; a string $s$ is a Kripke model of $\cC$ over
a chain which is a frame of $L$, and the evaluation at its first node
is a model $\hat G_s$ of $\HH_L\cC$.

\begin{thm}\label{thm:B}
Let $\cC$ be a pretopos with finite sorts such that $\Mod(\cC)$ is an
EI-category, and $L$ an intermediate logic proving linearity with
$n_L$ finite. Then every model of $\HH_L\cC$ is isomorphic to $\hat
G_s$ for a string $s$, the homomorphisms $\hat G_s\to\hat G_{s'}$ are
the composites of an initial part of $s$ with the first component of
an isomorphism from the remaining final segment of $s$ onto $s'$, and
$U\mapsto(s\mapsto\hat G_s(U))$ is a conservative Heyting functor from
$\HH_L\cC$ to the category of functors from the strings to sets.
\end{thm}

For a finite category of models this is \cite[Prop.~4.2]{Strings}.
It holds also for $n_L=\infty$ if the chains of prime filters of
sentences are bounded, and then $\HH_L\cC$ is the free category for
the logic of a finite chain (Proposition~\ref{prop:Lm}). The next result
needs neither finite sorts nor linearity. For a pretopos
$\cC$ and a prime filter $Q$ of the lattice $\Sub(1)$ of sentences let
$\cC_{(Q)}$ be the pretopos of the theory of $\cC$ with the additional
axioms $s'\vdash\bot$ for the sentences $s'\notin Q$; its models are
the models of $\cC$ whose true sentences lie in $Q$.

\begin{thm}\label{thm:C}
Let $\cC$ be a pretopos and $L$ an intermediate logic proving the weak
excluded middle $\neg p\vee\neg\neg p$. Then $c^L_\cC$ is an
equivalence if and only if $c^L_{\cC_{(Q)}}$ is an equivalence for
every prime filter $Q$ of $\Sub(1)$.
\end{thm}

If every prime filter of $\Sub(1)$ contains only finitely many prime
filters, this gives a criterion without finite languages
(Corollary~\ref{cor:locfin}). The first example with collapsing children is the following. Let $G$
be a profinite abelian group, $H$ a closed subgroup and
$(S_k)_{k\in\mathbb N}$ non-zero closed subgroups \emph{converging to
$0$}: every open subgroup contains all but finitely many $S_k$. Let
$\cC(G;H;S_\bullet)$ be the pretopos of the continuous functors to
finite sets on the category
\[
\begin{tikzcd}[column sep=large]
0_k\ar[r,"G/S_k"]\ar[rr,bend right=22,"G/(H+S_k)"']&1\ar[r,"G/H"]
\ar[loop above,"G"]&2
\end{tikzcd}\qquad(k\in\mathbb N)
\]
with composition the addition of $G$; precisely, it is a directed
colimit of pretoposes $\Fin^{\cX_N}$ over the open subgroups $N$ of
$G$, and in $\cX_N$ the objects $0_k$ with $S_k\subseteq N$ are
identified with $1$ (Section~\ref{sec:family}).

\begin{thm}\label{thm:D}
Let $L$ prove linearity with $n_L\ge3$. Then $c^L_\cC$ is an
equivalence for $\cC=\cC(G;H;S_\bullet)$ if and only if
\begin{equation}\label{eq:cofintro}
H\cap\overline{\sum_{k\in K}S_k}=\overline{\sum_{k\in K}(H\cap S_k)}
\qquad\text{for every cofinite }K\subseteq\mathbb N,
\end{equation}
the bars denoting the closed subgroups generated.
\end{thm}

For $G=(\mathbb Z/2)^{\mathbb N}$, $S_k$ generated by the $k$-th unit
vector and $H$ generated by $(1,1,1,\dots)$, the condition holds for
all finite sets of indices and fails for $K=\mathbb N$
(Example~\ref{ex:conv}): the obstruction uses the legs from all the
models $0_k$ together, and at each finite stage, for all but finitely
many $k$ they have become automorphisms of the model $1$.

The condition of Theorem~\ref{thm:A} depends on $L$ through $n_L$. For
$m\ge1$ let $L_m$ be the logic of the chain with $m$ elements, so that
$n_{L_m}=m$.

\begin{thm}\label{thm:E}
\begin{enumerate}
\item[(a)] Let $\cC$ be a coherent category with finite sorts, $L$ a
non-classical logic proving linearity and $L'$ a logic proving
linearity with $n_{L'}\le n_L$. If $c^L_\cC$ is an equivalence, so is
$c^{L'}_\cC$.
\item[(b)] For every $m\ge3$ there is a finite EI-category $\cX$ such
that, for $L$ proving linearity, $c^L_\cC$ is an equivalence for
$\cC=\Fin^\cX$ if and only if $n_L\le m$.
\item[(c)] There is a pretopos $\cC$ with finite sorts such that, for
$L$ proving linearity, $c^L_\cC$ is an equivalence if and only if
$n_L$ is finite. So the translation commutes with the imaginaries of
$\cC$ for all the logics $L_m$, and not for G\"odel--Dummett logic.
\end{enumerate}
\end{thm}

For $m=3$ the category of (b) is the example of
\cite[Ex.~6.19]{Strings}; the others generalise it
(Section~\ref{sec:depth}). The pretopos of (c) is the colimit of the
$\Fin^{\cX}$ along a tower of these categories. Its obstruction is at
a string with one model; the pulled back trees of
Theorem~\ref{thm:A} become deeper with the finite language, and by
(c) their depth cannot be bounded in terms of the logics $L_m$.

The groups $S_{(\Lambda,T)}$ and $R_{(\Lambda,T)}$ are defined through
the finite languages $\Lambda'\supseteq\Lambda$, and a homomorphism of
models of $\bT_{\Lambda'}$ need not be the restriction of a
homomorphism of models of $\cC$. They can be computed in $\Mod(\cC)$.
The pullback of $T$ to $\Mod(\cC)$ has as nodes strings of models
which extend $\tilde s$ to the left, and a node can have infinitely
many children. A \emph{solution} is a family of equivalence relations
on the sets of homomorphisms from the first models of these strings to
the first model of $\tilde s$, one for each node, such that two
homomorphisms are related at a node whenever their composites with a
child are related at the child; it is \emph{of finite level} if it
depends only on the restrictions to a finite language
(Section~\ref{sec:intrinsic}; for $n_L=\infty$ the nodes are
compatible families of strings of the finite languages).

\begin{thm}\label{thm:F}
Let $\cC$ be a pretopos with finite sorts such that $\Mod(\cC)$ is an
EI-category, $L$ a logic proving linearity, $\tilde s$ a string of
models and $T$ a leg tree of a finite language $\Lambda$ at $\tilde
s|_\Lambda$.
\begin{enumerate}
\item[(a)] $S_{(\Lambda,T)}$ is the set of the automorphisms of the
first model of $\tilde s$ which are related to the identity, at the
root, in every solution of finite level; and $R_{(\Lambda,T)}$ is
described in the same way with the homomorphisms of strings.
\item[(b)] If $n_L$ is finite, the relations defining the groups are
given by a recursion over the pulled back tree in $\Mod(\cC)$: the
relation at a node is the intersection of the open equivalence
relations containing the relations induced by its children. They form
the least solution which is closed, as a subset of the space of the
pairs of homomorphisms at all nodes, and consists of intersections of
open equivalence relations.
\item[(c)] If $n_L$ is finite, $\HH_L\cC$ is a pretopos if and only
if for every finite language $\Lambda$, every object $W$ of the
pretopos completion of $\HH_L\cC_\Lambda$, where $\cC_\Lambda$ is the
pretopos of $\bT_\Lambda$, every string $\tilde s$ of models, $a\in
W(\tilde s)$ and every automorphism $\Theta$ of the first model which
extends to an automorphism of the string $\tilde s$: if
$\id\mathrel{E^W_{(\tilde s,a)}}\Theta$, then $\Theta a=a$. Here $W(\tilde s)$ is the value at $W$ of the model of
$\HH_L\cC_\Lambda$ given by the restriction of $\tilde s$ to
$\Lambda$, extended to the pretopos completion, and the relations $E^W_{(\tilde t,y)}$
on the homomorphisms from the first model of $\tilde t$ to that of
$\tilde s$ are defined by the recursion of \emph{(b)} over the pairs
of a string $\tilde t$ of models and $y\in W(\tilde t)$, the children
of $(\tilde t,y)$ being the one-step homomorphisms $\tilde u\to\tilde
t$ with an element of $W(\tilde u)$ over $y$.
\end{enumerate}
For limits of inverse systems of finite EI-categories, the least
solution does not give the groups if $n_L=\infty$, and in \emph{(b)}
the open equivalence relations cannot be replaced by the closed ones
\textup{(Examples~\ref{ex:towersol} and~\ref{ex:interval})}.
\end{thm}

For theories with finitely many sentences (b) is
\cite[Prop.~8.3]{Proftrees}. In general a sequence of children which
are not isomorphisms can converge to an isomorphism, and the space of
the children of a node is not compact.

\subsection*{The method}
If $\Mod(\cC)$ is an EI-category, $\cC$ is the directed colimit of the
pretoposes $\Fin^{\cX_\Lambda}$ over the finite languages $\Lambda$
for which $\cX_\Lambda$ is an EI-category, and $\Mod(\cC)$ is the
limit of the $\cX_\Lambda$ \cite{Finsorts}.
\[
\begin{tikzcd}[column sep=large,row sep=small]
\Mod(\cC)\simeq\varprojlim\cX_\Lambda\ar[r]&\cX_{\Lambda'}\ar[r,"q"]&
\cX_\Lambda\\
\cC=\varinjlim\Fin^{\cX_\Lambda}&\Fin^{\cX_{\Lambda'}}\ar[l]&
\Fin^{\cX_\Lambda}\ar[l,"q^*"']
\end{tikzcd}
\]
By \cite{Strings}, $\PP\HH_L\Fin^{\cX}$ is the category of the functors
$W$ from the category $\cS_n(\cX)$ of the strings of $\cX$ to finite
sets, and $\HH_L\Fin^\cX$ consists of those $W$ for which a unit map
is injective. We work with an arbitrary inverse system of finite
EI-categories and arbitrary functors $q$. Such a functor acts on
strings by deleting the arrows which become invertible
(Lemma~\ref{lem:collapse}), and the comparison is an equivalence iff
every $W$ at a stage has an injective unit at some later stage
(Lemma~\ref{lem:eventual}). The kernel of the unit of $W$ is
controlled by the largest leg tree realised at each element of $W$
\cite{Proftrees}. The new point is that the largest tree for the
pullback of $W$ along $q$ is the pullback of the largest tree for $W$
\emph{with the collapsing children} (Lemma~\ref{lem:pull}). A
compactness argument on the finite sets of the elements at which a
unit fails to be injective gives the criterion
(Theorem~\ref{thm:general}), in terms of the compatible families of
strings of the stages, the threads. These are the strings of the limit
category if their lengths are bounded (Lemma~\ref{lem:threads}), and
in general they describe the models of $\HH_L\cC$
(Proposition~\ref{prop:models}). The relation which a leg tree
defines on all arrows depends only on the first object of the string
at which it is formed; so omitting objects from the strings of a
thread gives a string of the limit category at which a failure of the
criterion persists (Theorem~\ref{thm:bounded}). For
Theorem~\ref{thm:E}, the strings with at most $m$ objects form a
subterminal object of the category for $n_L>m$, and the category for
$L_m$ is a colimit of slices (Theorem~\ref{thm:mono}); the examples
are categories with one object with two automorphisms above a poset,
for which the criterion is a recursion on equivalence relations
(Lemma~\ref{lem:cone}). For Theorem~\ref{thm:F}, a solution of finite level
is pulled back to a late stage, where it contains the relations of
the pulled back tree because these are the least solution there
(Theorem~\ref{thm:level}); only finite configurations of a fixed
shape have to be lifted to the limit. For bounded strings a closed
solution has neighbourhoods which are solutions
(Theorem~\ref{thm:closed}), and the relations of the recursion form a
closed solution (Theorem~\ref{thm:rec}). For the largest leg trees the
tree in the category of models is the tree of the elements of $W$
(Theorem~\ref{thm:elements}).

\subsection*{What is new and what is known}
Known: the classification of \cite{Finsorts}, with the necessity of
the EI property; the criterion for a finite category of models
\cite[Thm.~1.1, Thm.~6.15]{Strings}; the criterion for finitely many
sentences \cite[Thm.~1.1, Thm.~8.2]{Proftrees}, which is the case of
Theorem~\ref{thm:general} in which the functors reflect isomorphisms
(Remark~\ref{rem:special}); Theorem~\ref{thm:B} for a finite category
of models \cite[Prop.~4.2]{Strings}, and for propositional theories,
where it amounts to the description of the free G\"odel algebras over
distributive lattices by chains of prime filters \cite{AGM08,Carai}; the implication from $\cC$ to the
$\cC_{(Q)}$ in Theorem~\ref{thm:C} \cite[Thm.~1.3]{Proftrees}; the
case $m=3$ of Theorem~\ref{thm:E}(b) \cite[Ex.~6.19]{Strings}, and
that for a finite category of models the comparison is an equivalence
if $n_L\le2$ \cite[Cor.~6.3]{Strings}. Theorem~\ref{thm:F}(b) is known for the
towers of \cite{Proftrees}, whose functors reflect isomorphisms
\cite[Prop.~8.3]{Proftrees}. New, to
the author's knowledge: Theorems~\ref{thm:A}--\ref{thm:F} beyond these
cases.

\section{The setting and the results used}\label{sec:setting}

Conventions are those of \cite{Finsorts} and \cite{Proftrees}. All
coherent categories are small. A \emph{pretopos} is a coherent
category with disjoint finite coproducts and effective quotients of
equivalence relations; $\Fin$ is the category of finite sets. An
\emph{$L$-category} is a Heyting category whose lattices of subobjects
are Heyting algebras in which the formulas of $L$ are valid;
$\eta\colon\cC\to\HH_L\cC$ is the free $L$-category on $\cC$: for an
$L$-category $\cE$, composition with $\eta$ is an equivalence from the
groupoid of the Heyting functors $\HH_L\cC\to\cE$ and natural
isomorphisms to that of the coherent functors $\cC\to\cE$. We omit
$\eta$. Every object of $\HH_L\cC$ is a subobject $U\le A$ of an
object $A$ of $\cC$, and these subobjects are obtained from those in
$\cC$ by finite meets and joins, $\Rightarrow$, and $f^*$,
$\exists_f$, $\forall_f$ for the morphisms $f$ of $\cC$. The models of
$\PP\cC$ are those of $\cC$, and $c^L_\cC$ is an equivalence iff
$c^L_{\PP\cC}$ is. So \emph{from now on $\cC$ is a pretopos}; then
$c^L_\cC$ is an equivalence iff $\HH_L\cC$ is a pretopos
\cite[Lemma~2.2]{Interm}. For a model $M$ let $\Sent(M)$ be the prime
filter of the sentences true in $M$, and $A(M)=M(A)$. A finite chain
is a frame of $L$ proving linearity iff it has at most $n_L$ elements.
Directed colimits of categories are taken along functors compatible up
to coherent isomorphisms; a directed colimit of pretoposes and
coherent functors is a pretopos. If $\cC$ has finite sorts, every
model is isomorphic to a \emph{standard} one, whose values at the
objects are sets $\{1,\dots,e\}$; the standard models form a small
category equivalent to $\Mod(\cC)$, and we take strings of standard
models. We use:

\begin{itemize}
\item[(R1)] (Completeness.) $U\le V$ for subobjects of $A$ iff
$U(M)\subseteq V(M)$ for all models $M$; every prime filter of
$\Sub(1)$ is $\Sent(M)$ for some $M$.
\item[(R2)] (Types, \cite[Lemma~4.3]{Finsorts}.) Let $N$, $N'$ be
models of a coherent category which are finite on every object, and
$(\alpha_i)_{i\in I}$, $(\alpha'_i)_{i\in I}$ families of elements
$\alpha_i\in A_i(N)$, $\alpha'_i\in A_i(N')$ such that for all
$i_1,\dots,i_r\in I$, $r\ge0$, and all $S\le A_{i_1}\times\dots\times
A_{i_r}$: $(\alpha_{i_1},\dots,\alpha_{i_r})\in S(N)$ implies
$(\alpha'_{i_1},\dots,\alpha'_{i_r})\in S(N')$. Then there is a
homomorphism $\lambda\colon N\to N'$ with $\lambda(\alpha_i)=
\alpha'_i$ for all $i$.
\item[(R3)] (\cite[Lemma~5.3 and its proof]{Finsorts}.) If a pretopos
$\cC$ is the colimit of a directed system of pretoposes $\cC_i$ and
coherent functors, then $\HH_L\cC$ is the colimit of the $\HH_L\cC_i$.
\item[(R4)] (\cite[Remark after Thm.~5.5]{Finsorts}.) If $\cC$ has
finite sorts, $L$ is not classical and $c^L_\cC$ is an equivalence,
then $\Mod(\cC)$ is an EI-category.
\item[(R5)] (\cite[Prop.~3.1]{Proftrees}.) Let $\cC$ have finite sorts
and $\Mod(\cC)$ be an EI-category. There is a homomorphism $M\to N$
iff $\Sent(M)\subseteq\Sent(N)$, and it is invertible iff
$\Sent(N)\subseteq\Sent(M)$.
\end{itemize}

\subsection*{Models of the translated theory}
A \emph{model} of $\HH_L\cC$ is a coherent functor $\hat G\colon
\HH_L\cC\to\Set$, not necessarily Heyting; it is \emph{over} its
restriction $M$ to $\cC$, and $\hat G(U)\subseteq A(M)$ for $U\le A$.
A homomorphism $\hat G\to\hat G'$ of models over $M$, $M'$ is a
homomorphism $\varphi\colon M\to M'$ with $\varphi(\hat G(U))\subseteq
\hat G'(U)$ for all $U$ \cite[\S2]{Finsorts}. If $\cC$ has finite
sorts, the models of $\HH_L\cC$ are finite on every object. A sentence
$s$ of $\cC$ is true in $\hat G$ iff it is true in $M$.
\begin{itemize}
\item[(R6)] (Truth lemma, \cite[Lemma~2.1]{Finsorts}.) Let $\hat G$ be
a model of $\HH_L\cC$ over $M$, $f\colon B\to A$ in $\cC$, $U,V\le B$
in $\HH_L\cC$, $\alpha\in A(M)$, $\beta\in B(M)$. Then $\alpha\in\hat
G(\forall_fU)$ iff for every homomorphism $\varphi\colon\hat G\to\hat
G'$, over $M\to M'$, and every $\beta'\in B(M')$ with $f(\beta')=
\varphi(\alpha)$ one has $\beta'\in\hat G'(U)$; and $\beta\in\hat
G(U\Rightarrow V)$ iff for every homomorphism $\varphi\colon\hat G\to
\hat G'$, $\varphi(\beta)\in\hat G'(U)$ implies $\varphi(\beta)\in\hat
G'(V)$.
\item[(R7)] (Kripke models, \cite[Lemma~2.3]{Interm}.) Let
$c=(M_0\xrightarrow{a_1}M_1\to\dots\xrightarrow{a_r}M_r)$ be a
sequence of homomorphisms of models of $\cC$ with $r+1\le n_L$. The
functor $j\mapsto M_j$ on the chain $0<\dots<r$ extends, uniquely up
to isomorphism, to a Heyting functor from $\HH_L\cC$ to the functors
from the chain to sets; its value at $U\le A$ and the node $j$ is the
subset $\|U\|^c_j\subseteq A(M_j)$ given by the forcing clauses:
$U(M_j)$ for $U$ in $\cC$; nodewise for finite meets and joins, $f^*$
and $\exists_f$; $\beta\in\|U\Rightarrow V\|_j$ iff for all $j'\ge j$,
$a_{j'}\cdots a_{j+1}(\beta)\in\|U\|_{j'}$ implies $a_{j'}\cdots
a_{j+1}(\beta)\in\|V\|_{j'}$; $\alpha\in\|\forall_fU\|_j$ iff for all
$j'\ge j$ and $\beta'$ with $f(\beta')=a_{j'}\cdots a_{j+1}(\alpha)$,
$\beta'\in\|U\|_{j'}$. So $U\mapsto\|U\|^c_j$ is a model of $\HH_L\cC$
over $M_j$, and the $a_j$ are homomorphisms between these models.
\end{itemize}
A \emph{string} is a sequence $s$ as in (R7) with all $a_j$
non-invertible; $\hat G_s$ is the model $U\mapsto\|U\|^s_0$, its
\emph{final segments} are the strings $s_{\ge j}=(M_j\to\dots\to
M_r)$, so that $\|U\|^s_j=\hat G_{s_{\ge j}}(U)$, and an
\emph{isomorphism of strings} is a family of isomorphisms of the
models commuting with the arrows.

\subsection*{Finite categories of models}
Let $\cX$ be a finite EI-category, $L$ a logic proving linearity and
$n=n_L$. In an EI-category a composite is invertible only if its
factors are. The models of $\Fin^\cX$ are the evaluations at the
objects and the homomorphisms are the arrows of $\cX$; so the strings
are the sequences $s=(m_0\to\dots\to m_r)$ of non-invertible arrows
with at most $n$ objects, and their number of objects is at most the
number of isomorphism classes of objects of $\cX$. The category
$\cS_n(\cX)$ has the strings as objects, an arrow $s\to t$ being an
arrow $g\colon s_0\to t_0$ of $\cX$ of the form $\iota_0p$ with $p$
the composite $s_0\to s_j$ of arrows of $s$ and $\iota$ an isomorphism
of strings from $s_{\ge j}$ onto $t$; the index $j$ is determined by
$t$. Let $\pi\colon\cS_n(\cX)\to\cX$ be the first object. An arrow of
$\cS_n(\cX)$ between strings with the same number of objects is the
first component of an isomorphism of strings; we call it an
\emph{isomorphism of strings} too, and it is invertible in
$\cS_n(\cX)$.

\begin{itemize}
\item[(R8)] (\cite[Prop.~4.2]{Strings}.) Every model of
$\HH_L\Fin^\cX$ is isomorphic to $\hat G_s$ for a string $s$, and the
homomorphisms $\hat G_s\to\hat G_t$ are the arrows $s\to t$ of
$\cS_n(\cX)$.
\item[(R9)] (\cite[Thm.~1.1]{Strings}.) $\PP\HH_L\Fin^\cX\simeq
\Fin^{\cS_n(\cX)}$, by the functor extending $U\mapsto(s\mapsto\hat
G_s(U))$; $\HH_L\Fin^\cX$ corresponds to the full subcategory of the
subfunctors of the functors $\pi^*A=A\pi$, $A$ in $\Fin^\cX$; and
$\HH_L\Fin^\cX$ is a pretopos iff for every $W\colon\cS_n(\cX)\to\Fin$
the unit $W\to\pi^*\Lan_\pi W$ of the adjunction $\Lan_\pi\dashv\pi^*$
is injective. A functor $W$ is isomorphic to a subfunctor of some
$\pi^*A$ iff its unit is injective.
\end{itemize}

For strings $t$, $s$ let $V_t=\cX(t_0,s_0)$ and $K_t=\cS_n(t,s)
\subseteq V_t$; $t$ is \emph{over} $s$ if $K_t\ne\emptyset$. A
\emph{one-step arrow} is an arrow $\beta\colon t'\to t$ of $\cS_n$
where $t'$ has one object more than $t$; if $t$ is over $s$, then $t'$
is over $s$. A \emph{leg tree} $T$ at $t$ \cite[Def.~6.14]{Strings} is
a finite rooted tree whose root is labelled by $t$ and whose other
nodes are labelled by one-step arrows into the string of the parent
node; so the root has children $(\beta,T')$ with $\beta\colon t'\to t$
one-step and $T'$ a leg tree at $t'$. Equivalence relations $E_T$ on
$V_t$ and $F_T$ on $K_t$ are defined by recursion: the equality if the
root has no children, and otherwise
\[
E_T=\bigvee\beta^*E_{T'},\qquad F_T=\bigvee\beta^*F_{T'},
\]
the joins of equivalence relations over the children of the root,
where $g\mathrel{\beta^*E}g'$ iff $g\beta\mathrel Eg'\beta$. A
\emph{realisation} of $T$ in $W\colon\cS_n\to\Fin$ is a family
$y_\nu\in W(t_\nu)$, $\nu$ a node with string $t_\nu$, with
$W(\beta)y_\nu=y_{\nu'}$ for $\nu$ a child of $\nu'$ with label
$\beta$. A \emph{morphism} of leg trees at the same string is a map of
trees preserving roots, parents and labels. For $t=s$ we have
$V_s=G_s:=\Aut_\cX(s_0)$ and $K_s=H_s:=\Aut_{\cS_n}(s)$, the group of
the first components of the automorphisms of the string $s$ (which
need not determine the other components).

\begin{itemize}
\item[(R10)] (\cite[proof of Thm.~6.15]{Strings}.) For $W\colon
\cS_n(\cX)\to\Fin$, a string $s$ and $a,b\in W(s)$: the unit of $W$
identifies $a$ and $b$ if and only if there are $\Theta\in H_s$ and a
leg tree $T$ at $s$ realised in $W$ with $a$ at the root such that
$b=W(\Theta)a$ and $\id\mathrel{E_T}\Theta$. If $T$ is realised in $W$
with $a$ at the root and $\id\mathrel{F_T}\Theta$, then
$W(\Theta)a=a$. For every leg tree $T$ at $s$ there is $W_T\colon
\cS_n(\cX)\to\Fin$ with a realisation of $T$, with $a$ at the root,
such that for $\Theta\in H_s$: $W_T(\Theta)a=a$ iff
$\id\mathrel{F_T}\Theta$.
\item[(R11)] (\cite[Lemma~8.1(a),(c)]{Proftrees}.) A morphism of leg
trees $T\to T''$ gives $E_T\subseteq E_{T''}$ and $F_T\subseteq
F_{T''}$. For $W$, $s$ and $a\in W(s)$ let $\Max(s,a)$ be the leg tree
at $s$ whose root has as children the pairs $(\beta,y)$ of a one-step
arrow $\beta\colon t'\to s$ and $y\in W(t')$ with $W(\beta)y=a$, with
label $\beta$ and subtree $\Max(t',y)$. It is realised in $W$ with $a$
at the root, and every leg tree at $s$ realised in $W$ with $a$ at the
root has a morphism to it. The unit of $W$ is injective at $s$ iff for
all $a\in W(s)$ and $\Theta\in H_s$ with $\id\mathrel{E_{\Max(s,a)}}
\Theta$ we have $W(\Theta)a=a$.
\end{itemize}
The relations at the root $s$ of a leg tree at $s$ (over $s$) are
invariant under left multiplication: $E_T$ is $g^{-1}g'\in S_T$ and
$F_T$ is $h^{-1}h'\in R_T$ for subgroups $R_T\le H_s\cap S_T$ of $G_s$
\cite[\S6]{Strings}. In \cite{Strings} these statements are made for a
category of strings with one string in each isomorphism class; they
are transported along the equivalence with $\cS_n(\cX)$.

\subsection*{Finite languages}
Let $\cC$ have finite sorts. For every object $A$ there is $n_A$ such
that $A$ has at most $n_A$ elements in every model
\cite[Lemma~4.1]{Finsorts}. For a finite language $\Lambda=(S,R)$ (in
\cite{Finsorts} it is denoted by $T$) let $\bT_\Lambda$ be the
coherent theory with the sorts $S$, the relation symbols $R$ and as
axioms all coherent sequents in this language valid in $\cC$, and
$\cX_\Lambda$ the category of the models of $\bT_\Lambda$ whose sorts
are sets $\{1,\dots,e_A\}$, $e_A\le n_A$. It is a finite category
equivalent to the category of all models of $\bT_\Lambda$, and for
$\Lambda\subseteq\Lambda'$ the restriction of structures is a functor
$q_{\Lambda'\Lambda}\colon\cX_{\Lambda'}\to\cX_\Lambda$; these
functors are strictly compatible.
\begin{itemize}
\item[(R12)] (\cite[Prop.~5.1, Lemma~5.2]{Finsorts}.) The pretopos
$\cC_\Lambda$ of $\bT_\Lambda$ is equivalent to $\Fin^{\cX_\Lambda}$
by the evaluation, and $\cC$ is the directed colimit of the
$\cC_\Lambda$ over any cofinal set of finite languages, along coherent
functors which correspond to the functors $q_{\Lambda'\Lambda}^*$.
\item[(R13)] (\cite[Prop.~5.4]{Finsorts}.) If $\Mod(\cC)$ is an
EI-category, the finite languages $\Lambda$ for which $\cX_\Lambda$ is
an EI-category are cofinal.
\end{itemize}

\subsection*{Localisation}
For a coherent category $\cC$ and $s\le1$ the slice $\cC/s$ is the
full subcategory of the objects whose map to $1$ factors through $s$;
a slice of an $L$-category is an $L$-category, a slice of a pretopos is
a pretopos, and the pullback functor $A\mapsto A\times s$ is coherent,
resp.\ Heyting. Let $D=\Sub_\cC(1)$. For a filter $F$ and an ideal $I$
of $D$ let $\cC_{F,I}$ be the pretopos of the theory of $\cC$ with the
additional axioms $\top\vdash s$, $s\in F$, and $s'\vdash\bot$, $s'\in
I$ \cite[\S4]{Proftrees}; we write $\cC_{\neg s'}$ if $F=\{1\}$ and
$I$ is generated by $s'$.
\begin{itemize}
\item[(R14)] (\cite[Lemma~4.1]{Proftrees}.) $\HH_L(\cC/s)\simeq
(\HH_L\cC)/s$; and for an $L$-category $\cH$ and $\tau\le1$,
composition with the pullback functor is an equivalence from the
Heyting functors on $\cH/\tau$ to the Heyting functors $G$ on $\cH$
with $G(\tau)\cong1$.
\item[(R15)] (\cite[Lemma~4.2, Thm.~4.3]{Proftrees}.) $\HH_L(\cC_{\neg
s'})\simeq(\HH_L\cC)/\neg s'$. If $\HH_L\cC$ is a pretopos, so is
$\HH_L\cC_{F,I}$ for every filter $F$ and ideal $I$ of $D$.
\end{itemize}

\section{The models of the translated theory are strings}\label{sec:strings}

In this section $\cC$ is a pretopos with finite sorts such that
$\Mod(\cC)$ is an EI-category, $L$ proves linearity and
$\cH=\HH_L\cC$. The argument follows \cite[\S\S3--4]{Strings}, where
$\Mod(\cC)$ is a finite category; the finiteness of the category is
replaced by the lemma on types and by the sentences, and the bound on
the length (Lemma~\ref{lem:classes}(c)) needs a new argument.

\begin{lem}[linearity]\label{lem:lin}
Let $\varphi\colon\hat G\to\hat G'$ and $\psi\colon\hat G\to\hat G''$
be homomorphisms of models of $\cH$. Then there is a homomorphism
$\chi\colon\hat G'\to\hat G''$ with $\chi\varphi=\psi$, or a
homomorphism $\chi'\colon\hat G''\to\hat G'$ with $\chi'\psi=\varphi$.
\end{lem}
\begin{proof}
Let $\hat G$ be over $M$. Apply (R2) to the coherent category $\cH$,
the models $\hat G'$, $\hat G''$, and the families of the elements
$\varphi(x)$ and $\psi(x)$, indexed by the elements $x$ of the sets
$A(M)$, $A$ in $\cC$. A homomorphism $\chi\colon\hat G'\to\hat G''$
with $\chi(\varphi(x))=\psi(x)$ for all $x$ satisfies $\chi\varphi=
\psi$. So if there is no such $\chi$, the hypothesis of (R2) fails:
there are a tuple $\alpha\in B(M)$, $B$ a product of objects of $\cC$,
and $E'\le B$ in $\cH$ with $\varphi(\alpha)\in\hat G'(E')$ and
$\psi(\alpha)\notin\hat G''(E')$. If there is no $\chi'$ either, there
are $\alpha_1\in B_1(M)$ and $E''\le B_1$ with $\psi(\alpha_1)\in\hat
G''(E'')$, $\varphi(\alpha_1)\notin\hat G'(E'')$. Replacing $B$, $B_1$
by $B\times B_1$, the tuples by $(\alpha,\alpha_1)$ and $E'$, $E''$ by
their pullbacks, we may assume $B=B_1$, $\alpha=\alpha_1$. In $\cH$,
$(E'\Rightarrow E'')\vee(E''\Rightarrow E')$ is the top subobject of
$B$; so $\alpha\in\hat G(E'\Rightarrow E'')$, say. Then
$\varphi(\alpha)\in\hat G'(E'\Rightarrow E'')\cap\hat G'(E')\subseteq
\hat G'(E'')$, a contradiction; the other case is symmetric.
\end{proof}

Write $\varphi\preceq\psi$ if $\psi=\chi\varphi$ for a homomorphism
$\chi$ of models of $\cH$; by Lemma~\ref{lem:lin} this is a total
preorder on the homomorphisms out of a model.

\begin{lem}\label{lem:iso}
A homomorphism of models of $\cH$ whose underlying homomorphism of
models of $\cC$ is invertible is an isomorphism of models of $\cH$.
\end{lem}
\begin{proof}
The proof of \cite[Lemma~4.1]{Strings} applies. One shows, for all
homomorphisms $\chi\colon\hat G'\to\hat G''$ with invertible
underlying homomorphism at once and all $U\le A$, that $b\in\hat
G'(U)$ if $\chi(b)\in\hat G''(U)$, by induction on the construction of
$U$. For $\forall_fU$ and $U\Rightarrow V$ one uses (R6): given a
homomorphism $\lambda\colon\hat G'\to\hat G'''$, by
Lemma~\ref{lem:lin} there is a homomorphism $\mu$ with
$\mu\chi=\lambda$, and then $\lambda$ factors through $\chi$, or with
$\mu\lambda=\chi$. In the second case the underlying homomorphism of
$\lambda$ has the left inverse $\chi^{-1}\mu$, so it is invertible,
$\Mod(\cC)$ being an EI-category, and so is that of $\mu\colon\hat
G'''\to\hat G''$; the induction hypothesis is applied to $\mu$.
\end{proof}

\begin{lem}\label{lem:classes}
Let $\hat G$ be a model of $\cH$ and $\varphi\colon\hat G\to\hat G'$,
$\psi\colon\hat G\to\hat G''$ homomorphisms, over $M'$ and $M''$.
\begin{enumerate}
\item[(a)] If $\varphi\preceq\psi\preceq\varphi$, then
$\psi=\chi\varphi$ with $\chi$ an isomorphism of models of $\cH$, and
$\Sent(M')=\Sent(M'')$.
\item[(b)] If $\varphi\preceq\psi$ and not $\psi\preceq\varphi$, then
$\Sent(M')\subsetneq\Sent(M'')$.
\item[(c)] Suppose $n_L$ is finite. Then the homomorphisms out of
$\hat G$ form at most $n_L$ classes for the equivalence
$\varphi\preceq\psi\preceq\varphi$.
\end{enumerate}
\end{lem}
\begin{proof}
(a) Let $\psi=\chi\varphi$ and $\varphi=\chi'\psi$. There are
homomorphisms $M'\to M''\to M'$, so $\Sent(M')=\Sent(M'')$ and $\chi$
is invertible by (R5), hence an isomorphism of models of $\cH$ by
Lemma~\ref{lem:iso}. (b) Let $\psi=\chi\varphi$. If $\chi$ were
invertible as a homomorphism of models of $\cC$, it would be an
isomorphism of models of $\cH$ and $\varphi=\chi^{-1}\psi$. So $\chi$
is not invertible, and $\Sent(M')\subsetneq\Sent(M'')$ by (R5).

(c) Let $n=n_L$ and suppose there are homomorphisms $\varphi_0\prec
\varphi_1\prec\dots\prec\varphi_n$ out of $\hat G$, with targets over
$M^0,\dots,M^n$. By (b) choose sentences $s_i$, $1\le i\le n$, true in
$M^i$ and false in $M^{i-1}$; replacing $s_i$ by $s_1\wedge\dots\wedge
s_i$ we have $s_1\ge\dots\ge s_n$, and $s_i$ is true in $M^j$ iff
$j\ge i$, the sets $\Sent(M^j)$ being increasing. Put $s_0=1$ and
$s_{n+1}=0$. For a homomorphism $\psi$ out of $\hat G$ with target
$\hat G^\psi$ over $M^\psi$ let $\ell(\psi)\in\{0,\dots,n\}$ be the
largest $i$ with $s_i$ true in $M^\psi$; then $s_i$ is true in
$M^\psi$ iff $i\le\ell(\psi)$; $\ell$ is monotone for $\preceq$; and
$\ell(\varphi_j)=j$.

Let $\mu$ be a propositional formula in variables $p_1,\dots,p_k$ and
$v$ a valuation of the variables by up-sets $\{j\mid j\ge i_e\}$,
$0\le i_e\le n+1$, of the chain $\{0<\dots<n\}$; let $\mu^v\le1$ be
the sentence of $\cH$ obtained by substituting $s_{i_e}$ for $p_e$.
\emph{Claim: for every homomorphism $\psi$ out of $\hat G$, the
sentence $\mu^v$ is true in $\hat G^\psi$ iff the node $\ell(\psi)$ of
the chain forces $\mu$ under $v$.} By induction on $\mu$. For a
variable this is the property of $\ell$; the cases of $\wedge$,
$\vee$, $\bot$ are clear, a model being a coherent functor. Let
$\mu=\mu_1\to\mu_2$. By (R6), $\mu^v$ is true in $\hat G^\psi$ iff for
every homomorphism $\rho$ out of $\hat G^\psi$, $\mu_1^v$ true in the
target of $\rho$ implies $\mu_2^v$ true there; by induction, iff for
every $\rho$ the node $\ell(\rho\psi)$ does not force $\mu_1$ or
forces $\mu_2$. The numbers $\ell(\rho\psi)$ are $\ge\ell(\psi)$, and
every $l\ge\ell(\psi)$ occurs: for $l=\ell(\psi)$ take $\rho=\id$; for
$l>\ell(\psi)$ we have not $\varphi_l\preceq\psi$, $\ell$ being
monotone, so $\psi\preceq\varphi_l$ by Lemma~\ref{lem:lin},
$\varphi_l=\rho\psi$. So the condition is the forcing of
$\mu_1\to\mu_2$ at the node $\ell(\psi)$.

Now let $\mu$ be a formula of $L$. Then $\mu^v$ is the top element of
$\Sub_\cH(1)$, so it is true in $\hat G^{\varphi_j}$ for every $j$,
and by the claim every node $j=\ell(\varphi_j)$ of the chain forces
$\mu$ under every valuation. So the chain with $n+1$ elements is a
frame of $L$, contradicting $n=n_L$.
\end{proof}

\begin{thm}\label{thm:strings}
Let $\cC$ be a pretopos with finite sorts such that $\Mod(\cC)$ is an
EI-category, and $L$ a logic proving linearity with $n_L$ finite.
\begin{enumerate}
\item[(a)] Every model of $\HH_L\cC$ is isomorphic to $\hat G_s$ for a
string $s$.
\item[(b)] For strings $s=(M_0\xrightarrow{a_1}\dots\xrightarrow{a_r}
M_r)$ and $s'$, the homomorphisms $\hat G_s\to\hat G_{s'}$ are the
$\iota_0a_j\cdots a_1$, for $0\le j\le r$ and $\iota$ an isomorphism
of strings from $s_{\ge j}$ onto $s'$; the index $j$ is determined by
$s'$. In particular $\hat G_s\cong\hat G_{s'}$ iff $s\cong s'$.
\end{enumerate}
\end{thm}
\[
\begin{tikzcd}[column sep=large]
M_0\ar[r,"a_1"]&\cdots\ar[r,"a_j"]&M_j\ar[r,"a_{j+1}"]\ar[d,"\iota_0",
"\cong"']&\cdots\ar[r,"a_r"]&M_r\ar[d,"\iota_{r-j}","\cong"']\\
&&M'_0\ar[r]&\cdots\ar[r]&M'_{r-j}
\end{tikzcd}
\]
\begin{proof}
\emph{Step 1: the string of a model.} Let $\hat G$ be a model over
$M_0$. By Lemma~\ref{lem:classes} choose representatives $\id=
\varphi_0\prec\varphi_1\prec\dots\prec\varphi_r$ of the classes of the
homomorphisms out of $\hat G$, $r+1\le n_L$ ($\id$ is in the least
class), $\varphi_j\colon\hat G\to\hat G_j$ over $M_j$, and
homomorphisms $a_j\colon\hat G_{j-1}\to\hat G_j$ with
$\varphi_j=a_j\varphi_{j-1}$. By Lemma~\ref{lem:classes}(b) and (R5)
the $a_j$ are not invertible; so $s=(M_0\xrightarrow{a_1}\dots
\xrightarrow{a_r}M_r)$ is a string.

\emph{Step 2: every homomorphism $\lambda$ out of $\hat G_j$ is $\chi
a_{j'}\cdots a_{j+1}$ for some $j'\ge j$ and an isomorphism of models
$\chi$.} The homomorphism $\lambda\varphi_j$ out of $\hat G$ is
equivalent to some $\varphi_{j'}$, and $j'\ge j$ because
$\varphi_j\preceq\lambda\varphi_j$; so $\lambda\varphi_j=\iota
a_{j'}\cdots a_{j+1}\varphi_j$ with an isomorphism of models $\iota$,
by Lemma~\ref{lem:classes}(a). By Lemma~\ref{lem:lin} at $\hat G_j$,
$\lambda$ and $\iota a_{j'}\cdots a_{j+1}$ differ by a homomorphism
$\theta$ between their targets, in one of the two directions; the
targets have the same sentences (they are targets of homomorphisms in
the class of $\varphi_{j'}$), so $\theta$ is invertible by (R5), and
an isomorphism of models by Lemma~\ref{lem:iso}.

\emph{Step 3: $\hat G_j(U)=\|U\|^s_j$ for all $j$ and $U\le A$.} By
induction on the construction of $U$, as in
\cite[Prop.~4.2, Step~3]{Strings}: by (R6) and Step 2 the conditions
for $\forall_fU$ and $U\Rightarrow V$ at $\hat G_j$ are, after applying
the inverses of the isomorphisms $\chi$, the forcing clauses at the
node $j$. As $r+1\le n_L$, $\hat G_s$ is defined and $\hat G=\hat
G_0=\hat G_s$. This proves (a).

\emph{Step 4: \emph{(b)}.} Let $s$ be a string. The homomorphisms
$a_i\cdots a_1$ out of $\hat G_s$ are pairwise inequivalent, their
targets having different sentences. We show that every homomorphism
$\psi\colon\hat G_s\to\hat G'$, over $M'$, is $\chi a_i\cdots a_1$
with $\chi$ an isomorphism of models. Let $i$ be maximal with
$a_i\cdots a_1\preceq\psi$, and $\psi=\lambda a_i\cdots a_1$ with
$\lambda\colon\hat G_{s_{\ge i}}\to\hat G'$; by Lemma~\ref{lem:iso}
it suffices to show that $\lambda\colon M_i\to M'$ is invertible.
Suppose not; by (R5) there is a sentence $t$ true in $M'$ and false in
$M_i$. If $i=r$: the node $r$ is the top of the chain, $t$ is false
there, so $\neg t$ is true in $\hat G_{s_{\ge r}}$, hence in $\hat
G'$, where $t$ is true; then $0=t\wedge\neg t$ would be true in $\hat
G'$. Let $i<r$. By Lemma~\ref{lem:lin} and the maximality of $i$ there
is a homomorphism $\mu\colon\hat G'\to\hat G_{s_{\ge i+1}}$ with
$\mu\lambda=a_{i+1}$. If $\mu$ were invertible, it would be an
isomorphism of models and $a_{i+1}\cdots a_1\preceq\psi$. So by (R5)
there is a sentence $u$ true in $M_{i+1}$ and false in $M'$. The node
$i$ of $s$ forces $t\Rightarrow u$: $t$ is false at the node $i$, and
$u$ is true at the nodes $j>i$. So $t\Rightarrow u$ is true in $\hat
G_{s_{\ge i}}$, hence in $\hat G'$, where $t$ is true and $u$ is
false: a contradiction.

It remains to show that the isomorphisms of models $\chi_0\colon\hat
G_u\to\hat G_t$ are the first components of the isomorphisms of
strings $u\to t$. An isomorphism of strings preserves forcing.
Conversely, $\chi_0$ induces a bijection between the homomorphisms out
of $\hat G_u$ and out of $\hat G_t$ preserving $\preceq$, whose
classes are those of the initial parts $a_i\cdots a_1$ of $u$ and
$b_i\cdots b_1$ of $t$; so $u$ and $t$ have the same number of models,
and if it is more than one, $b_1\chi_0=\chi_1a_1$ for an isomorphism
of models $\chi_1\colon\hat G_{u_{\ge1}}\to\hat G_{t_{\ge1}}$, which
extends to an isomorphism of strings by induction. The index $j$ in
(b) is determined because the targets of the $a_j\cdots a_1$ are
pairwise non-isomorphic.
\end{proof}

\begin{cor}\label{cor:cons}
In the situation of Theorem~\ref{thm:strings} let $\cS$ be the
category of the strings of standard models, with the homomorphisms of
\emph{(b)} as arrows. Then $U\mapsto(s\mapsto\hat G_s(U))$ is a
conservative Heyting functor $\HH_L\cC\to\Set^{\cS}$ with finite
values, sending $A$ in $\cC$ to $s\mapsto A(s_0)$.
\end{cor}
\begin{proof}
It is a coherent functor, conservative by the completeness theorem for
the coherent category $\HH_L\cC$ and (a). In $\Set^\cS$ the operations
$\Rightarrow$ and $\forall$ are computed by quantifying over the
arrows out of an object; so the functor preserves $U\Rightarrow V$ and
$\forall_fU$, for $f$ in $\cC$, by (R6), the homomorphisms out of
$\hat G_s$ being the arrows of $\cS$ out of $s$ by (b). For a morphism
$h\colon X\to Y$ of $\HH_L\cC$ with $X\le A$, $Y\le B$ and graph
$\Gamma_h\le A\times B$, and $W\le X$, we have $\forall_hW=
Y\wedge\forall_{p_B}(\Gamma_h\Rightarrow p_A^*W)$ in every Heyting
category; so the functor preserves $\forall_h$.
\end{proof}

This proves Theorem~\ref{thm:B}.

\begin{rem}\label{rem:infty}
If $n_L=\infty$ and the chains of prime filters of $\Sub_\cC(1)$ have
at most $m$ elements, Theorem~\ref{thm:strings} and
Corollary~\ref{cor:cons} hold with strings of at most $m$ models: in
Lemma~\ref{lem:classes} the number of classes is at most $m$ by (b),
and the rest of the proof is unchanged.
\end{rem}

For $m\ge1$ let $L_m$ be the logic of the chain with $m$ elements; it
proves linearity, $n_{L_m}=m$, and $L\subseteq L_m$ if $m\le n_L$.

\begin{prop}\label{prop:Lm}
Let $\cC$ be a pretopos with finite sorts such that $\Mod(\cC)$ is an
EI-category, $L$ a logic proving linearity and $m\le n_L$ a positive
integer such that $m=n_L$, or the chains of prime filters of
$\Sub_\cC(1)$ have at most $m$ elements. Then $\HH_L\cC$ is an
$L_m$-category, and $\eta\colon\cC\to\HH_L\cC$ is the free
$L_m$-category on $\cC$.
\end{prop}
\begin{proof}
All strings have at most $m$ models, by definition if $m=n_L$ and by
(R5) otherwise. By
Corollary~\ref{cor:cons} and Remark~\ref{rem:infty} the lattices of
subobjects of $\HH_L\cC$ are Heyting subalgebras of those of
$\Set^\cS$. For $F\colon\cS\to\Set$ the lattice $\Sub(F)$ is the
Heyting algebra of the up-sets of the preorder of the pairs $(s,x)$,
$x\in F(s)$, with $(s,x)\le(t,y)$ iff $y=F(g)x$ for an arrow $g\colon
s\to t$. Restriction to the principal up-sets ${\uparrow}(s,x)$ embeds
it into the product of their algebras of up-sets. By
Theorem~\ref{thm:strings}(b), for arrows $g=\iota_0a_j\cdots a_1$ and
$g'=\iota'_0a_{j'}\cdots a_1$ out of $s$ with $j\le j'$ we have
$g'=hg$ for the arrow $h=\iota'_0a_{j'}\cdots a_{j+1}\iota_0^{-1}$ of
$\cS$. So ${\uparrow}(s,x)$ is a total preorder with at most $m$
classes, and its algebra of up-sets is that of a chain with at most
$m$ elements, a generated subframe of the chain with $m$ elements. So
the formulas of $L_m$ are valid in $\Sub(F)$, and $\HH_L\cC$ is an
$L_m$-category. As $L\subseteq L_m$, the $L_m$-categories are the
$L$-categories in which the formulas of $L_m$ are valid, with the same
Heyting functors; so the universal property of $\HH_L\cC$ among the
$L$-categories gives that of the free $L_m$-category.
\end{proof}

\section{Reduction to the theories below a prime filter}\label{sec:prime}

In this section $\cC$ is any pretopos, $D=\Sub_\cC(1)$, $L$ is an
intermediate logic proving the weak excluded middle and
$\cH=\HH_L\cC$. For a prime filter $Q$ of $D$ let
$\cC_{(Q)}=\cC_{\{1\},D\setminus Q}$, the pretopos of the theory of
the models $M$ of $\cC$ with $\Sent(M)\subseteq Q$. It is the directed
colimit of the $\cC_{\neg s'}$, $s'\notin Q$; so by (R3), (R14) and
(R15), $\HH_L\cC_{(Q)}$ is the directed colimit of the slices
$\cH/\neg s'$, $s'\notin Q$, along the pullback functors.

\begin{lem}\label{lem:prod}
Let $\tau_1,\dots,\tau_k\le1$ in $\cH$ be pairwise disjoint with join
$1$. Then $A\mapsto(A\times\tau_i)_i$ is an equivalence
$\cH\simeq\prod_i\cH/\tau_i$, and $\cH$ is a pretopos if and only if
every $\cH/\tau_i$ is.
\end{lem}
\begin{proof}
An object $A$ is the disjoint join of its subobjects $A\times\tau_i$,
hence their coproduct; so a morphism $A\to B$ is a family of morphisms
$A\times\tau_i\to B\times\tau_i$, and the functor is full and
faithful. Let $B_i$ be objects of $\cH/\tau_i$, $B_i\le A_i$ with
$A_i$ in $\cC$. In the pretopos $\cC$ the coproduct $A=\coprod_iA_i$
exists, and $\eta$ preserves it, a coherent functor preserving the
property of being a disjoint join of subobjects. The $B_i$ are
pairwise disjoint subobjects of $A$, and their join $B\le A$ has
$B\times\tau_i=B_i$, since $B_j\times\tau_i=0$ for $j\ne i$. So the
functor is essentially surjective. A finite product of categories with
terminal objects is a pretopos iff every factor is: a diagram in one
factor is completed by terminal objects in the others, and limits,
colimits, images and joins in the product are computed factorwise.
\end{proof}

\begin{lem}\label{lem:partition}
Let $D_0\subseteq D$ be a finite sublattice containing $0$ and $1$.
For a prime filter $Q_0$ of $D_0$ let
\[
\tau_{Q_0}=\bigwedge_{s\in Q_0}\neg\neg s\ \wedge\bigwedge_{s'\in
D_0\setminus Q_0}\neg s'\ \le1\quad\text{in }\cH.
\]
The $\tau_{Q_0}$ are pairwise disjoint with join $1$. For finite
sublattices $D_0\subseteq D_1$ and a prime filter $Q_1$ of $D_1$,
$\tau_{Q_1}\le\tau_{Q_1\cap D_0}$.
\end{lem}
\begin{proof}
In a Heyting algebra satisfying the weak excluded middle, $\neg
a\vee\neg b=\neg(a\wedge b)$; so the elements $\neg a$ form a
sublattice $B$ which is a Boolean algebra with complement $\neg$, and
$h\colon s\mapsto\neg\neg s$ is a homomorphism of bounded lattices
from $D_0$ to $B\subseteq\Sub_\cH(1)$: it preserves meets in every
Heyting algebra, and $\neg\neg(a\vee b)=\neg(\neg a\wedge\neg
b)=\neg\neg a\vee\neg\neg b$. In the Boolean algebra $B$, $1$ is the
join of the pairwise disjoint elements $\epsilon_X=\bigwedge_{s\in
X}h(s)\wedge\bigwedge_{s'\in D_0\setminus X}\neg h(s')$,
$X\subseteq D_0$, and $\neg h(s')=\neg s'$. If $X$ is not a prime
filter of $D_0$, then $\epsilon_X=0$: this is clear if $1\notin X$ or
$0\in X$; if $a\in X$, $a\le b\notin X$, then $h(a)\wedge\neg h(b)=0$;
if $a,b\in X$ and $a\wedge b\notin X$, then $h(a)\wedge h(b)\wedge\neg
h(a\wedge b)=0$; if $a\vee b\in X$ and $a,b\notin X$, then $h(a\vee
b)\wedge\neg h(a)\wedge\neg h(b)=0$. The last statement is clear, the
meet for $Q_1$ having more terms.
\end{proof}

For a prime filter $Q$ of $D$ let $\cH_{(Q)}$ be the directed colimit
of the slices $\cH/\tau_{Q\cap D_0}$ over the finite sublattices $D_0$
of $D$ containing $0$ and $1$, along the pullback functors.

\begin{thm}\label{thm:stalk}
Let $\cC$ be a pretopos and $L$ an intermediate logic proving the weak
excluded middle. The following are equivalent:
\begin{itemize}
\item[(i)] $\HH_L\cC$ is a pretopos;
\item[(ii)] $\cH_{(Q)}$ is a pretopos for every prime filter $Q$ of
$D$;
\item[(iii)] $\HH_L\cC_{(Q)}$ is a pretopos for every prime filter $Q$
of $D$.
\end{itemize}
\end{thm}
\begin{proof}
(i)$\Rightarrow$(iii) is (R15).

(iii)$\Rightarrow$(ii). For a finite sublattice $D_0$ let $s'_0$ be
the join of $D_0\setminus Q$; then $s'_0\notin Q$, $Q$ being prime,
$\neg s'_0=\bigwedge_{s'\in D_0\setminus Q}\neg s'$, and
$\tau_{Q\cap D_0}=\neg s'_0\wedge\sigma_{Q\cap D_0}$, where
$\sigma_X=\bigwedge_{s\in X}\neg\neg s$ for a finite $X\subseteq Q$.
The pairs $(s'_0,Q\cap D_0)$ are cofinal among the pairs $(s',X)$ of a
sentence $s'\notin Q$ and a finite $X\subseteq Q$, ordered by $s'\le
s'_1$ and $X\subseteq X_1$: take for $D_0$ the sublattice generated by
$X\cup\{s'\}$, which is finite. So $\cH_{(Q)}$ is the directed colimit
of the slices $\cH/(\neg s'\wedge\sigma_X)$ over these pairs, that is
the directed colimit over $X$ of the slices at $\sigma_X$ of
$\varinjlim_{s'\notin Q}\cH/\neg s'=\HH_L\cC_{(Q)}$; here we use that
an object of a directed colimit whose map to $1$ factors through
$\sigma_X$ is represented by such an object at a stage. These slices
are pretoposes by (iii), and so is their directed colimit.

(ii)$\Rightarrow$(i). Let $E$ be an equivalence relation on an object
$Z$ of $\cH$. In a coherent category, a cover $Z\to P$ with kernel
pair $E$ is an effective quotient of $E$, stable under pullback. For
every prime filter $Q$ the image of $E$ in the pretopos $\cH_{(Q)}$
has a quotient; the object $P$, the morphism, and the finitely many
relations between subobjects saying that it is a cover with kernel
pair $E$ are represented at a stage: there is a finite sublattice
$D_0(Q)$ such that $E\times\tau_{Q\cap D_0(Q)}$ has a quotient in
$\cH/\tau_{Q\cap D_0(Q)}$. The sets $V_Q=\{Q'\mid Q'\cap D_0(Q)=Q\cap
D_0(Q)\}$ are open in the patch topology of the set of prime filters
of $D$, which is compact; so finitely many $V_{Q_1},\dots,V_{Q_k}$
cover. Let $D_0$ be the finite sublattice generated by the
$D_0(Q_i)$. Every prime filter $Q_0$ of $D_0$ is $Q'\cap D_0$ for a
prime filter $Q'$ of $D$; if $Q'\in V_{Q_i}$, then
$\tau_{Q_0}\le\tau_{Q'\cap D_0(Q_i)}=\tau_{Q_i\cap D_0(Q_i)}$ by
Lemma~\ref{lem:partition}, so $E\times\tau_{Q_0}$ has a quotient in
$\cH/\tau_{Q_0}$, a slice of a slice. By Lemmas~\ref{lem:partition}
and~\ref{lem:prod}, $E$ has a quotient in $\cH$. The same argument
gives the disjoint finite coproducts: the coproduct of two objects
exists in each $\cH_{(Q)}$, hence at a stage, where the property of
being a disjoint join of the two subobjects is represented.
\end{proof}

Theorem~\ref{thm:C} is the equivalence of (i) and (iii). The open
parts do not behave in this way. For $L$ proving linearity with
$n_L\ge3$ let $\cC$ be the
pretopos $\cC(G;A;T_1,T_2)$ of \cite[Thm.~1.2]{Proftrees} with
$A\cap(T_1+T_2)\ne(A\cap T_1)+(A\cap T_2)$, so that $\HH_L\cC$ is not
a pretopos, and $s_j$ the sentence which is false exactly in the model
$0_j$. Then $s_1\vee s_2=1$, and the slices $\cC/s_j$ are the
pretoposes of the families with one subgroup, for which
$\HH_L(\cC/s_j)$ is a pretopos.

\begin{cor}\label{cor:locfin}
Let $\cC$ be a pretopos with finite sorts such that every prime filter
of $D$ contains only finitely many prime filters, and $L$ a
non-classical logic proving linearity. Then $\HH_L\cC$ is a pretopos
if and only if, for every prime filter $Q$, the pretopos $\cC_{(Q)}$,
which has finite sorts and finitely many sentences, satisfies the
criterion of \cite[Thm.~1.1]{Proftrees}: $\Mod(\cC_{(Q)})$ is an
EI-category and no profinite leg tree with strings of at most $n_L$
models is an obstruction.
\end{cor}
\begin{proof}
A logic proving linearity proves the weak excluded middle. The lattice
of sentences of $\cC_{(Q)}$ is the quotient $D_Q$ of $D$ by the
congruence generated by the ideal $D\setminus Q$; its prime filters
are the prime filters of $D$ contained in $Q$, finitely many, so it is
finite. The models of $\cC_{(Q)}$ are models of $\cC$, so it has
finite sorts. Theorem~\ref{thm:stalk} and
\cite[Thm.~1.1]{Proftrees} apply.
\end{proof}

\section{Functors which do not reflect isomorphisms}\label{sec:collapse}

From now on $L$ proves linearity and $n=n_L$. Let $q\colon\cX'\to\cX$
be a functor between finite EI-categories. For a string
$s'=(m_0\xrightarrow{a_1}\dots\xrightarrow{a_r}m_r)$ of $\cX'$ let
$j_0=0<j_1<\dots<j_e$ be the indices $j$ such that $j=0$ or $q(a_j)$
is not invertible, and $\tilde q(s')$ the sequence
\[
q(m_{j_0})\to q(m_{j_1})\to\dots\to q(m_{j_e})
\]
whose arrows are the images of the composites of the arrows of $s'$.
Each of these composites contains exactly one arrow with
non-invertible image, so $\tilde q(s')$ is a string of $\cX$, with
first object $q(m_0)$.
\[
\begin{tikzcd}[column sep=large,row sep=small]
s'\colon&m_0\ar[r,"a_1"]&m_1\ar[r,"a_2"]&m_2\ar[r,"a_3"]&m_3\\
\tilde q(s')\colon&q(m_0)\ar[rr,"q(a_2a_1)"]&&q(m_2)\ar[r,"q(a_3)"]&q(m_3)
\end{tikzcd}
\]
(here $q(a_1)$ is invertible and $q(a_2)$, $q(a_3)$ are not).

\begin{lem}\label{lem:collapse}
\begin{enumerate}
\item[(a)] For a sequence $c$ as in \emph{(R7)} in which $a_{j+1}$ is
invertible, $\|U\|^c_j=a_{j+1}^{-1}\|U\|^c_{j+1}$ for all $U$, and
$\|U\|^c_i=\|U\|^{c'}_i$ for $i\le j$, where $c'$ is obtained by
deleting the node $j+1$ and composing $a_{j+2}a_{j+1}$ (deleting the
last arrow if $j+1=r$).
\item[(b)] For a string $s'$ of $\cX'$ and $U\le A$ in $\HH_L\Fin^\cX$:
$\hat G_{s'}(\HH_L(q^*)U)=\hat G_{\tilde q(s')}(U)$ as subsets of
$A(q(m_0))$.
\item[(c)] $q$ sends the arrows $s'\to t'$ of $\cS_n(\cX')$ to arrows
$\tilde q(s')\to\tilde q(t')$ of $\cS_n(\cX)$; so $\tilde q\colon
\cS_n(\cX')\to\cS_n(\cX)$ is a functor with $\pi\tilde q=q\pi$, and
under the equivalences of \emph{(R9)} the functor $\PP\HH_L(q^*)$ is
$W\mapsto W\tilde q$.
\item[(d)] For composable functors, $\widetilde{qq'}=\tilde q\tilde
q'$. If $\beta'\colon t'\to s'$ is a one-step arrow of $\cS_n(\cX')$,
then $q(\beta')\colon\tilde q(t')\to\tilde q(s')$ is a one-step arrow
or an isomorphism of strings.
\end{enumerate}
\end{lem}
\begin{proof}
(a) Write $a=a_{j+1}$. We show $\beta\in\|U\|_j$ iff $a\beta\in
\|U\|_{j+1}$ by induction on the construction of $U$. For $U$ in $\cC$
this holds because $a$ is an isomorphism of models; it is compatible
with meets, joins, $f^*$, and with $\exists_f$ because $a$ is
bijective on every object. For $U\Rightarrow V$: the clause at the
node $j$ is the conjunction of the conditions for $j'=j$ and for
$j'\ge j+1$; the latter is the clause for $a\beta$ at the node $j+1$,
and the condition for $j'=j$ is, by the induction hypothesis,
equivalent to the condition for $j'=j+1$. The case of $\forall_f$ is
the same. For the second statement, by induction on $U$: in the
clauses at a node $i\le j$ of $c$ the condition for $j'=j+1$ is
equivalent to the one for $j'=j$ by the first statement, and the
remaining conditions are those of $c'$.

(b) The composite of $\HH_L(q^*)$ with the Heyting functor of the
Kripke model $s'$ is a Heyting functor extending the Kripke model
$q(s')$ of $\Fin^\cX$, the sequence of the $q(a_j)$, since the
evaluation at an object $m$ composed with $q^*$ is the evaluation at
$q(m)$. By the uniqueness in (R7), $\hat G_{s'}(\HH_L(q^*)U)=
\|U\|^{q(s')}_0$, and by (a), applied repeatedly from the right, this
is $\|U\|^{\tilde q(s')}_0$.

(c) By (R8) an arrow $g\colon s'\to t'$ of $\cS_n(\cX')$ is a
homomorphism $\hat G_{s'}\to\hat G_{t'}$; so $q(g)$ maps $\hat
G_{s'}(\HH_L(q^*)U)$ into $\hat G_{t'}(\HH_L(q^*)U)$ for all $U$, that
is, it is a homomorphism $\hat G_{\tilde q(s')}\to\hat G_{\tilde
q(t')}$ by (b), an arrow of $\cS_n(\cX)$ by (R8). For an object $Z$ of
$\PP\HH_L\Fin^\cX$ corresponding to $W$, the value of $\PP\HH_L(q^*)Z$
at $s'$ is the value at $Z$ of the model $\hat G_{s'}\HH_L(q^*)$,
extended to the pretopos completion, which is $W(\tilde q(s'))$ by
(b), naturally in $s'$.

(d) In an EI-category a composite is invertible iff its factors are;
so deleting the arrows which become invertible under $q'$ and then
those which become invertible under $q$ deletes the arrows which
become invertible under $qq'$. The final segment of $t'$ after its
first object is isomorphic to $s'$, so its image under $\tilde q$ has
as many objects as $\tilde q(s')$, and $\tilde q(t')$ has the same
number of objects or one more, according to whether the image of the
first arrow of $t'$ is invertible or not.
\end{proof}

\begin{lem}[the eventual criterion]\label{lem:eventual}
Let $(\cX_i)_{i\in I}$ be an inverse system of finite EI-categories
and functors $q_{i'i}\colon\cX_{i'}\to\cX_i$ over a directed set, and
$\cC=\varinjlim\Fin^{\cX_i}$ along the functors $q^*$. Then $\HH_L\cC$
is a pretopos if and only if for every $i$ and every
$W\colon\cS_n(\cX_i)\to\Fin$ there is $i'\ge i$ such that the unit of
$W\tilde q_{i'i}$ is injective.
\end{lem}
\begin{proof}
For functors which reflect isomorphisms this is
\cite[Lemma~7.3]{Proftrees}. The proof given there uses only the
following: by (R3) $\HH_L\cC$ is the colimit of the
$\HH_L\Fin^{\cX_i}$; by (R9) these are the full subcategories of the
pretoposes $\Fin^{\cS_n(\cX_i)}$ of the functors with injective unit;
and the transitions are the restrictions of the coherent functors
$W\mapsto W\tilde q_{i'i}$, which now holds by
Lemma~\ref{lem:collapse}(c). In outline: if the condition holds, the
quotient $W$ of an equivalence relation at a stage $i$, formed in
$\Fin^{\cS_n(\cX_i)}$, becomes an object $W\tilde q_{i'i}$ of
$\HH_L\Fin^{\cX_{i'}}$ and is the quotient there and at all later
stages, hence in the colimit, and likewise for finite coproducts.
Conversely, a functor $W$ at the stage $i$ is the quotient of a finite
coproduct of representable functors $\cS_n(\cX_i)(s,-)$, which are
objects of $\HH_L\Fin^{\cX_i}$, being subfunctors of the functors
$\pi^*\cX_i(s_0,-)$; if $\HH_L\cC$ is a pretopos, the quotient formed
in $\HH_L\cC$ is represented at a stage $i'\ge i$ by an object of
$\HH_L\Fin^{\cX_{i'}}$ which is identified with $W\tilde q_{i'i}$ in
the pretopos $\Fin^{\cS_n(\cX_{i'})}$.
\end{proof}

\section{Pullbacks of leg trees}\label{sec:pull}

Fix a functor $q\colon\cX'\to\cX$ of finite EI-categories. In this
section all leg trees are taken relative to a base string: for a leg
tree at $t$ the relations $E_T$, $F_T$ are on $V_t=\cX(t_0,s_0)$ and
$K_t=\cS_n(t,s)$ for a fixed string $s$ with $t$ over $s$; in $\cX'$
the base string is a string $s'$ and in $\cX$ it is $\tilde q(s')$.
As $\tilde q$ is a functor preserving first objects, $q$ maps $V_{t'}$
to $V_{\tilde q(t')}$ and $K_{t'}$ to $K_{\tilde q(t')}$.

\emph{Transport.} Let $T$ be a leg tree at $t$ and $\iota\colon t_1\to
t$ an isomorphism of strings. The leg tree $\iota^*T$ at $t_1$ has as
children of the root the $(\iota^{-1}\beta,T_1)$ for the children
$(\beta,T_1)$ of the root of $T$. Then $h\mathrel{E_{\iota^*T}}h'$ iff
$h\iota^{-1}\mathrel{E_T}h'\iota^{-1}$, and the same for $F$.

\emph{Pullback.} Let $T$ be a leg tree of $\cX$ at $t$ and $t'$ a
string of $\cX'$ with $\tilde q(t')=t$. The leg tree $T^{t'}$ of
$\cX'$ at $t'$ is defined by recursion on the number of strings of
$\cX'$ with more objects than $t'$: the children of its root are
\begin{itemize}
\item the pairs $(\beta',\nu)$ of a child $\nu=(\beta,T_1)$ of the
root of $T$ and a one-step arrow $\beta'\colon t''\to t'$ of $\cX'$
with $q(\beta')=\beta$ as arrows of $\cS_n(\cX)$ (so $\tilde q(t'')$
is the string of $\nu$), with label $\beta'$ and subtree $T_1^{t''}$;
\item the one-step arrows $\beta'\colon t''\to t'$ such that
$\iota=q(\beta')\colon\tilde q(t'')\to t$ is an isomorphism of
strings, with label $\beta'$ and subtree $(\iota^*T)^{t''}$: the
\emph{collapsing children}.
\end{itemize}
\[
\begin{tikzcd}[column sep=small,row sep=small]
&t''_1\ar[drr,"\beta'_1"']&t''_2\ar[dr,"\beta'_2"]&&t''_3\ar[dl,
"\beta'_3"']&&&\tilde q(t''_1)=\tilde q(t''_2)\ar[d,"\beta"]&\tilde
q(t''_3)\ar[dl,"\iota","\cong"']\\
T^{t'}\colon&&&t'&&&T\colon&t&
\end{tikzcd}
\]
Here $\beta'_1$, $\beta'_2$ lie over the child $\beta$ of $T$ and
$\beta'_3$ is a collapsing child, above which the whole tree $T$ is
pulled back again. If $q$ is an identity, $T^t=T$.

\begin{lem}\label{lem:pull}
\begin{enumerate}
\item[(P1)] For $W\colon\cS_n(\cX)\to\Fin$, a string $s'$ of $\cX'$
and $a\in W(\tilde q(s'))$: the tree $\Max(s',a)$ for $W\tilde q$ is
$\Max(\tilde q(s'),a)^{s'}$, where $\Max(\tilde q(s'),a)$ is taken for
$W$.
\item[(P2)] Let $q'\colon\cX''\to\cX'$ be another functor of finite
EI-categories, $T$ a leg tree of $\cX$ at $t$, $t''$ a string of
$\cX''$, $t'=\tilde q'(t'')$ and $\tilde q(t')=t$. Then
$(T^{t'})^{t''}=T^{t''}$, the left side formed with $q$ and then $q'$,
the right side with $qq'$.
\item[(P3)] $q(E_{T^{t'}})\subseteq E_T$ and $q(F_{T^{t'}})\subseteq
F_T$.
\item[(P4)] A morphism of leg trees $T\to T_0$ at $t$ induces a
morphism $T^{t'}\to T_0^{t'}$.
\end{enumerate}
Here trees are identified along the evident bijections of nodes
preserving labels.
\end{lem}
\begin{proof}
First, for an isomorphism of strings $\iota'\colon t'_1\to t'$ of
$\cX'$ with $q(\iota')=\iota\colon t_1\to t$:
\begin{equation}\label{eq:tp}
(\iota^*T)^{t'_1}=\iota'^*(T^{t'}),
\end{equation}
by $\beta'\mapsto\iota'\beta'$ on the labels of the children of the
roots: $q(\beta')$ is $\iota^{-1}\beta$ for a child $(\beta,T_1)$ of
$T$ iff $q(\iota'\beta')=\beta$, and $q(\beta')$ is an isomorphism of
strings iff $q(\iota'\beta')$ is; in the first case the subtrees are
$T_1^{t''}$ on both sides, in the second they are $(q(\beta')^*
\iota^*T)^{t''}$ and $(q(\iota'\beta')^*T)^{t''}$, which are equal.

(P1) The children of the root of $\Max(s',a)$ for $W\tilde q$ are the
$(\beta',y)$ with $\beta'\colon t''\to s'$ one-step, $y\in W(\tilde
q(t''))$ and $W(q(\beta'))y=a$. By Lemma~\ref{lem:collapse}(d),
$q(\beta')$ is one-step or an isomorphism of strings. If $q(\beta')$
is one-step, the condition says that $(q(\beta'),y)$ is a child $\nu$
of the root of $\Max(\tilde q(s'),a)$, and $(\beta',y)$ corresponds to
$(\beta',\nu)$; the subtrees are $\Max(t'',y)$ and $\Max(\tilde
q(t''),y)^{t''}$, equal by induction. If $\iota=q(\beta')$ is an
isomorphism of strings, $W(\iota)$ is bijective and
$y=W(\iota)^{-1}a$ is the only possibility; so each collapsing
$\beta'$ gives one child. Its subtree is $\Max(t'',y)=\Max(\tilde
q(t''),y)^{t''}$ by induction, and $\Max(\tilde
q(t''),W(\iota)^{-1}a)=\iota^*\Max(\tilde q(s'),a)$: $(\beta_1,z)$ is
a child of the root of the left side iff $W(\iota\beta_1)z=a$, iff
$(\iota\beta_1,z)$ is a child of the root of $\Max(\tilde q(s'),a)$,
and the subtrees are the same.

(P2) Let $\beta''\colon u''\to t''$ be one-step in $\cX''$. Then
$q'(\beta'')$ is one-step or an isomorphism of strings, and so is
$qq'(\beta'')$; it is an isomorphism of strings if $q'(\beta'')$ is.
If $q'(\beta'')=\beta'$ is one-step: $\beta''$ gives children of the
root of $(T^{t'})^{t''}$, one for each child of the root of $T^{t'}$
with label $\beta'$; these are the children $\nu$ of the root of $T$
with label $q(\beta')$, or the one collapsing child if $q(\beta')$ is
an isomorphism of strings; and these index the children of the root of
$T^{t''}$ with label $\beta''$. The subtrees are $(T_1^{u'})^{u''}$
and $T_1^{u''}$, resp.\ $((\iota^*T)^{u'})^{u''}$ and
$(\iota^*T)^{u''}$, with $u'=\tilde q'(u'')$, equal by induction. If
$\iota'=q'(\beta'')$ is an isomorphism of strings, $\beta''$ is a
collapsing child on both sides, with subtrees
$(\iota'^*(T^{t'}))^{u''}$ and $(q(\iota')^*T)^{u''}$; by
\eqref{eq:tp} the first is $((q(\iota')^*T)^{u'})^{u''}$, equal to the
second by induction.

(P3) By induction. Let $(g,g')\in\beta'^*E$ for a child of the root of
$T^{t'}$ with label $\beta'$ and relation $E$ of its subtree. For a
child $(\beta',\nu)$, by induction $q(g)\beta\mathrel{E_{T_1}}q(g')
\beta$, so $(q(g),q(g'))\in\beta^*E_{T_1}\subseteq E_T$. For a
collapsing child, by induction $q(g)\iota\mathrel{E_{\iota^*T}}q(g')
\iota$, that is $q(g)\mathrel{E_T}q(g')$. So $q$ maps the generators
of $E_{T^{t'}}$ into the equivalence relation $E_T$. The same for $F$,
as $q$ maps $K_{t'}$ to $K_t$.

(P4) A child $(\beta',\nu)$ is sent to $(\beta',\nu_0)$, $\nu_0$ the
image of $\nu$, with the induced morphism of the subtrees; a
collapsing child to the same collapsing child, with the morphism
induced by $\iota^*T\to\iota^*T_0$.
\end{proof}

\section{Inverse systems of finite categories}\label{sec:inverse}

Let $(\cX_i)_{i\in I}$ be an inverse system of finite EI-categories
and arbitrary functors $q_{i'i}$, $q_{i'i}q_{i''i'}=q_{i''i}$, over a
directed set, $\cM$ its limit (objects and arrows are the compatible
families; the sets of arrows between two objects are profinite),
$q_i\colon\cM\to\cX_i$ the projections and $\cC=\varinjlim
\Fin^{\cX_i}$, with the functors $J_i\colon\Fin^{\cX_i}\to\cC$. A
\emph{thread} is a family $\sigma=(s_i)_{i\in I}$ of strings $s_i$ of
$\cX_i$ with $\tilde q_{i'i}(s_{i'})=s_i$ for $i'\ge i$. The first
objects of the $s_i$ form an object $\sigma_0$ of $\cM$; let
$G_\sigma=\Aut_\cM(\sigma_0)=\varprojlim G_{s_i}$, a profinite group
with projections $\rho_i$, and
\[
H_\sigma=\{\Theta\in G_\sigma\mid\rho_i(\Theta)\in H_{s_i}\text{ for
all }i\},
\]
a closed subgroup; $q_{i'i}$ maps $H_{s_{i'}}$ into $H_{s_i}$ by
Lemma~\ref{lem:collapse}(c). The numbers of objects of the $s_i$ are
non-decreasing in $i$, and $\sigma$ is \emph{bounded} if they are
bounded; every thread is bounded if $n=n_L$ is finite. Strings of
$\cM$ (with at most $n$ objects), their isomorphisms and
automorphisms, and $\tilde q_i(\tilde s)$ are defined as for a finite
EI-category.

\begin{lem}\label{lem:threads}
\begin{enumerate}
\item[(a)] $\cM$ is an EI-category; an arrow $f$ of $\cM$ is
invertible iff all $q_i(f)$ are, and if $q_i(f)$ is not invertible,
neither is $q_{i'}(f)$ for $i'\ge i$.
\item[(b)] $\tilde s\mapsto(\tilde q_i(\tilde s))_i$ is a bijection
from the strings of $\cM$ to the bounded threads; and for a bounded
thread $(s_i)$ there is $i_0$ such that for $i'\ge i\ge i_0$ the
functor $q_{i'i}$ sends the arrows of $s_{i'}$ to the arrows of $s_i$.
\item[(c)] For a string $\tilde s$ of $\cM$ with thread $\sigma$, the
group $H_\sigma$ is the group $H_{\tilde s}$ of the first components
of the automorphisms of the string $\tilde s$. We write $G_{\tilde
s}=G_\sigma$.
\end{enumerate}
\end{lem}
\begin{proof}
(a) The inverse of an invertible arrow is unique, so the inverses of
the components are compatible; functors preserve invertibility.

(b) $\tilde q_i(\tilde s)=\tilde q_{i'i}\tilde q_{i'}(\tilde s)$ as in
Lemma~\ref{lem:collapse}(d), so the family is a thread, bounded by the
number of objects of $\tilde s$. By (a), for $i$ large no arrow of
$\tilde s$ has an invertible image, so $\tilde s$ is determined by the
family. Conversely let $(s_i)$ be a bounded thread. The numbers of
objects of the $s_i$ are constant for $i\ge i_0$; then $\tilde
q_{i'i}(s_{i'})=s_i$ means that $q_{i'i}$ sends the arrows of $s_{i'}$
to those of $s_i$, and the families of arrows are the arrows of a
string $\tilde s$ of $\cM$, non-invertible by (a), with $\tilde
q_i(\tilde s)=s_i$ for $i\ge i_0$, hence for all $i$. It has at most
$n$ objects, as the $s_i$ have.

(c) Let $i_0$ be as in (b). For $i\ge i_0$ the automorphisms of the
string $\tilde s$ restrict to automorphisms of the string $s_i$; so
$\rho_i(H_{\tilde s})\subseteq H_{s_i}$ for $i\ge i_0$, and then for
all $i$. Conversely let $\rho_i(\Theta)\in H_{s_i}$ for all $i$. For
$i\ge i_0$ the set of the automorphisms of the string $s_i$, that is
of the families of automorphisms of its objects commuting with its
arrows, with first component $\rho_i(\Theta)$ is finite and non-empty;
these sets form an inverse system, and an element of the limit is an
automorphism of the string $\tilde s$ with first component $\Theta$.
\end{proof}

Let $\sigma=(s_i)$ be a thread, $i\in I$ and $T$ a leg tree of $\cX_i$
at $s_i$, over $s_i$. For $i'\ge i$ let $T^{s_{i'}}$ be the pullback
along $q_{i'i}$, a leg tree of $\cX_{i'}$ at $s_{i'}$ over $s_{i'}$,
with the subgroups $R_{T^{s_{i'}}}\le H_{s_{i'}}\cap S_{T^{s_{i'}}}$
of $G_{s_{i'}}$. Put
\[
S_{(i,T)}=\bigcap_{i'\ge i}\rho_{i'}^{-1}\bigl(S_{T^{s_{i'}}}\bigr),
\qquad
R_{(i,T)}=\bigcap_{i'\ge i}\rho_{i'}^{-1}\bigl(R_{T^{s_{i'}}}\bigr).
\]
The pair $(i,T)$ is an \emph{obstruction at $\sigma$} if
$R_{(i,T)}\ne H_\sigma\cap S_{(i,T)}$, and an obstruction at a string
of $\cM$ if it is one at its thread.

\begin{lem}\label{lem:groups}
\begin{enumerate}
\item[(a)] For $i''\ge i'\ge i$, $q_{i''i'}$ maps $S_{T^{s_{i''}}}$
into $S_{T^{s_{i'}}}$ and $R_{T^{s_{i''}}}$ into $R_{T^{s_{i'}}}$. So
$S_{(i,T)}$ and $R_{(i,T)}$ are closed subgroups of $G_\sigma$,
the intersections are directed, and $R_{(i,T)}\le H_\sigma\cap
S_{(i,T)}$.
\item[(b)] $S_{(i',T^{s_{i'}})}=S_{(i,T)}$ and
$R_{(i',T^{s_{i'}})}=R_{(i,T)}$ for $i'\ge i$.
\item[(c)] $\rho_i(H_\sigma\cap S_{(i,T)})=\bigcap_{i'\ge
i}q_{i'i}(H_{s_{i'}}\cap S_{T^{s_{i'}}})$.
\item[(d)] No pair $(i,T)$ is an obstruction at $\sigma$ if and only
if
\begin{equation}\label{eq:crit}
\rho_i(H_\sigma\cap S_{(i,T)})\subseteq R_T
\end{equation}
for every $i$ and every leg tree $T$ of $\cX_i$ at $s_i$.
\end{enumerate}
\end{lem}
\begin{proof}
(a) By (P2), $T^{s_{i''}}=(T^{s_{i'}})^{s_{i''}}$, and (P3) applies to
$q_{i''i'}$ and the tree $T^{s_{i'}}$. As $R_{T^{s_{i'}}}\le
H_{s_{i'}}$ for all $i'\ge i$, and $\rho_j=q_{i'j}\rho_{i'}$ for
$i'\ge i,j$, we get $R_{(i,T)}\le H_\sigma$. (b) By (P2) the groups on
the left are the intersections over $i''\ge i'$ of the same families,
which by (a) is the intersection over $i''\ge i$. (c) The left side is
contained in the right side. Let $\theta$ be in the right side. The
sets of the $\Theta'\in H_{s_{i'}}\cap S_{T^{s_{i'}}}$ with
$q_{i'i}(\Theta')=\theta$ are finite and non-empty, and form an
inverse system by (a) and Lemma~\ref{lem:collapse}(c); an element of
the limit is $\Theta\in G_\sigma$ with $\rho_i(\Theta)=\theta$ and
$\rho_{i'}(\Theta)\in H_{s_{i'}}\cap S_{T^{s_{i'}}}$ for $i'\ge i$, so
$\Theta\in H_\sigma\cap S_{(i,T)}$. (d) If no pair is an obstruction, then
$H_\sigma\cap S_{(i,T)}=R_{(i,T)}\subseteq\rho_i^{-1}(R_{T^{s_i}})$,
and $T^{s_i}=T$. Conversely, \eqref{eq:crit} for the pair
$(i',T^{s_{i'}})$ and (b) give $H_\sigma\cap
S_{(i,T)}\subseteq\rho_{i'}^{-1}(R_{T^{s_{i'}}})$ for all $i'\ge i$,
so $H_\sigma\cap S_{(i,T)}\subseteq R_{(i,T)}$.
\end{proof}

\begin{thm}\label{thm:general}
Let $(\cX_i)_{i\in I}$ be an inverse system of finite EI-categories,
$\cC=\varinjlim\Fin^{\cX_i}$, and $L$ a logic proving linearity. Then
$\HH_L\cC$ is a pretopos if and only if no pair $(i,T)$ of an index
$i$ and a leg tree $T$ of $\cX_i$ is an obstruction at a thread. If
all threads are bounded, in particular if $n_L$ is finite, the threads
are the strings of the limit $\cM$.
\end{thm}
\begin{proof}
Fix $i$ and $W\colon\cS_n(\cX_i)\to\Fin$. For $i'\ge i$ let
$\mathcal B_{i'}$ be the finite set of the triples $(s',a,\theta)$ of
a string $s'$ of $\cX_{i'}$, $a\in W(\tilde q_{i'i}s')$ and
$\theta=q_{i'i}(\Theta')$ for some $\Theta'\in H_{s'}$ with
$\id\mathrel{E_{\Max(\tilde q_{i'i}s',a)^{s'}}}\Theta'$, such that
$W(\theta)a\ne a$. By (R11) and (P1), the unit of $W\tilde q_{i'i}$ is
injective iff $\mathcal B_{i'}=\emptyset$. For $i''\ge i'$ the map
$(s'',a,\theta)\mapsto(\tilde q_{i''i'}s'',a,\theta)$ sends $\mathcal
B_{i''}$ to $\mathcal B_{i'}$: if $\Theta''\in H_{s''}$ with
$\id\mathrel{E_{M^{s''}}}\Theta''$, $M=\Max(\tilde q_{i''i}s'',a)$,
then with $s'=\tilde q_{i''i'}s''$ we have $M^{s''}=(M^{s'})^{s''}$ by
(P2), so $\Theta'=q_{i''i'}(\Theta'')\in H_{s'}$ satisfies
$\id\mathrel{E_{M^{s'}}}\Theta'$ by (P3), and
$q_{i'i}(\Theta')=\theta$. These maps form an inverse system of finite
sets, so some $\mathcal B_{i'}$ is empty iff the limit is empty. An
element of the limit is a family of strings $s_{i'}$, $i'\ge i$, with
$\tilde q_{i''i'}(s_{i''})=s_{i'}$, which extends uniquely to a thread
$\sigma$ (put $s_j=\tilde q_{i'j}(s_{i'})$ for $i'\ge i,j$), together
with $a\in W(s_i)$ and $\theta\in\bigcap_{i'\ge i}q_{i'i}(H_{s_{i'}}
\cap S_{\Max(s_i,a)^{s_{i'}}})$ with $W(\theta)a\ne a$; by
Lemma~\ref{lem:groups}(c) the condition on $\theta$ is $\theta\in
\rho_i(H_\sigma\cap S_{(i,\Max(s_i,a))})$. With
Lemma~\ref{lem:eventual}: $\HH_L\cC$ is a pretopos iff for all $i$,
$W$, threads $\sigma$ and $a\in W(s_i)$,
\begin{equation}\label{eq:W}
\rho_i\bigl(H_\sigma\cap S_{(i,\Max(s_i,a))}\bigr)\ \text{fixes }a.
\end{equation}
By Lemma~\ref{lem:groups}(d) it remains to show that \eqref{eq:W} for
all $i$, $W$, $\sigma$, $a$ is equivalent to \eqref{eq:crit} for all
$i$, $\sigma$, $T$. If \eqref{eq:crit} holds, then \eqref{eq:W}
holds: $\Max(s_i,a)$ is realised in $W$ with $a$ at the root, so
$R_{\Max(s_i,a)}$ fixes $a$ by (R10). Conversely assume \eqref{eq:W}
and let $T$ be a leg tree at $s_i$. Let $W=W_T$ as in (R10), with $a$
at the root. There is a morphism $T\to\Max(s_i,a)$ by (R11), hence
morphisms $T^{s_{i'}}\to\Max(s_i,a)^{s_{i'}}$ by (P4), so
$S_{(i,T)}\subseteq S_{(i,\Max(s_i,a))}$ by (R11). So
$\rho_i(H_\sigma\cap S_{(i,T)})$ fixes $a$, and it is contained in
$R_T$ by the property of $W_T$. The last statement is
Lemma~\ref{lem:threads}(b).
\end{proof}

\begin{prop}\label{prop:models}
In the situation of Theorem~\ref{thm:general}:
\begin{enumerate}
\item[(a)] For a thread $\sigma=(s_i)$ there is a model $\hat
G_\sigma$ of $\HH_L\cC$ over the evaluation at $\sigma_0$, with $\hat
G_\sigma(\HH_L(J_i)U)=\hat G_{s_i}(U)$ for $U$ in $\HH_L\Fin^{\cX_i}$.
\item[(b)] Every model of $\HH_L\cC$ is isomorphic to $\hat G_\sigma$
for a thread $\sigma$.
\item[(c)] $H_\sigma$ is the group of the automorphisms $\Theta$ of
$\sigma_0$ with $\Theta(\hat G_\sigma(U))=\hat G_\sigma(U)$ for all
$U$, that is, the automorphism group of the model $\hat G_\sigma$.
\end{enumerate}
\end{prop}
\begin{proof}
A model $F$ of $\cC$ is isomorphic to the evaluation at an object of
$\cM$: the sets of the pairs of an object $x_i$ of $\cX_i$ and an
isomorphism from $FJ_i$ to the evaluation at $x_i$ are finite and
non-empty, they form an inverse system, and an element of the limit is
an object $x$ of $\cM$ with an isomorphism from $F$ to the evaluation
at $x$. By (R3), $\HH_L\cC$ is the colimit of the $\HH_L\Fin^{\cX_i}$;
so a model of $\HH_L\cC$ over the evaluation at $x$ is a family of
models $\hat G_i$ of $\HH_L\Fin^{\cX_i}$ over $q_ix$ with $\hat
G_{i'}(\HH_L(q_{i'i}^*)U)=\hat G_i(U)$ for $i'\ge i$ and all $U$.

(a) The $\hat G_{s_i}$ are such a family by
Lemma~\ref{lem:collapse}(b). (b) Let $\hat G=(\hat G_i)$ be a model
over the evaluation at $x$, and $\Sigma_i$ the set of the strings $s$
of $\cX_i$ with first object $q_ix$ and $\hat G_s=\hat G_i$. It is
finite, and non-empty: by (R8) there are a string $t$ and an
isomorphism of models $\hat G_i\to\hat G_t$, given by an isomorphism
$\varphi\colon q_ix\to t_0$; replacing the first object of $t$ by
$q_ix$ and its first arrow $a_1$ by $a_1\varphi$ (if $t$ has one
object: $t$ by $(q_ix)$) gives a string $s$ and an isomorphism of
strings $s\to t$ with first component $\varphi$, so $\hat
G_s(U)=\varphi^{-1}\hat G_t(U)=\hat G_i(U)$. By
Lemma~\ref{lem:collapse}(b), $\tilde q_{i'i}$ maps $\Sigma_{i'}$ to
$\Sigma_i$: for $s'\in\Sigma_{i'}$, $\hat G_{\tilde q_{i'i}(s')}(U)=
\hat G_{s'}(\HH_L(q_{i'i}^*)U)=\hat G_i(U)$. An element of the limit
of the $\Sigma_i$ is a thread $\sigma$ with $\hat G_\sigma=\hat G$.
(c) An automorphism $\Theta$ of $\sigma_0$ preserves all $\hat
G_\sigma(U)$ iff every $\rho_i(\Theta)$ preserves all $\hat
G_{s_i}(U)$, iff $\rho_i(\Theta)\in H_{s_i}$ for all $i$, by (R8): an
automorphism of the model $\hat G_{s_i}$ is an arrow $s_i\to s_i$ of
$\cS_n(\cX_i)$.
\end{proof}

\subsection*{Obstructions occur at the strings of the limit}
The relation $E_T$ of a leg tree depends on the base string only
through its first object $m_*$: it is a relation on $\cX(t_0,m_*)$,
defined by composition with the arrows of $\cX$ underlying the labels
(the sets $K_t$ and the relations $F_T$ depend on the whole base
string). We use that two arrows of $\cS_n(\cX)$ with the same source,
the same target and the same underlying arrow of $\cX$ are equal, and
that an invertible arrow of $\cS_n(\cX)$ whose underlying arrow is an
identity is the first component of an isomorphism of strings
$(\id,\mu_1,\dots)$.

\begin{lem}\label{lem:first}
Let $q\colon\cX'\to\cX$ be a functor of finite EI-categories,
$W\colon\cS_n(\cX)\to\Fin$, and $u$, $v$ strings of $\cX'$ with the
same first object such that $v$ has at most as many objects as $u$.
Let $\lambda\colon\tilde q(v)\to\tilde q(u)$ be an invertible arrow of
$\cS_n(\cX)$ whose underlying arrow is an identity, and $a\in W(\tilde
q(u))$. Then, for the maximal trees of $W\tilde q$,
\[
E_{\Max(u,a)}\subseteq E_{\Max(v,W(\lambda)^{-1}a)}
\]
as relations on $\cX'(u_0,m_*)$, for every object $m_*$.
\end{lem}
\begin{proof}
By induction on the largest number $N(u_0)$ of composable
non-invertible arrows of $\cX'$ ending at $u_0$; it is finite, as in
an EI-category such a sequence does not meet an isomorphism class of
objects twice. Let $(\beta,y)$ be a child of the root of $\Max(u,a)$:
$\beta\colon x\to u$ is one-step, $x=(y_0\xrightarrow{p}w)$,
$\iota\colon w\to u$ is an isomorphism of strings, the underlying
arrow of $\beta$ is $g=\iota_0p$, and $y\in W(\tilde q(x))$ with
$W(\tilde q\beta)y=a$. Let $\bar w$ be the string obtained from $v$ by
replacing its first object by $w_0$ and its first arrow $c$ by
$c\iota_0$ ($\bar w=(w_0)$ if $v$ has one object);
$\kappa=(\iota_0,\id,\dots,\id)$ is an isomorphism of strings $\bar
w\to v$. Let $\bar x=(y_0\xrightarrow{p}\bar w)$, a string with at
most as many objects as $x$, and $\bar\beta\colon\bar x\to v$ the
one-step arrow with underlying arrow $\kappa_0p=g$.
\[
\begin{tikzcd}[column sep=large,row sep=small]
x=(y_0\xrightarrow{p}w)\ar[r,"\beta"]&u&\tilde q(x)\ar[r,"\tilde
q\beta"]&\tilde q(u)\\
\bar x=(y_0\xrightarrow{p}\bar w)\ar[r,"\bar\beta"]&v&\tilde q(\bar
x)\ar[r,"\tilde q\bar\beta"]\ar[u,"\bar\lambda"]&\tilde
q(v)\ar[u,"\lambda"']
\end{tikzcd}
\]
The arrows $\tilde q(\iota_0)\colon\tilde q(w)\to\tilde q(u)$ and
$\tilde q(\kappa_0)\colon\tilde q(\bar w)\to\tilde q(v)$ of
$\cS_n(\cX)$ are invertible with underlying arrow $q(\iota_0)$; so
$\mu=\tilde q(\iota_0)^{-1}\lambda\,\tilde q(\kappa_0)\colon\tilde
q(\bar w)\to\tilde q(w)$ is invertible with underlying arrow the
identity of $q(w_0)$, the first component of an isomorphism of strings
$(\id,\mu_1,\dots)$. The strings $\tilde q(x)$ and $\tilde q(\bar x)$
are obtained from $\tilde q(w)$ and $\tilde q(\bar w)$ in the same
way: by adding the object $q(y_0)$ and the arrow $q(p)$ in front if
$q(p)$ is not invertible, and by replacing the first object by
$q(y_0)$ and composing the first arrow, if there is one, with $q(p)$
if $q(p)$ is invertible. Accordingly $(\id_{q(y_0)},\id,\mu_1,\dots)$,
resp.\ $(\id_{q(y_0)},\mu_1,\dots)$, is an isomorphism of strings
$\tilde q(\bar x)\to\tilde q(x)$; let $\bar\lambda$ be the
corresponding arrow of $\cS_n(\cX)$, invertible, with underlying arrow
the identity of $q(y_0)$. The arrows $\tilde q(\beta)\bar\lambda$ and
$\lambda\tilde q(\bar\beta)$ from $\tilde q(\bar x)$ to $\tilde q(u)$
have the same underlying arrow $q(g)$, so they are equal. Hence for
$\bar y=W(\bar\lambda)^{-1}y\in W(\tilde q(\bar x))$
\[
W(\tilde q\bar\beta)\bar y=W(\lambda)^{-1}W(\tilde
q\beta)W(\bar\lambda)\bar y=W(\lambda)^{-1}a,
\]
and $(\bar\beta,\bar y)$ is a child of the root of
$\Max(v,W(\lambda)^{-1}a)$, with subtree $\Max(\bar x,\bar y)$. The
strings $x$, $\bar x$ have the first object $y_0$, with
$N(y_0)<N(u_0)$, and $x$, $\bar x$, $\bar\lambda$, $y$ satisfy the
hypotheses; by induction $E_{\Max(x,y)}\subseteq E_{\Max(\bar x,\bar
y)}$ on $\cX'(y_0,m_*)$. As $\beta$ and $\bar\beta$ have the same
underlying arrow $g$,
\[
\beta^*E_{\Max(x,y)}\subseteq\bar\beta^*E_{\Max(\bar x,\bar
y)}\subseteq E_{\Max(v,W(\lambda)^{-1}a)} .
\]
The join over the children $(\beta,y)$ gives the claim (if there are
no children, the left side is the equality).
\end{proof}

\begin{lem}\label{lem:sub}
Let $\sigma=(s_i)$ be a thread and $i\in I$. There is a bounded thread
$\tau=(t_j)$ such that $t_i=s_i$, the strings $t_j$ and $s_j$ have the
same first object for all $j$, for $j\ge i$ the string $t_j$ has as
many objects as $s_i$ and $\tilde q_{ji}(t_j)=s_i$, and
$H_\sigma\subseteq H_\tau$ in $G_\sigma=G_\tau$.
\end{lem}
\begin{proof}
Let $j\ge i$ and $s_j=(m_0\xrightarrow{a_1}\dots\xrightarrow{a_r}
m_r)$. Let $k_0=0<k_1<\dots<k_e$ be the indices $k$ such that $k=0$ or
$q_{ji}(a_k)$ is not invertible; so $\tilde q_{ji}(s_j)=s_i$ is the
sequence of the $q_{ji}(m_{k_l})$ and the images of the composites
$b_l\colon m_{k_{l-1}}\to m_{k_l}$ of arrows of $s_j$. Let $t_j$ be
the sequence of the $m_{k_l}$ and the $b_l$. The $b_l$ are not
invertible, as their images are not; so $t_j$ is a string with as many
objects as $s_i$, $q_{ji}$ sends its arrows to those of $s_i$, and
$\tilde q_{ji}(t_j)=s_i$; $t_i=s_i$.

Let $j'\ge j\ge i$ and $s_{j'}=(m'_0\xrightarrow{a'_1}\dots)$. An
arrow $a'_k$ whose image in $\cX_i$ is not invertible has a
non-invertible image in $\cX_j$. The string $\tilde
q_{j'j}(s_{j'})=s_j$ has the objects $q_{j'j}(m'_k)$, for $k=0$ and
the $k$ with $q_{j'j}(a'_k)$ not invertible, and its arrows are the
images of the composites of the arrows of $s_{j'}$ between them. Among
these $k$ are the indices $k'_l$ of the objects of $t_{j'}$; so the
$q_{j'j}(m'_{k'_l})$ are objects of $s_j$, the image of the arrow
$b'_l$ of $t_{j'}$ is the composite of the arrows of $s_j$ between
them, and the last of these arrows of $s_j$ is the only one among them
whose image in $\cX_i$ is not invertible, as $q_{j'i}(a'_k)$ is
invertible for $k'_{l-1}<k<k'_l$. Hence the $q_{j'j}(m'_{k'_l})$ are
the objects $m_{k_l}$ of $t_j$, and $q_{j'j}$ sends the arrows of
$t_{j'}$ to those of $t_j$, which are not invertible: $\tilde
q_{j'j}(t_{j'})=t_j$. The family $(t_j)_{j\ge i}$ extends uniquely to
a thread $\tau$ (put $t_j=\tilde q_{j'j}(t_{j'})$ for $j'\ge i,j$),
which is bounded. For $j\ge i$, $t_j$ and $s_j$ have the same first
object; for the other $j$ the first objects of $t_j$ and of
$s_j=\tilde q_{j'j}(s_{j'})$, $j'\ge i,j$, are the image of the first
object of $t_{j'}$ and $s_{j'}$. So $\sigma_0=\tau_0$ and
$G_\sigma=G_\tau$. An automorphism of the string $s_j$, $j\ge i$,
gives one of $t_j$ with the same first component (its components at
the $m_{k_l}$ commute with the composites $b_l$); so $H_{s_j}\subseteq
H_{t_j}$ for $j\ge i$. If $\Theta\in H_\sigma$, then
$\rho_j(\Theta)\in H_{t_j}$ for $j\ge i$, and for the other $j$,
$\rho_j(\Theta)=q_{j'j}(\rho_{j'}(\Theta))\in H_{t_j}$ with $j'\ge
i,j$, by Lemma~\ref{lem:collapse}(c); so $\Theta\in H_\tau$.
\end{proof}

\begin{thm}\label{thm:bounded}
In the situation of Theorem~\ref{thm:general}, with $n_L$ finite or
infinite: if some pair $(i,T)$ is an obstruction at a thread, then
some pair is an obstruction at a bounded thread. Consequently
$\HH_L\cC$ is a pretopos if and only if no pair $(i,T)$ of an index
$i$ and a leg tree $T$ of $\cX_i$ is an obstruction at a string of the
limit $\cM$.
\end{thm}
\begin{proof}
Suppose that no pair is an obstruction at a bounded thread; we show
\eqref{eq:W} for all $i$, $W$, $\sigma$ and $a$, so that $\HH_L\cC$ is
a pretopos and no pair is an obstruction at any thread, by
Theorem~\ref{thm:general} and its proof. Let $i\in I$,
$W\colon\cS_n(\cX_i)\to\Fin$, $\sigma=(s_j)$ a thread, $a\in W(s_i)$,
$M=\Max(s_i,a)$ and $\Theta\in H_\sigma\cap S_{(i,M)}$; we have to
show $W(\rho_i\Theta)a=a$. Let $\tau=(t_j)$ be as in
Lemma~\ref{lem:sub}; $t_i=s_i$, so $M$ is a leg tree at $t_i$, and
$\Theta\in H_\tau$. For $j\ge i$, $\rho_j(\Theta)\in S_{M^{s_j}}$,
where $M^{s_j}$ is by (P1) the tree $\Max(s_j,a)$ of $W\tilde q_{ji}$.
The strings $s_j$ and $t_j$ of $\cX_j$ have the same first object,
$t_j$ has at most as many objects as $s_j$, and $\tilde
q_{ji}(t_j)=s_i=\tilde q_{ji}(s_j)$; by Lemma~\ref{lem:first} for
$q_{ji}$, with $u=s_j$, $v=t_j$, $\lambda$ the identity of $s_i$ and
$m_*$ the first object of $s_j$, $E_{\Max(s_j,a)}\subseteq
E_{\Max(t_j,a)}$, and the latter tree is $M^{t_j}$ by (P1). So
$\rho_j(\Theta)\in S_{M^{t_j}}$ for all $j\ge i$: $\Theta$ is in the
group $S_{(i,M)}$ formed for the thread $\tau$. As no pair is an
obstruction at $\tau$, by Lemma~\ref{lem:groups}(d)
$\rho_i(\Theta)\in R_M$, which fixes $a$ by (R10). The last statement
follows with Lemma~\ref{lem:threads}(b).
\end{proof}

For $n_L=\infty$ the strings of $\cM$ have arbitrary finite length.
The proof moves a violation of \eqref{eq:W} from a thread to a bounded
thread whose strings are obtained by omitting objects; it does not
show that the same pair $(i,T)$ is an obstruction there.

\begin{rem}\label{rem:special}
(a) For a constant system, $S_{(i,T)}=S_T$ and $R_{(i,T)}=R_T$, and
Theorem~\ref{thm:general} is \cite[Thm.~6.15]{Strings}.

(b) If the functors $q_{i'i}$ reflect isomorphisms, there are no
collapsing children, and the pullback $T^{s_{i'}}$ consists of the
branches of $\cX_{i'}$ over those of $T$. For the towers of
\cite[\S7]{Proftrees}, whose functors are bijective on objects and
full, and whose threads are bounded, the pairs $(i,T)$ are the profinite leg
trees and Theorem~\ref{thm:general} is \cite[Thm.~8.2]{Proftrees}.

(c) If $n\le2$, there are no obstructions: a string with two objects
has no one-step arrows into it, and for a string $s=(m)$ and a
one-step arrow $\beta$ into $s$ the relations $\beta^*(=)$ on
$V_s=G_s$ and on $K_s=H_s=G_s$ coincide; the same holds for the
pullbacks.
\end{rem}

\section{Theories with finite sorts}\label{sec:theories}

Let $\cC$ be a pretopos with finite sorts such that $\Mod(\cC)$ is an
EI-category, and $I$ the set of the finite languages $\Lambda$ for
which $\cX_\Lambda$ is an EI-category, cofinal by (R13). The
categories $\cX_\Lambda$, $\Lambda\in I$, with the restriction
functors form an inverse system of finite EI-categories, and
$\cC\simeq\varinjlim_{\Lambda\in I}\Fin^{\cX_\Lambda}$ by (R12). Its
limit $\cM$ is the category of the standard models of $\cC$: a
compatible family of structures interprets all objects of $\cC$ by
sets $\{1,\dots,e\}$ and all subobjects of finite products of objects
by relations, so that all coherent sequents valid in $\cC$ hold; this
is a standard model, the morphisms of $\cC$ being interpreted through
their graphs; and a compatible family of homomorphisms is a natural
transformation. A thread $\sigma=(s_\Lambda)$ is a family of strings
$s_\Lambda$ of $\cX_\Lambda$, $\Lambda\in I$, with $\tilde
q_{\Lambda'\Lambda}(s_{\Lambda'})=s_\Lambda$; $\sigma_0$ is a standard
model of $\cC$, and for a leg tree $T$ of $\cX_\Lambda$ at $s_\Lambda$
we have the closed subgroups $R_{(\Lambda,T)}\le H_\sigma\cap
S_{(\Lambda,T)}$ of $G_\sigma=\Aut(\sigma_0)$ of
Section~\ref{sec:inverse}. The strings of $\cM$ are the strings of
standard models of $\cC$, and such a string $\tilde s$ has the thread
of the strings $\tilde s|_\Lambda=\tilde q_\Lambda(\tilde s)$.

\begin{cor}\label{cor:final}
Let $\cC$ be a pretopos with finite sorts and $L$ a non-classical
logic proving linearity. Then $\HH_L\cC$ is a pretopos if and only if
$\Mod(\cC)$ is an EI-category and
\[
R_{(\Lambda,T)}=H_{\tilde s}\cap S_{(\Lambda,T)}
\]
for every string $\tilde s$ of standard models, every $\Lambda\in I$
and every leg tree $T$ of $\cX_\Lambda$ at $\tilde s|_\Lambda$, the
groups being those of the thread of $\tilde s$ and $H_{\tilde s}$ the
group of the automorphisms of $\tilde s_0$ which extend to an
automorphism of the string $\tilde s$. The same holds with all
threads $\sigma=(s_\Lambda)$ and the groups $H_\sigma$ in place of the
strings of standard models. If $\Mod(\cC)$ is an EI-category, every
model of $\HH_L\cC$ is isomorphic to $\hat G_\sigma$ for a thread
$\sigma$, and $H_\sigma=\Aut(\hat G_\sigma)$.
\end{cor}
\begin{proof}
If $\HH_L\cC$ is a pretopos, $\Mod(\cC)$ is an EI-category by (R4). If
$\Mod(\cC)$ is an EI-category, Theorems~\ref{thm:general}
and~\ref{thm:bounded}, Lemma~\ref{lem:threads}(c) and
Proposition~\ref{prop:models} apply to the inverse system above.
\end{proof}

This is Theorem~\ref{thm:A}. By Lemma~\ref{lem:groups}(d) the
condition can be tested in the finite groups $H_{s_\Lambda}$: the
image of $H_\sigma\cap S_{(\Lambda,T)}$ is contained in $R_T$. By
Theorem~\ref{thm:stalk} it suffices to test the pretoposes
$\cC_{(Q)}$.

\begin{rem}\label{rem:strings}
Suppose that $n_L$ is finite, or that the chains of prime filters of
$\Sub_\cC(1)$ have at most $m$ elements. Then every thread is bounded;
so by Lemma~\ref{lem:threads} the threads are the families $(\tilde
s|_\Lambda)_\Lambda$ for the strings $\tilde s$ of standard models,
$H_\sigma$ is the group $H_{\tilde s}$ of the automorphisms of $\tilde
s_0$ which extend to an automorphism of the string $\tilde s$, and
by Corollary~\ref{cor:final} every model of $\HH_L\cC$ is isomorphic
to the model of the thread of a string of standard models. For
$n_L$ finite this is clear. In the other case let $\sigma=(s_\Lambda)$
be a thread. By Theorem~\ref{thm:strings}(a) and
Remark~\ref{rem:infty}, $\hat G_\sigma\cong\hat G_{\tilde s}$ for a
string $\tilde s$ of standard models with at most $m$ models; changing
the first model and the first arrow of $\tilde s$ along this
isomorphism, we may assume $\tilde s_0=\sigma_0$ and $\hat G_{\tilde
s}=\hat G_\sigma$. Let $\tau=(t_\Lambda)$ be the thread of $\tilde s$.
Then $\hat G_{\tilde s}(\HH_L(J_\Lambda)U)=\hat G_{t_\Lambda}(U)$, as
in Lemma~\ref{lem:collapse}(b): the composite of $\HH_L(J_\Lambda)$
with the Heyting functor of the Kripke model $\tilde s$ extends the
Kripke model of the restrictions to $\Lambda$, and
Lemma~\ref{lem:collapse}(a) applies. So $\hat G_{t_\Lambda}=\hat
G_{s_\Lambda}$, and by (R8) the identity of the first object is an
arrow $\iota_0p\colon s_\Lambda\to t_\Lambda$ of
$\cS_n(\cX_\Lambda)$; as it is invertible in $\cX_\Lambda$, $p$ is an
identity, and $s_\Lambda$ is isomorphic to $t_\Lambda$, which has at
most $m$ objects.
\end{rem}

\section{A family with legs converging to automorphisms}\label{sec:family}

Let $G$ be a profinite abelian group, $H$ a closed subgroup and
$(S_k)_{k\in\mathbb N}$ non-zero closed subgroups converging to $0$.
For an open subgroup $N$ let $I_N$ be the finite set of the $k$ with
$S_k\not\subseteq N$, and $\cX_N$ the category with the objects $0_k$,
$k\in I_N$, $1$, $2$, with
\begin{gather*}
\Aut(1)=G/N,\qquad\cX_N(0_k,1)=G/(S_k+N),\\
\cX_N(1,2)=G/(H+N),\qquad\cX_N(0_k,2)=G/(H+S_k+N),
\end{gather*}
trivial groups $\Aut(0_k)$, $\Aut(2)$, no other arrows, and the
addition as composition. This is a finite EI-category. For $N'\subseteq N$
let $q\colon\cX_{N'}\to\cX_N$ be the identity on $1$, $2$ and on the
$0_k$, $k\in I_N$, send $0_k$ to $1$ for $k\in I_{N'}\setminus I_N$,
and be given on the arrows by the quotient maps $G/(S_k+N')\to G/N$,
$G/(H+S_k+N')\to G/(H+N)$ for these $k$ (defined as $S_k\subseteq N$)
and by the evident quotient maps otherwise. These functors form an
inverse system; for $k\in I_{N'}\setminus I_N$ the non-invertible
arrows $0_k\to1$ become automorphisms of $1$. Let
$\cC(G;H;S_\bullet)=\varinjlim_N\Fin^{\cX_N}$.

\begin{thm}\label{thm:collapse}
Let $L$ prove linearity with $n_L\ge3$ and $\cC=\cC(G;H;S_\bullet)$.
Then $\HH_L\cC$ is a pretopos if and only if for every cofinite
$K\subseteq\mathbb N$
\begin{equation}\label{eq:cof}
H\cap\overline{\sum_{k\in K}S_k}=\overline{\sum_{k\in K}(H\cap S_k)}.
\end{equation}
\end{thm}
\begin{proof}
We use: for directed decreasing families of closed subsets of $G$,
intersections commute with finite sums and with binary intersections;
a decreasing family of subsets of a finite set is eventually constant;
and $\overline X=\bigcap_N(X+N)$ for a subgroup $X$.

\emph{Strings of $\cX_N$.} They have at most three objects. For a
string $s$ whose first object is $0_k$ or $2$ the group $G_s$ is
trivial, so the units are injective at $s$ by (R11). For $s=(1)$ the
one-step arrows into $s$ are the $x\colon(0_k\xrightarrow{x}1)\to s$,
and the relations $x^*(=)$ on $V_s=G_s$ and on $K_s=H_s=G_s$ coincide;
so $E_T=F_T$ for every leg tree at $s$, and the units are injective at
$s$ by (R10). There remain the strings
$t_\gamma=(1\xrightarrow{\gamma}2)$, $\gamma\in G/(H+N)$. An
automorphism of the string $t_\gamma$ is $g\in G/N$ with $\gamma+g=
\gamma$; so $G_{t_\gamma}=G/N$ and $H_{t_\gamma}=(H+N)/N$. The
one-step arrows into $t_\gamma$ are the arrows $y\colon(0_k
\xrightarrow{x}1\xrightarrow{\gamma'}2)\to t_\gamma$, $y=g+x\in
G/(S_k+N)$ with $\gamma+g=\gamma'$; every $y\in G/(S_k+N)$ occurs, and
these strings have no one-step arrows into them. Over $s=t_\gamma$,
for $u=(0_k\to1\to2)$ we have $V_u=\cX_N(0_k,1)$ with the equality,
and the leg $y$ gives on $V_s=G/N$ the relation $g+y=g'+y$, that is
$g-g'\in(S_k+N)/N$, and on $K_s=(H+N)/N$ its restriction. So for a leg
tree $T$ at $t_\gamma$, with $J$ the set of the $k$ such that the root
has a child with first object $0_k$,
\begin{equation}\label{eq:SR}
S_T=\sum_{k\in J}(S_k+N)/N,\qquad R_T=\sum_{k\in J}\bigl((H+N)\cap
(S_k+N)\bigr)/N.
\end{equation}

\emph{The pullback of a functor.} Let $W\colon\cS_n(\cX_N)\to\Fin$,
$N'\subseteq N$, $W'=W\tilde q$, $t'=(1\xrightarrow{\gamma'}2)$ a
string of $\cX_{N'}$ over $t=\tilde q(t')=(1\xrightarrow{\gamma}2)$,
and $a\in W'(t')=W(t)$. Let $J(a)\subseteq I_N$ be the set of the $k$
such that $a=W(\bar y)b$ for a one-step arrow $\bar y$ into $t$ from a
string with first object $0_k$ and some $b$. We determine the first
objects of the children of the root of $\Max(t',a)$ for $W'$. Let
$y'\colon u'\to t'$ be one-step, $u'=(0_k\xrightarrow{x}1\to2)$. If
$k\in I_{N'}\setminus I_N$, then $q(x)$ is invertible, $\tilde q(u')$
has two objects and $q(y')\colon\tilde q(u')\to t$ is an isomorphism
of strings; so $W'(y')=W(q(y'))$ is bijective and $a$ is in its image:
these are the collapsing children of Lemma~\ref{lem:pull}. If $k\in
I_N$, then $\tilde q(u')=q(u')$ and $q(y')$ is a one-step arrow into
$t$, and every one-step arrow $\bar y$ into $t$ from a string with
first object $0_k$ is $q(y')$ for such a $y'$, the quotient maps being
surjective; so some child has first object $0_k$ iff $k\in J(a)$. By
\eqref{eq:SR} for $\cX_{N'}$ and (R11), the unit of $W'$ is injective
at $t'$ iff every $a\in W(t)$ is fixed by $W(\theta)$ for $\theta$ in
\[
B_{N'}(a)=\text{the image in }G/N\text{ of }(H+N')\cap D_{N'}(a),
\]
where $D_{N'}(a)=N'+\sum_{k\in J(a)\cup(I_{N'}\setminus I_N)}S_k$.
For $N''\subseteq N'$ we have $D_{N''}(a)\subseteq D_{N'}(a)$, since
$S_k+N''\subseteq N'$ for $k\in I_{N''}\setminus I_{N'}$. As
$S_k\subseteq N'$ for $k\notin I_{N'}$, $D_{N'}(a)$ is $\sum_{k\in
K(a)}S_k+N'$ with $K(a)=J(a)\cup(\mathbb N\setminus I_N)$, a cofinite
set; the intersection over $N'$ of these groups is
$\overline{\sum_{k\in K(a)}S_k}$. Hence for all sufficiently small
$N'$, $B_{N'}(a)$ is the image $B_\infty(a)$ of $H\cap\overline{
\sum_{k\in K(a)}S_k}$. With the first part and
Lemma~\ref{lem:eventual}: $\HH_L\cC$ is a pretopos iff for every $N$,
$W$, $\gamma$ and $a\in W(t_\gamma)$ the group $B_\infty(a)$ fixes
$a$.

\emph{Sufficiency.} Assume \eqref{eq:cof}. Then $B_\infty(a)$ is the
image of $\overline{\sum_{k\in K(a)}(H\cap S_k)}$, which, $G/N$ being
finite and $S_k\subseteq N$ for $k\notin I_N$, is the image of
$\sum_{k\in J(a)}(H\cap S_k)$, contained in $R_T$ for the leg tree $T$
at $t_\gamma$ with one child for each $k\in J(a)$, labelled by a
one-step arrow with $a$ in the image; $T$ is realised in $W$ with $a$
at the root, so $R_T$ fixes $a$ by (R10).

\emph{Necessity.} Let $K$ be cofinite and $x\in H\cap\overline{
\sum_{k\in K}S_k}$ outside $\overline{\sum_{k\in K}(H\cap
S_k)}=\bigcap_{N''}(\sum_{k\in K\cap I_{N''}}(H\cap S_k)+N'')$. Choose
$N''$ with $x\notin\sum_{k\in K\cap I_{N''}}(H\cap S_k)+N''$ and, the
$S_k$ being non-zero and $K$ cofinite, with $\mathbb N\setminus
I_{N''}\subseteq K$. The closed sets $\sum_{k\in K\cap I_{N''}}((H+N)
\cap(S_k+N))+N''$, $N\subseteq N''$, decrease to $\sum_{k\in K\cap
I_{N''}}(H\cap S_k)+N''$; so there is $N\subseteq N''$ such that $x$
is outside one of them. Let $J=K\cap I_N$. For $k\in J\setminus
I_{N''}$ we have $S_k+N\subseteq N''$; so
\[
x\notin N+\sum_{k\in J}\bigl((H+N)\cap(S_k+N)\bigr),
\]
the inverse image of $R_T$ for the leg tree $T$ of $\cX_N$ at
$t_0=(1\xrightarrow{0}2)$ with one child for each $k\in J$, labelled
by $0\colon(0_k\xrightarrow{0}1\xrightarrow{0}2)\to t_0$. Let $W=W_T$
as in (R10), with $a$ at the root: $W(\bar x)a\ne a$ for the image
$\bar x$ of $x$ in $(H+N)/N$. But $J(a)\supseteq J$, so $K(a)\supseteq
J\cup(\mathbb N\setminus I_N)=K$, as $\mathbb N\setminus I_N\subseteq
\mathbb N\setminus I_{N''}\subseteq K$; hence $\bar x\in B_\infty(a)$,
and $\HH_L\cC$ is not a pretopos.
\end{proof}

This proves Theorem~\ref{thm:D}.

\begin{ex}\label{ex:conv}
Let $G=(\mathbb Z/2)^{\mathbb N}$, $S_k=\langle e_k\rangle$ ($e_k$ with
$k$-th coordinate $1$ and the others $0$) and $H=\langle h\rangle$,
$h=(1,1,1,\dots)$. The elements with finitely many non-zero
coordinates are dense, so for $K=\mathbb N$ the left side of
\eqref{eq:cof} is $H$, and the right side is $0$; so $c^L_\cC$ is not
an equivalence for $n_L\ge3$. For finite $J$ both sides are $0$. In
terms of Section~\ref{sec:inverse}: at the stage $N=G$ the category
$\cX_G$ has the objects $1$ and $2$ only. For the string $(1\to2)$ and
the tree $T$ without children we have $S_T=R_T=0$, while the pullbacks
of $T$ consist of collapsing children, $S_{(G,T)}$ is the closed
subgroup generated by all $S_k$, and $H\cap S_{(G,T)}=H\ne0=
R_{(G,T)}$. At each finite stage, for all but finitely many $k$ the
homomorphisms $0_k\to1$ have become automorphisms of the model $1$.
\end{ex}

\section{The dependence on the logic}\label{sec:depth}

The criterion depends on $L$ through $n_L$ only. We show that it
becomes weaker when $n_L$ decreases, that every $n_L\ge3$ is the exact
bound for some finite category, and that the criterion for
$n_L=\infty$ is not the conjunction of those for the finite $n_L$.

\subsection*{Monotonicity}

\begin{lem}\label{lem:slice}
Let $\cX$ be a finite EI-category, $n=n_L$ and $m\le n$ a positive
integer. Then $\cS_m(\cX)$ is a full subcategory of $\cS_n(\cX)$
closed under the arrows out of its objects; so there is a subfunctor
$\tau_m\le1$ in $\Fin^{\cS_n(\cX)}$ with $\tau_m(s)=1$ iff $s$ has at
most $m$ objects, it is an object of $\HH_L\Fin^\cX$, and restriction
to $\cS_m(\cX)$ is an equivalence from $(\HH_L\Fin^\cX)/\tau_m$ to the
full subcategory of $\Fin^{\cS_m(\cX)}$ of the functors isomorphic to
subfunctors of the functors $A\pi$, $A$ in $\Fin^\cX$.
\end{lem}
\begin{proof}
The arrows of $\cS_m$ and $\cS_n$ between strings with at most $m$
objects are defined in the same way, and the target of an arrow $s\to
t$ is isomorphic to a final segment of $s$, so it has at most as many
objects as $s$. Hence $\tau_m$ is a functor, a subfunctor of
$1=\pi^*1$, and an object of $\HH_L\Fin^\cX$ by (R9). Restriction is
an equivalence from the functors on $\cS_n$ which are empty outside
$\cS_m$ to the functors on $\cS_m$, the inverse being the extension by
the empty set, which is a functor because $\cS_m$ is closed under the
arrows out of its objects. A functor $W$ on $\cS_n$, empty outside
$\cS_m$, is isomorphic to a subfunctor of $A\pi$ iff its restriction
is isomorphic to a subfunctor of the restriction of $A\pi$: an
injective transformation on $\cS_m$ extends by the empty maps.
\end{proof}

\begin{thm}\label{thm:mono}
Let $(\cX_i)_{i\in I}$ be an inverse system of finite EI-categories
over a directed set, $\cC=\varinjlim\Fin^{\cX_i}$, and $L$, $L'$
logics proving linearity with $n_{L'}\le n_L$. If $\HH_L\cC$ is a
pretopos, so is $\HH_{L'}\cC$.
\end{thm}
\begin{proof}
Let $n=n_L$, $m=n_{L'}$, $\cE_i=\HH_L\Fin^{\cX_i}$ and
$\cE=\HH_L\cC=\varinjlim\cE_i$ by (R3). By (R9) and
Lemma~\ref{lem:collapse}(c), $\cE_i$ is the category of the functors
on $\cS_n(\cX_i)$ isomorphic to subfunctors of the $A\pi$, and the
transition functors are $W\mapsto W\tilde q_{i'i}$; the same holds for
$L'$ and $\cS_m$. If $m=n$, $\HH_{L'}\cC\simeq\cE$ and we are done.
Let $m<n$, so $m$ is finite, and $\tau_i=\tau_m\le1$ in $\cE_i$ as in
Lemma~\ref{lem:slice}. For $i'\ge i$ the functor $\tilde q_{i'i}$ maps
strings with at most $m$ objects to such strings, and its restriction
to $\cS_m(\cX_{i'})$ is the functor $\tilde q_{i'i}$ of $\cS_m$. So
$\tau_{i'}\le\tau_i\tilde q_{i'i}$, the image of $\tau_i$ in
$\cE_{i'}$, and under the equivalences of Lemma~\ref{lem:slice} the
functor
\[
\cE_i/\tau_i\to\cE_{i'}/\tau_{i'},\qquad W\mapsto W\tilde
q_{i'i}\times\tau_{i'},
\]
corresponds to $W\mapsto W\tilde q_{i'i}$ on the functors on
$\cS_m(\cX_i)$, that is to $\HH_{L'}(q_{i'i}^*)$. By (R3) for $L'$,
\[
\HH_{L'}\cC\simeq\varinjlim_i\cE_i/\tau_i .
\]
Let $t_i\le1$ be the image of $\tau_i$ in $\cE$; then $t_{i'}\le t_i$
for $i'\ge i$. For $j\le i$ let $\tau_{j,i}$ be the image of $\tau_j$
in $\cE_i$, so $\tau_i\le\tau_{j,i}$. The slice $\cE/t_j$ is the
colimit over $i\ge j$ of the $\cE_i/\tau_{j,i}$: an object of $\cE$
whose map to $1$ factors through $t_j$ is represented at some stage
$i\ge j$ by an object whose map to $1$ factors through $\tau_{j,i}$,
and the slices are full subcategories. So $\varinjlim_j\cE/t_j$, along
the functors $Z\mapsto Z\times t_{j'}$, is the colimit of the
categories $\cE_i/\tau_{j,i}$ over the directed set of the pairs
$j\le i$, along the functors $Z\mapsto Z|_{i'}\times\tau_{j',i'}$ for
$j\le j'$, $i\le i'$. The pairs $(i,i)$ are cofinal, and on them this
is the system of the $\cE_i/\tau_i$. Hence
\[
\HH_{L'}\cC\simeq\varinjlim_j\cE/t_j .
\]
If $\cE$ is a pretopos, this is a directed colimit of pretoposes and
coherent functors (Section~\ref{sec:setting}, Localisation), hence a
pretopos.
\end{proof}

\begin{cor}\label{cor:mono}
Let $\cC$ be a coherent category with finite sorts, $L$ a
non-classical logic proving linearity and $L'$ a logic proving
linearity with $n_{L'}\le n_L$. If $c^L_\cC$ is an equivalence, so is
$c^{L'}_\cC$. In particular, if the translation into G\"odel--Dummett
logic commutes with the imaginaries of $\cC$, so do the translations
into the logics $L_m$ of the finite chains.
\end{cor}
\begin{proof}
We may assume that $\cC$ is a pretopos (Section~\ref{sec:setting});
then $\HH_L\cC$ is a pretopos, $\Mod(\cC)$ is an EI-category by (R4),
and Theorem~\ref{thm:mono} applies to the inverse system of
Section~\ref{sec:theories}.
\end{proof}

So for a finite EI-category $\cX$ the set of the $n$ such that
$\HH_L\Fin^\cX$ is a pretopos for $n_L=n$ is an initial segment of
$\{1,2,\dots,\infty\}$; it contains $1$ and $2$
\cite[Cor.~6.3]{Strings}.

\subsection*{Cone categories}
Let $Z=\{1,\omega\}$ be the group with two elements, $P$ a finite
poset and $\Phi\colon P^{op}\to\Fin$ a functor with an action of $Z$
on every $\Phi(y)$ which is free and commutes with the restriction
maps $x\mapsto x|_{y'}$, $y'\le y$. The \emph{cone category}
$\cX(P,\Phi)$ has the objects of $P$ and one more object $c$; the
arrows of $P$ (one arrow $y'\to y$ iff $y'\le y$), $\cX(y,c)=\Phi(y)$,
$\Aut(c)=Z$, and no arrows from $c$ to $P$; the composites are given
by the restrictions and the actions. It is a finite EI-category, and
its non-invertible arrows are the arrows $y'<y$ and the elements of
the $\Phi(y)$. For $y\in P$ and $k\ge0$ define equivalence relations
$E^k_y$ on $\Phi(y)$ by $E^0_y=$ equality and, for $k\ge1$,
\[
E^k_y=\bigvee_{y'<y}\{(x_1,x_2)\mid x_1|_{y'}\mathrel{E^{k-1}_{y'}}
x_2|_{y'}\}
\]
(the equality if $y$ is minimal). They increase with $k$; let
$E^\infty_y=\bigcup_kE^k_y$ and $\infty-2=\infty$.

\begin{lem}\label{lem:cone}
Let $\cX=\cX(P,\Phi)$ and $n=n_L\ge2$. Then $\HH_L\Fin^\cX$ is a
pretopos if and only if there are no $y\in P$ and $x\in\Phi(y)$ with
$x\mathrel{E^{n-2}_y}\omega x$.
\end{lem}
\begin{proof}
\emph{Strings.} An arrow out of $c$ is invertible; so a string is a
chain $y_0<\dots<y_r$ of $P$, or such a chain followed by an arrow
$x\in\Phi(y_r)$, for which we write $(y_0,\dots,y_r;x)$, a string
with $r+2\le n$ objects, or $(c)$. For a string $s$ with first object
in $P$ the group $G_s$ is trivial and the units are injective at $s$
by (R11). There remains $s=(c)$, with $G_s=H_s=Z$; we take it as the
base string. The strings $t\ne(c)$ over it are the
$t=(y_0,\dots,y_r;x)$: an arrow $t\to(c)$ of $\cS_n$ is $\theta p$
with $p=x|_{y_0}$ the composite of $t$ and $\theta\in Z$; so
$V_t=\Phi(y_0)$ and $K_t=\{x|_{y_0},\omega x|_{y_0}\}$, two different
elements. An isomorphism of strings onto $t$ is the identity on the
$y_j$, so its source is $t$ or $(y_0,\dots,y_r;\omega x)$. Hence the
one-step arrows into $t$ are the arrows $y'<y_0$ of $\cX$, as arrows
from $(y',y_0,\dots,y_r;x)$ and from $(y',y_0,\dots,y_r;\omega x)$ (if
$r+3\le n$), and the one-step arrows into $(c)$ are the $x\in\Phi(y)$,
$y\in P$, as arrows from $(y;x)$ and from $(y;\omega x)$.

\emph{The relations $F$.} For every leg tree $T$ at a string $t\ne(c)$
over $(c)$, $F_T$ is the equality on the two-element set $K_t$: by
induction, for a child $(\beta,T')$, $\beta$ the arrow $y'<y_0$, the
map $g\mapsto g\beta=g|_{y'}$ from $K_t$ to $K_{t'}$ is injective, the
action being free, and $F_{T'}$ is the equality. The same holds at
$(c)$: for a child $(x,T')$ the map $\theta\mapsto\theta x$ from $Z$
to $K_{t'}$ is injective. So $R_T=1$ for every leg tree at $(c)$.

\emph{The relations $E$.} For a leg tree $T$ of depth at most $k$ at
$t=(y_0,\dots,y_r;x)$ we have $E_T\subseteq E^k_{y_0}$, by induction:
$E_T$ is the join of the relations $\beta^*E_{T'}$, $\beta$ an arrow
$y'<y_0$ and $T'$ of depth at most $k-1$ at a string with first object
$y'$, and $\beta^*E_{T'}$ is the set of the $(x_1,x_2)$ with
$x_1|_{y'}\mathrel{E_{T'}}x_2|_{y'}$. Conversely, if $r+2+k\le n$, the
leg tree $T^k_t$ at $t$ whose root has all one-step arrows into $t$ as
children, with the subtrees $T^{k-1}_{t'}$ at their sources (no
children if $k=0$), has $E_{T^k_t}=E^k_{y_0}$, by the same induction.
A leg tree at $(y;x)$ has depth at most $n-2$.

\emph{Conclusion.} For a leg tree $T$ at $(c)$, $E_T$ is a relation on
the two-element set $Z$; so $\omega\in S_T$ iff for some child
$(x,T')$ of the root we have $x\mathrel{E_{T'}}\omega x$, $T'$ being a
leg tree at $(y;x)$ or $(y;\omega x)$. If there are no $y$, $x$ with
$x\mathrel{E^{n-2}_y}\omega x$, then $S_T=1$ for all $T$ and the units
are injective at $(c)$ by (R11). If $x\mathrel{E^{n-2}_y}\omega x$,
let $k\le n-2$ be finite with $x\mathrel{E^k_y}\omega x$ and $T$ the
leg tree at $(c)$ with the one child $(x,T^k_{(y;x)})$; then
$\omega\in S_T$ and $R_T=1$. For $W=W_T$ and $a$ as in (R10),
$W(\omega)a\ne a$, and there is a morphism $T\to\Max((c),a)$, so
$\omega\in S_{\Max((c),a)}$ by (R11): the unit of $W$ is not injective
at $(c)$. With (R9) this proves the lemma.
\end{proof}

\subsection*{The categories $\cX_d$}
Let $d\ge1$. Let $P_d$ be the poset with the elements $b_0,\dots,b_d$,
$a'$, $z^k_1,z^k_2$ ($1\le k\le d$) and the order generated by
\[
b_d<b_{d-1}<\dots<b_0,\qquad a'<b_0,\qquad z^k_i<b_k .
\]
Let $\Gamma$ be the set of the symbols $\beta$, $e_0,\dots,e_d$,
$e'_0,\dots,e'_{d-1}$; put $e'_d=e_d$, and $\Gamma_{\le j}=\{\beta\}
\cup\{e_k,e'_k\mid k\le j\}$. Let
\begin{align*}
\Phi_d(b_j)&=Z\times\Gamma_{\le j},\\
\Phi_d(z^k_1)&=Z\times\Gamma_{\le k}/\bigl((\theta,e_{k-1})\sim
(\theta,e_k)\bigr),\\
\Phi_d(z^k_2)&=Z\times\Gamma_{\le k}/\bigl((\theta,e'_{k-1})\sim
(\theta,e'_k)\bigr),\\
\Phi_d(a')&=Z\times\Gamma_{\le0}/\bigl((\theta,\beta)\sim(\theta,e_0),
\ (\theta,e'_0)\sim(\omega\theta,\beta)\bigr),
\end{align*}
with the action of $Z$ on the first factor, the inclusions
$\Phi_d(b_j)\to\Phi_d(b_l)$ ($l\ge j$, so $b_l\le b_j$) as
restrictions, and the inclusions followed by the quotient maps as the
other restrictions. The equivalence relations are invariant under $Z$
and no class contains $(\theta,e)$ and $(\omega\theta,e)$ (the classes
of $\Phi_d(a')$ are $\{(\theta,\beta),(\theta,e_0),(\omega\theta,
e'_0)\}$); so the actions are free. Let $\cX_d=\cX(P_d,\Phi_d)$. We
write $\theta e$ for $(\theta,e)$.
\[
\begin{tikzcd}[column sep=small,row sep=small]
z^d_1,z^d_2\ar[d]&&z^2_1,z^2_2\ar[d]&z^1_1,z^1_2\ar[d]&a'\ar[d]&\\
b_d\ar[r]&\cdots\ar[r]&b_2\ar[r]&b_1\ar[r]&b_0\ar[r,"{\beta,e_0,e'_0}"]&c
\ar[loop right,"\omega"]
\end{tikzcd}
\]
For $d=1$ this is the category of \cite[Ex.~6.19]{Strings}, with
$e_0$, $e'_0$, $e_1$, $\omega$ for $g$, $g'$, $h$, $\sigma$ there.

\begin{thm}\label{thm:Xd}
Let $L$ prove linearity and $d\ge1$. Then $\HH_L\Fin^{\cX_d}$ is a
pretopos if and only if $n_L\le d+2$.
\end{thm}
\begin{proof}
If $n_L=1$, the strings have one object and there are no one-step
arrows; so the trees $\Max(s,a)$ have no children, and the units are
injective by (R11). Let $n_L\ge2$. We compute the relations $E^k_y$.
At the minimal elements they are the equality. Let $\ell_j$ be the
equivalence relation on $\Phi_d(b_j)$ generated by $\theta
e_{j-1}\sim\theta e_j$, $\theta e'_j\sim\theta e'_{j-1}$ if $j\ge1$,
and by $\theta\beta\sim\theta e_0$, $\theta e'_0\sim\omega\theta\beta$
if $j=0$ ($\theta\in Z$).

\emph{Claim.} For $k\ge1$ and $0\le j\le d$, $E^k_{b_j}$ is generated
by $\ell_j$ and, if $j<d$ and $k\ge d-j+1$, the pairs $\theta
e_j\sim\theta e'_j$.

Proof by induction on $k$. The elements below $b_j$ are $z^j_1,z^j_2$
(for $j\ge1$), $a'$ (for $j=0$), and the $b_l$, $z^l_i$ with $l>j$. As
$E^{k-1}$ is the equality at the minimal elements, they contribute the
kernels of the restrictions. The kernels of the restrictions to
$z^j_1$, $z^j_2$ (to $a'$ if $j=0$) generate $\ell_j$. For $l>j$ the
restriction to $z^l_i$ is injective on $\Phi_d(b_j)$: the non-trivial
classes of $\Phi_d(z^l_1)$ are the $\{\theta e_{l-1},\theta e_l\}$, and
$e_l\notin\Gamma_{\le j}$ ($e_l\ne e'_k$ for $k\le j<l$, also if
$l=d$); likewise for $z^l_2$. The element $b_l$, $l>j$, contributes
the restriction of $E^{k-1}_{b_l}$ to the subset $\Phi_d(b_j)$. For
$k=1$ this is the equality. For $k\ge2$, by the induction hypothesis
every class of $E^{k-1}_{b_l}$ with more than one element is contained
in $\{\theta e_{l-1},\theta e_l,\theta e'_l,\theta e'_{l-1}\}$ for one
$\theta$; the elements of this set in $\Phi_d(b_j)$ are $\theta
e_{l-1}$, $\theta e'_{l-1}$ if $l=j+1$, and none if $l\ge j+2$. For
$l=j+1$: $\theta e_j$ and $\theta e'_j$ are related by
$E^{k-1}_{b_{j+1}}$ iff $\theta e_{j+1}$ and $\theta e'_{j+1}$ are
equal or related by one of the generators, that is, iff $j+1=d$, or
$j+1<d$ and $k-1\ge d-j$. As $k\ge2$, this is the condition $k\ge
d-j+1$. This proves the claim.

By the claim, for $j\ge1$ every class of $E^k_{b_j}$ is contained in
some $\{\theta\}\times\Gamma$; so $x$ and $\omega x$ are not related.
For $j=0$: if $1\le k\le d$, the classes of $E^k_{b_0}$ are the
$\{\theta\beta,\theta e_0,\omega\theta e'_0\}$ and singletons; if
$k\ge d+1$, $\beta\sim e_0\sim e'_0\sim\omega\beta$. So there are $y$,
$x$ with $x\mathrel{E^{k}_y}\omega x$ iff $k\ge d+1$, and
Lemma~\ref{lem:cone} gives: $\HH_L\Fin^{\cX_d}$ is a pretopos iff
$n_L-2\le d$.
\end{proof}

\subsection*{The tower}
For $d\ge1$ let $q_d\colon\cX_{d+1}\to\cX_d$ be the functor which is
the identity on $c$, $a'$, the $b_j$ and $z^k_i$ with $j,k\le d$,
sends $b_{d+1}$, $z^{d+1}_1$, $z^{d+1}_2$ to $b_d$, is the identity on
$Z$, and is given on the arrows into $c$ by $q_d(\theta e)=\theta
\varepsilon(e)$, where $\varepsilon$ maps the symbols of $\cX_{d+1}$
to those of $\cX_d$ by
\[
\varepsilon(e_d)=\varepsilon(e'_d)=\varepsilon(e_{d+1})=e_d
\]
and the identity on the other symbols. This is well defined and a
functor: the map on the objects is monotone; $\varepsilon$ maps the
symbols $\Gamma_{\le j}$ of $\cX_{d+1}$ into those of $\cX_d$ for
$j\le d$ and $\Gamma_{\le d+1}$ into $\Gamma_{\le d}$; it sends the
generating pairs of the relations for $z^k_i$, $k<d$, and $a'$ to the
same pairs, the pair $(\theta e_{d-1},\theta e_d)$ for $z^d_1$ to
itself, the pair $(\theta e'_{d-1},\theta e'_d)$ for $z^d_2$ to
$(\theta e'_{d-1},\theta e_d)$, the pair of $z^d_2$ in $\cX_d$, and
the pairs for $z^{d+1}_i$ to pairs of equal elements of $\Phi_d(b_d)$;
so the maps $\Phi_{d+1}(y)\to\Phi_d(q_dy)$ are defined, equivariant
and compatible with the restrictions. The $q_d$ form an inverse system
indexed by the positive integers; for $d'\ge d$ the composite
$q_{d'd}\colon\cX_{d'}\to\cX_d$ sends $b_j$ ($j\ge d$) and $z^k_i$
($k>d$) to $b_d$, and $e_k$, $e'_k$ ($k\ge d$) to $e_d$. Let
\[
\cC_\infty=\varinjlim_d\Fin^{\cX_d},
\]
the colimit along the functors $q_d^*$.

\begin{thm}\label{thm:tower}
Let $L$ prove linearity. Then $\HH_L\cC_\infty$ is a pretopos if and
only if $n_L$ is finite.
\end{thm}
\begin{proof}
Let $n_L$ be finite. For $d'\ge n_L-2$, $\HH_L\Fin^{\cX_{d'}}$ is a
pretopos by Theorem~\ref{thm:Xd}, so the unit of every functor
$\cS_n(\cX_{d'})\to\Fin$ is injective by (R9), in particular that of
$W\tilde q_{d'd}$ for $W$ at a stage $d\le d'$;
Lemma~\ref{lem:eventual} applies.

Let $n_L=\infty$. In $\cX_1$ let $T$ be the leg tree at $(c)$ whose
root has the one child $\beta\colon(b_0;\beta)\to(c)$, which has the
children $a'<b_0$, from $(a',b_0;\beta)$, without children, and
$b_1<b_0$, from $t_1=(b_1,b_0;\beta)$, with the two children
$z^1_i<b_1$, from $(z^1_i,b_1,b_0;\beta)$, without children; let $T_1$
be its subtree at $t_1$. Let $W=W_T$ and $a$ be as in (R10); as
$R_T=1$ (proof of Lemma~\ref{lem:cone}), $W(\omega)a\ne a$. Fix
$d\ge1$, let $q=q_{d1}$ and $T^{(c)}$ the pullback along $q$, a leg
tree of $\cX_d$ at $(c)$. We show $\omega\in S_{T^{(c)}}$. Then, as
there is a morphism $T\to\Max((c),a)$, by (P4), (P1) and (R11)
$\omega$ is in $S_{\Max((c),a)^{(c)}}$, the group of the tree
$\Max((c),a)$ for $W\tilde q$, while $W\tilde
q(\omega)a=W(\omega)a\ne a$; so the unit of $W\tilde q_{d1}$ is not
injective, for every $d$, and $\HH_L\cC_\infty$ is not a pretopos by
Lemma~\ref{lem:eventual}.

Let $U$ be the leg tree of $\cX_d$ at $(c)$ whose nodes other than the
root are: $\beta$, from $u_0=(b_0;\beta)$; above it $a'<b_0$, from
$(a',b_0;\beta)$; and for $1\le k\le d$ the node $b_k<b_{k-1}$, from
$u_k=(b_k,\dots,b_0;\beta)$, above the node of $u_{k-1}$, and above it
the two nodes $z^k_i<b_k$, from $(z^k_i,b_k,\dots,b_0;\beta)$.

\emph{$U$ is contained in $T^{(c)}$}, as a set of nodes closed under
parents. We use the definition of the pullback in
Section~\ref{sec:pull}; note that $\tilde q(u_0)=(b_0;\beta)$ and
$\tilde q(u_k)=t_1$ for $k\ge1$, as $q$ maps the arrows $b_k<b_{k-1}$,
$k\ge2$, to the identity of $b_1$. The arrow $\beta$ of $\cX_d$ lies
over the child $\beta$ of $T$; above it, the subtree is the pullback
at $u_0$, in which $a'<b_0$ lies over the child $a'<b_0$ and $b_1<b_0$
over the child $b_1<b_0$, with subtree $T_1^{u_1}$. In $T_1^{u_1}$ the
arrows $z^1_i<b_1$ lie over the children of $T_1$. For $1\le k<d$: in
$T_1^{u_k}$ the one-step arrow $b_{k+1}<b_k$ from $u_{k+1}$ is mapped
by $q$ to the identity of $t_1$, so it is a collapsing child with
subtree $T_1^{u_{k+1}}$; and for $2\le k\le d$ the one-step arrows
$z^k_i<b_k$ into $u_k$ are mapped to the identity of $t_1$ as well
($q(z^k_i)=b_1$), so they are collapsing children of the root of
$T_1^{u_k}$.

\emph{$\omega\in S_U$.} Let $E_k$ be the relation on $V_{u_k}=
\Phi_d(b_k)$ of the subtree of $U$ at $u_k$. The children $z^k_1$,
$z^k_2$ give $e_{k-1}\sim e_k$ and $e'_k\sim e'_{k-1}$. So
$e_{d-1}\mathrel{E_d}e'_{d-1}$, as $e_d=e'_d$; and for $1\le k<d$, if
$e_k\mathrel{E_{k+1}}e'_k$, the child $b_{k+1}$ gives $e_k\sim e'_k$
and so $e_{k-1}\mathrel{E_k}e'_{k-1}$. Hence $e_0\mathrel{E_1}e'_0$,
and at $u_0$, with the child $a'$: $\beta\sim e_0\sim e'_0\sim
\omega\beta$. So $\id\mathrel{E_U}\omega$, and $S_U\subseteq
S_{T^{(c)}}$ by (R11), the inclusion of a set of nodes closed under
parents being a morphism of leg trees.
\end{proof}

\begin{cor}\label{cor:tower}
There is a pretopos $\cC$ with finite sorts such that $c^L_\cC$ is an
equivalence for every logic $L$ proving linearity with $n_L$ finite,
in particular for the logics $L_m$ of the finite chains, and is not an
equivalence for G\"odel--Dummett logic.
\end{cor}
\begin{proof}
$\cC_\infty$ is a directed colimit of pretoposes and coherent
functors, hence a pretopos. An object of $\cC_\infty$ comes from some
$\Fin^{\cX_d}$, and a coherent functor $F\colon\Fin^{\cX_d}\to\Set$
has finite values: for an object $A$ with at most $N$ elements at
every object of $\cX_d$, the power $A^{N+1}$ is the union of the
subobjects on which two coordinates agree, and $F$ preserves this
finite union. Theorem~\ref{thm:tower} applies.
\end{proof}

Corollary~\ref{cor:mono}, Theorem~\ref{thm:Xd} for $d=m-2$ and
Corollary~\ref{cor:tower} give Theorem~\ref{thm:E}.

\begin{rem}\label{rem:tower}
(a) By Corollary~\ref{cor:mono} the criterion for G\"odel--Dummett
logic implies those for all the logics $L_m$; by
Corollary~\ref{cor:tower} it is strictly stronger than all of them
together. (b) The limit of the $\cX_d$ has the objects $c$, $a'$,
$b_j$, $z^k_i$ and one more object $b_\infty$ below all $b_j$. The
obstruction for $n_L=\infty$ is at the string $(c)$ with one object,
in accordance with Theorem~\ref{thm:bounded}; the depth of the pulled
back trees $T^{(c)}$ is not bounded. For the logics $L_m$ the trees
have depth less than $m$ and there is no obstruction; so the depth of
the pulled back trees in Theorem~\ref{thm:A} cannot be bounded. (c) The chain $b_\infty<\dots<b_1<b_0$ gives an
unbounded chain of prime filters of sentences of $\cC_\infty$, as it
must: if $\Mod(\cC)$ is an EI-category and the chains of prime filters
of a pretopos $\cC$ with finite sorts are bounded by $m$, then $\HH_L\cC$ is the free $L_m$-category
on $\cC$ for all $L$ with $n_L\ge m$ (Proposition~\ref{prop:Lm}), and
the criterion is the same for these $L$.
\end{rem}

\section{The groups of a leg tree in the limit category}\label{sec:intrinsic}

The groups $S_{(i,T)}$ and $R_{(i,T)}$ of Section~\ref{sec:inverse}
are defined through the categories $\cX_{i'}$, $i'\ge i$, and an arrow
of $\cX_{i'}$ need not be the image of an arrow of the limit $\cM$. In
this section we describe them in terms of $\cM$, the threads and the
functor $q_i\colon\cM\to\cX_i$: as intersections over the solutions of
finite level of a system of implications
(Theorem~\ref{thm:level}), and, if the strings are bounded, by a
recursion over a tree in $\cM$ (Theorem~\ref{thm:rec}). Two examples
show that the least solution, and the solutions by closed equivalence
relations, do not give the groups.

Let $(\cX_k)_{k\in I}$ be an inverse system of finite EI-categories
and arbitrary functors with limit $\cM$, as in
Section~\ref{sec:inverse}, and $n=n_L\le\infty$. We fix a thread
$\sigma=(s_k)$, the \emph{base}, an index $i$ and a leg tree $T$ of
$\cX_i$ at $s_i$. Finite sets are discrete, and the sets of arrows of
$\cM$ and the sets of threads carry the topology of the inverse
limit; an equivalence relation on a profinite set $V$ is \emph{open}
if it is an open subset of $V\times V$, that is, if it is the inverse
image of an equivalence relation on a finite quotient of $V$.

\subsection*{Thread arrows}
For threads $\upsilon=(u_k)$ and $\tau=(t_k)$ a \emph{thread arrow}
$\gamma\colon\upsilon\to\tau$ is an arrow $\gamma\colon\upsilon_0\to
\tau_0$ of $\cM$ such that $\gamma_k:=q_k(\gamma)$ is an arrow
$u_k\to t_k$ of $\cS_n(\cX_k)$ for every $k$. It is \emph{short} if
every $\gamma_k$ is a one-step arrow or an isomorphism of strings.
Let
\[
V_\tau=\cM(\tau_0,\sigma_0),\qquad K_\tau=\{\text{thread arrows }\tau
\to\sigma\}\subseteq V_\tau,
\]
closed subsets of the set of arrows of $\cM$.

\begin{lem}\label{lem:tarrows}
\begin{enumerate}
\item[(a)] Thread arrows compose, and $K_\sigma=H_\sigma$. A thread
arrow $\gamma\colon\upsilon\to\tau$ which is invertible in $\cM$ is
short, every $\gamma_k$ is an isomorphism of strings, and
$\gamma^{-1}$ is a thread arrow $\tau\to\upsilon$; we call $\gamma$
an \emph{isomorphism}.
\item[(b)] Let $\upsilon$, $\tau$ be bounded, the threads of the
strings $\tilde u$, $\tilde t$ of $\cM$
\textup{(Lemma~\ref{lem:threads}(b))}. Then the thread arrows
$\upsilon\to\tau$ are the arrows $\tilde u\to\tilde t$ of
$\cS_n(\cM)$, that is, the arrows $\tilde\iota_0\tilde p$ with
$\tilde p$ the composite $\tilde u_0\to\tilde u_j$ of arrows of
$\tilde u$ and $\tilde\iota$ an isomorphism of strings from $\tilde
u_{\ge j}$ onto $\tilde t$. If a short thread arrow $\upsilon\to\tau$
is not an isomorphism, $\tilde u$ has exactly one object more than
$\tilde t$.
\end{enumerate}
\end{lem}
\begin{proof}
(a) The $\cS_n(\cX_k)$ are categories, and $\cS_n(\cX_k)(s_k,s_k)=
H_{s_k}$, an arrow $\sigma_0\to\sigma_0$ of $\cM$ being invertible.
If $\gamma$ is invertible, $\gamma_k=\iota_0p$ is invertible in
$\cX_k$, so the composite $p$ of non-invertible arrows is an
identity, and $\gamma_k$ is an isomorphism of strings; its inverse is
an arrow of $\cS_n(\cX_k)$.

(b) Let $i_0$ be such that for $k'\ge k\ge i_0$ the functor $q_{k'k}$
sends the arrows of $u_{k'}$, $t_{k'}$ to the arrows of $u_k$, $t_k$.
Let $\gamma$ be a thread arrow. For $k\ge i_0$ we have
$\gamma_k=\iota_0p_k$, with $p_k$ the composite of the first $j$
arrows of $u_k$, where $j$ is the difference of the numbers of
objects of $\tilde u$ and $\tilde t$; so $p_k=q_k(\tilde p)$. The set
of the isomorphisms of strings $\iota\colon(u_k)_{\ge j}\to t_k$ with
$\iota_0p_k=\gamma_k$ is finite and non-empty, and these sets form an
inverse system over the $k\ge i_0$; an element of the limit is an
isomorphism of strings $\tilde\iota\colon\tilde u_{\ge j}\to\tilde t$
with $\tilde\iota_0\tilde p=\gamma$. Conversely, $q_k(\tilde\iota_0
\tilde p)$ is an arrow $u_k\to t_k$ of $\cS_n(\cX_k)$ for $k\ge i_0$,
and then for all $k$ by Lemma~\ref{lem:collapse}(c). If $\gamma$ is
short and not invertible, then $\gamma_k$ is not invertible for $k$
large (Lemma~\ref{lem:threads}(a)), so it is one-step, and $u_k$ has
one object more than $t_k$.
\end{proof}

\subsection*{Nodes, children and solutions}
Let $T'$ be a leg tree of $\cX_i$ at a string $t$ and $\delta\colon
u\to t$ an arrow of $\cS_n(\cX_i)$ which is one-step or an
isomorphism of strings. Let $\Delta(T',\delta)$ be the following set
of leg trees at $u$: the set of the $T_1$ with $(\delta,T_1)$ a child
of the root of $T'$, if $\delta$ is one-step; and $\{\delta^*T'\}$,
the transport of Section~\ref{sec:pull}, if $\delta$ is an
isomorphism of strings. Let $\mathcal T$ be the set of the leg trees
$\iota^*T_1$, where $T_1$ is the subtree of $T$ at one of its nodes
and $\iota$ an isomorphism of strings of $\cX_i$ onto the string of
this node. It is finite, contains $T$, and $\Delta(T',\delta)\subseteq
\mathcal T$ for $T'\in\mathcal T$: the children of the root of
$\iota^*T_1$ are the $(\iota^{-1}\beta,T_2)$ for the children
$(\beta,T_2)$ of the root of $T_1$, and $\delta^*\iota^*T_1=
(\iota\delta)^*T_1$.

A \emph{node} is a pair $x=(\tau,T')$ of a thread $\tau=(t_k)$ and
$T'\in\mathcal T$ at $t_i$. A \emph{child} of the node $(\tau,T')$ is
a triple $(\upsilon,\gamma,T'')$ of a thread $\upsilon$, a short
thread arrow $\gamma\colon\upsilon\to\tau$ and $T''\in\Delta(T',
\gamma_i)$; then $(\upsilon,T'')$ is a node. An \emph{$E$-solution}
is a family $\mathcal F=(\mathcal F_x)$ of equivalence relations on
the sets $V_\tau$, indexed by the nodes $x=(\tau,T')$, such that for
every child $(\upsilon,\gamma,T'')$ of a node $(\tau,T')$ and
$g,g'\in V_\tau$:
\begin{equation}\label{eq:sol}
g\gamma\mathrel{\mathcal F_{(\upsilon,T'')}}g'\gamma\quad
\Longrightarrow\quad g\mathrel{\mathcal F_{(\tau,T')}}g'.
\end{equation}
An \emph{$F$-solution} is defined in the same way with the sets
$K_\tau$ in place of the $V_\tau$ (if $g\in K_\tau$, then $g\gamma\in
K_\upsilon$). A solution is \emph{of level $k$} ($k\ge i$) if the
truth of $g\mathrel{\mathcal F_{(\tau,T')}}g'$ depends only on
$(t_k,T',q_k(g),q_k(g'))$, and \emph{of finite level} if it is of
some level; its relations are then open. For $k\ge i$ let
$\varepsilon^k$ and $\varphi^k$ be the families
\[
g\mathrel{\varepsilon^k_{(\tau,T')}}g'\iff q_k(g)\mathrel{E_{T'^{t_k}}}
q_k(g'),\qquad
h\mathrel{\varphi^k_{(\tau,T')}}h'\iff q_k(h)\mathrel{F_{T'^{t_k}}}
q_k(h'),
\]
where $T'^{t_k}$ is the pullback of $T'$ along $q_{ki}$ and the
relations are taken relative to the base string $s_k$, so on
$\cX_k((t_k)_0,(s_k)_0)$ and on $\cS_n(\cX_k)(t_k,s_k)$. They
decrease with $k$, by (P2) and (P3). Let
\[
\bar E=\bigcap_{k\ge i}\varepsilon^k,\qquad\bar F=\bigcap_{k\ge i}
\varphi^k .
\]
By the definitions, $S_{(i,T)}=\{\Theta\in G_\sigma\mid\id\mathrel{
\bar E_{(\sigma,T)}}\Theta\}$ and $R_{(i,T)}=\{\Theta\in H_\sigma\mid
\id\mathrel{\bar F_{(\sigma,T)}}\Theta\}$.

\begin{thm}\label{thm:level}
\begin{enumerate}
\item[(a)] $\varepsilon^k$ is an $E$-solution of level $k$ and
$\varphi^k$ is an $F$-solution of level $k$, for every $k\ge i$.
\item[(b)] For every $E$-solution $\mathcal F$ of finite level there
is $l$ with $\varepsilon^l_x\subseteq\mathcal F_x$ for all nodes $x$;
and the same for the $F$-solutions and the $\varphi^l$.
\item[(c)] $\bar E$ is the intersection of the $E$-solutions of
finite level and $\bar F$ that of the $F$-solutions of finite level.
So $S_{(i,T)}$ is the set of the $\Theta\in G_\sigma$ with
$\id\mathrel{\mathcal F_{(\sigma,T)}}\Theta$ for every $E$-solution
$\mathcal F$ of finite level, and $R_{(i,T)}$ the set of the
$\Theta\in H_\sigma$ with this property for every $F$-solution of
finite level.
\end{enumerate}
\end{thm}
\begin{proof}
(a) The relations are equivalence relations depending only on the
images at the stage $k$. Let $(\upsilon,\gamma,T'')$ be a child of
$(\tau,T')$; note $q_{ki}(\gamma_k)=\gamma_i$. If $\gamma_k$ is
one-step, then by the definition of the pullback $\gamma_k$ is the
label of a child of the root of $T'^{t_k}$ with subtree $T''^{u_k}$:
if $\gamma_i$ is one-step, $T''$ is the subtree of a child
$\nu=(\gamma_i,T'')$ of the root of $T'$ and the child is
$(\gamma_k,\nu)$; if $\gamma_i$ is an isomorphism of strings,
$T''=\gamma_i^*T'$ and $\gamma_k$ is a collapsing child. Hence
$\gamma_k^*E_{T''^{u_k}}\subseteq E_{T'^{t_k}}$ and
$\gamma_k^*F_{T''^{u_k}}\subseteq F_{T'^{t_k}}$, which is
\eqref{eq:sol}. If $\gamma_k$ is an isomorphism of strings, so is
$\gamma_i$, $T''=\gamma_i^*T'$, and $T''^{u_k}=\gamma_k^*(T'^{t_k})$
by \eqref{eq:tp}; by the description of the relations of a
transported tree in Section~\ref{sec:pull}, $g\gamma_k\mathrel{
E_{\gamma_k^*(T'^{t_k})}}g'\gamma_k$ iff $g\mathrel{E_{T'^{t_k}}}g'$,
and the same for $F$.

(b) We use the following fact: \emph{for an inverse system of finite
sets $(B_{k'})_{k'\ge k}$ over a directed set there is $l\ge k$ such
that the image of $B_l$ in $B_k$ is the image of $\varprojlim
B_{k'}$.} Indeed, the images of the $B_{k'}$ in $B_k$ decrease and
are constant for $k'\ge l$; for $b$ in this image the sets of the
preimages of $b$ in the $B_{k'}$ are finite and non-empty and form an
inverse system, which has a non-empty limit.

Let $\mathcal F$ be an $E$-solution of level $k$. For $k'\ge k$ let
$\Pi^3_{k'}$ be the set of the tuples $(t,T',g_1,g_2,g_3)$ of a
string $t$ of $\cX_{k'}$, a tree $T'\in\mathcal T$ at $\tilde q_{k'i}
(t)$ and $g_1,g_2,g_3\in\cX_{k'}(t_0,(s_{k'})_0)$; and $\Pi'_{k'}$
the set of the tuples $(u,t,\gamma,T',T'',g_1,g_2)$ of strings $u$,
$t$ of $\cX_{k'}$, an arrow $\gamma\colon u\to t$ of $\cS_n(\cX_{k'})$
which is one-step or an isomorphism of strings, $T'\in\mathcal T$ at
$\tilde q_{k'i}(t)$, $T''\in\Delta(T',q_{k'i}(\gamma))$ and
$g_1,g_2\in\cX_{k'}(t_0,(s_{k'})_0)$. These sets are finite, and the
functors $q$ and $\tilde q$ map the sets at $k''$ to the sets at $k'$
for $k''\ge k'$, keeping the trees: by Lemma~\ref{lem:collapse}(c),(d)
the image of a one-step arrow or an isomorphism of strings is one or
the other, and $\Delta(T',q_{k''i}(\gamma))$ depends only on the image
at the stage $i$. An element of $\varprojlim\Pi^3_{k'}$ is a tuple
$(x,g_1,g_2,g_3)$ of a node $x=(\tau,T')$ and three elements of
$V_\tau$ (a family of strings $t_{k'}$, $k'\ge k$, extends uniquely
to a thread, and the arrows of $\cM$ are the compatible families); an
element of $\varprojlim\Pi'_{k'}$ is a tuple $(x,(\upsilon,\gamma,
T''),g_1,g_2)$ of a node $x=(\tau,T')$, a child of it and
$g_1,g_2\in V_\tau$. Choose $l\ge k$ as in the fact for both systems.

Let $\bar{\mathcal F}$ be the set of the tuples $(t_k,T',q_k(g),
q_k(g'))$ for the nodes $(\tau,T')$ and $g\mathrel{\mathcal F_{(\tau,
T')}}g'$; as $\mathcal F$ is of level $k$, $g\mathrel{\mathcal
F_{(\tau,T')}}g'$ iff $(t_k,T',q_k(g),q_k(g'))\in\bar{\mathcal F}$.
For a string $t$ of $\cX_l$ and $T'\in\mathcal T$ at $\tilde q_{li}
(t)$ define a relation $\Psi_{(t,T')}$ on $\cX_l(t_0,(s_l)_0)$ by
\[
g_1\mathrel{\Psi_{(t,T')}}g_2\iff(\tilde q_{lk}(t),T',q_{lk}(g_1),
q_{lk}(g_2))\in\bar{\mathcal F}.
\]

\emph{$\Psi_{(t,T')}$ is an equivalence relation.} It is symmetric
as $\bar{\mathcal F}$ is. Let $g_1,g_2,g_3$ be given. The image of
$(t,T',g_1,g_2,g_3)$ in $\Pi^3_k$ is the image of some
$((\tau,T'),h_1,h_2,h_3)$ of the limit: $t_k=\tilde q_{lk}(t)$ and
$q_k(h_a)=q_{lk}(g_a)$. So $g_a\mathrel\Psi g_b$ iff
$h_a\mathrel{\mathcal F_{(\tau,T')}}h_b$, for all $a,b$, and
$\mathcal F_{(\tau,T')}$ is reflexive and transitive.

\emph{Closure.} Let $\beta\colon u\to t$ be a one-step arrow of
$\cS_n(\cX_l)$, $T^\circ\in\Delta(T',q_{li}(\beta))$ and
$g_1\beta\mathrel{\Psi_{(u,T^\circ)}}g_2\beta$. The image of
$(u,t,\beta,T',T^\circ,g_1,g_2)\in\Pi'_l$ in $\Pi'_k$ is the
image of some $((\tau,T'),(\upsilon,\gamma,T^\circ),h_1,h_2)$ of the
limit. Then $(u_k,T^\circ,q_k(h_1\gamma),q_k(h_2\gamma))=(\tilde
q_{lk}(u),T^\circ,q_{lk}(g_1\beta),q_{lk}(g_2\beta))\in
\bar{\mathcal F}$, so $h_1\gamma\mathrel{\mathcal F_{(\upsilon,
T^\circ)}}h_2\gamma$, hence $h_1\mathrel{\mathcal F_{(\tau,T')}}h_2$
by \eqref{eq:sol}, and $g_1\mathrel{\Psi_{(t,T')}}g_2$.

\emph{Conclusion.} $E_{T'^t}\subseteq\Psi_{(t,T')}$ for all $t$ and
$T'$, by induction along the recursive definition of the pullback:
the children of the root of $T'^t$ have as labels one-step arrows
$\beta\colon u\to t$ and as subtrees trees $T^{\circ u}$ with
$T^\circ\in\Delta(T',q_{li}(\beta))\subseteq\mathcal T$; so
$E_{T'^t}$ is the join of relations $\beta^*E_{T^{\circ u}}\subseteq
\beta^*\Psi_{(u,T^\circ)}\subseteq\Psi_{(t,T')}$ (the equality if
there are no children). Now if $g\mathrel{\varepsilon^l_{(\tau,T')}}
g'$, then $q_l(g)\mathrel{\Psi_{(t_l,T')}}q_l(g')$, that is
$(t_k,T',q_k(g),q_k(g'))\in\bar{\mathcal F}$, and $g\mathrel{\mathcal
F_{(\tau,T')}}g'$.

For an $F$-solution the proof is the same, with $g_a\in\cS_n(\cX_{k'})
(t,s_{k'})$ in the sets $\Pi^3_{k'}$, $\Pi'_{k'}$: the elements of
the limits then have their $g_a$ in $K_\tau$ (the condition at the
indices $k'\ge k$ implies it at all indices, by
Lemma~\ref{lem:collapse}(c)), and $\Psi$ is defined on
$\cS_n(\cX_l)(t,s_l)$.

(c) By (a) and (b), and the description of the groups before the
theorem; recall $K_\sigma=H_\sigma$.
\end{proof}

The nodes are all the pairs $(\tau,T')$, not only those reached from
$(\sigma,T)$ by passing to children, and a node is a child of itself
by the identity. By Lemma~\ref{lem:tarrows}(b), if all threads are
bounded, the nodes reached from $(\sigma,T)$ and their children which
are not isomorphisms are those of the pullback of $T$ along $q_i\colon\cM\to\cX_i$, formed as in
Section~\ref{sec:pull} with the one-step arrows of $\cS_n(\cM)$: a
tree in which a node can have infinitely many children.

\subsection*{Bounded strings: closed solutions and the recursion}
Let $\Pi$ be the set of the triples $(x,g,g')$ of a node
$x=(\tau,T')$ and $g,g'\in V_\tau$. It is the limit of the finite
sets $\Pi_k$ of the tuples $(t,T',g,g')$ at the stages $k\ge i$
(defined as $\Pi^3_k$, with two arrows), a profinite space; let
$q_k\colon\Pi\to\Pi_k$ be the projections. A \emph{family} is a
subset $\mathcal U\subseteq\Pi$ such that every $\mathcal U_x=\{(g,g')
\mid(x,g,g')\in\mathcal U\}$ is an equivalence relation; it is of
level $k$ iff $\mathcal U=q_k^{-1}(q_k(\mathcal U))$, and of finite
level iff it is open and closed. So $\bar E$ is a closed family, and
every $\bar E_x$ is an intersection of open equivalence relations.
Let $\Pi'=\varprojlim\Pi'_k$ be the profinite space of the tuples
$(x,(\upsilon,\gamma,T''),g,g')$ of a node, a child of it and
$g,g'\in V_\tau$; the maps to $\Pi$ given by $(x,g,g')$ and by
$((\upsilon,T''),g\gamma,g'\gamma)$ are continuous.

\begin{thm}\label{thm:closed}
Suppose that there is $N$ such that every string of every $\cX_k$ has
at most $N$ objects; for instance, that $n_L$ is finite. Let
$\mathcal F$ be an $E$-solution such that
\begin{enumerate}
\item[(i)] $\mathcal F$ is closed in $\Pi$, and
\item[(ii)] every $\mathcal F_x$ is an intersection of open
equivalence relations.
\end{enumerate}
Then $\mathcal F$ is an intersection of $E$-solutions of finite
level. Hence $\bar E$ is the least $E$-solution with \emph{(i)} and
\emph{(ii)}. The same holds for the $F$-solutions and $\bar F$.
\end{thm}
\begin{proof}
A \emph{neighbourhood} is a family of finite level containing
$\mathcal F$; the intersection of two neighbourhoods is one.

\emph{Step 1: $\mathcal F$ is the intersection of its neighbourhoods.}
Let $(x,g,g')\notin\mathcal F$, $x=(\tau,T')$. By (ii) there is an
open equivalence relation $\mathcal E\supseteq\mathcal F_x$ on
$V_\tau$ with $(g,g')\notin\mathcal E$; let $k$ be such that
$\mathcal E$ is the inverse image of an equivalence relation on
$q_k(V_\tau)$, and $\mathcal E^+$ the equivalence relation on
$\cX_k((t_k)_0,(s_k)_0)$ obtained by adding the equality. For $l\ge
k$ let $N_l$ be the set of the nodes $(\upsilon,T')$ with $u_l=t_l$,
and $B_l$ the set of the $(y,h,h')\in\mathcal F$ with $y\in N_l$ and
$(q_k(h),q_k(h'))\notin\mathcal E^+$. The $B_l$ are closed, by (i),
and decrease, and a point of their intersection has $y=x$ (a thread
is determined by its components) and $(h,h')\in\mathcal
F_x\setminus\mathcal E$; so the intersection is empty and
$B_l=\emptyset$ for some $l$. The family which at the nodes of $N_l$
is the inverse image of $\mathcal E^+$ under $q_k$ and at the other
nodes the total relation is of level $l$, contains $\mathcal F$ and
does not contain $(x,g,g')$.

\emph{Step 2: for every neighbourhood $\mathcal U$ there is a
neighbourhood $\mathcal U'$ such that
\begin{equation}\label{eq:S}
g\gamma\mathrel{\mathcal U'_{(\upsilon,T'')}}g'\gamma\quad
\Longrightarrow\quad g\mathrel{\mathcal U_{(\tau,T')}}g'
\end{equation}
for all children $(\upsilon,\gamma,T'')$ of nodes $(\tau,T')$.} For a
neighbourhood $\mathcal U'$ let $A(\mathcal U')\subseteq\Pi'$ be the
set of the points with $((\upsilon,T''),g\gamma,g'\gamma)\in\mathcal
U'$ and $(x,g,g')\notin\mathcal U$. These sets are closed, the set of
them is directed downwards, and by Step~1 a point of their
intersection has $g\gamma\mathrel{\mathcal F_{(\upsilon,T'')}}
g'\gamma$, so $g\mathrel{\mathcal F_x}g'$ by \eqref{eq:sol},
contradicting $(x,g,g')\notin\mathcal U$. So some $A(\mathcal U')$ is
empty.

\emph{Step 3: every neighbourhood $\mathcal U$ contains an
$E$-solution of finite level containing $\mathcal F$.} By Step~2
choose neighbourhoods
\[
\mathcal U=\mathcal U_1\supseteq\mathcal U'_1\supseteq\mathcal U_2
\supseteq\mathcal U'_2\supseteq\dots\supseteq\mathcal U_N\supseteq
\mathcal U'_N
\]
such that \eqref{eq:S} holds for the pairs $(\mathcal U_m,\mathcal
U'_m)$, $1\le m\le N$, and $(\mathcal U'_m,\mathcal U_{m+1})$, $1\le
m<N$, in the place of $(\mathcal U,\mathcal U')$. Let $k$ be such
that all of them are of level $k$, and $l\ge k$ as in the fact in the
proof of Theorem~\ref{thm:level}(b) for the systems $\Pi^3$ and
$\Pi'$.

A \emph{stage node} is a pair $z=(t,T^\circ)$ of a string $t$ of
$\cX_l$ and $T^\circ\in\mathcal T$ at $\tilde q_{li}(t)$; let
$m(z)\le N$ be the number of objects of $t$ and $V_z=\cX_l(t_0,
(s_l)_0)$. For a stage node $z$ with $m(z)=m$ let
\[
\Omega_z=\{(g,g')\in V_z^2\mid q_{lk}(z,g,g')\in q_k(\mathcal U_m)\}.
\]
As in the proof of Theorem~\ref{thm:level}(b), $\Omega_z$ is an
equivalence relation: for $g_1,g_2,g_3\in V_z$ there is a point
$(x,h_1,h_2,h_3)$ of $\varprojlim\Pi^3_{k'}$ with the same image at
the stage $k$ as $(z,g_1,g_2,g_3)$, and then $g_a\mathrel{\Omega_z}
g_b$ iff $h_a\mathrel{(\mathcal U_m)_x}h_b$, because $\mathcal U_m$
is of level $k$. Let $\Xi_z$ be the equivalence relation on $V_z$
generated by the pairs
\begin{equation}\label{eq:gen}
\bigl(q_l(c)\kappa^{-1},\,q_l(c')\kappa^{-1}\bigr),
\end{equation}
where $y=(\upsilon,T'')$ is a node, $c\mathrel{(\mathcal U'_m)_y}c'$,
and $\kappa\colon u_l\to t$ is an isomorphism of strings of $\cX_l$
with $T''=q_{li}(\kappa)^*T^\circ$. Define the family $\mathcal H$ by
\[
g\mathrel{\mathcal H_{(\tau,T')}}g'\iff q_l(g)\mathrel{\Xi_{(t_l,T')}}
q_l(g').
\]
It is a family of level $l$.

(a) \emph{$\mathcal F\subseteq\mathcal H$}: for $g\mathrel{\mathcal
F_x}g'$ and $m=m(t_l,T')$ we have $g\mathrel{(\mathcal U'_m)_x}g'$,
and \eqref{eq:gen} with $y=x$, $\kappa=\id$.

(b) \emph{$\Xi_z\subseteq\Omega_z$; hence $\mathcal H_x\subseteq
(\mathcal U_m)_x$ for $m=m(t_l,T')$, and $\mathcal H\subseteq\mathcal
U$.} As $\Omega_z$ is an equivalence relation it suffices to consider
a pair \eqref{eq:gen}; put $g=q_l(c)\kappa^{-1}$, $g'=q_l(c')
\kappa^{-1}$. The tuple $(u_l,t,\kappa,T^\circ,T'',g,g')$ is in
$\Pi'_l$. Its image in $\Pi'_k$ is the image of a point of $\Pi'$: a
node $x^\circ=(\tau^\circ,T^\circ)$, a child $(\upsilon^\circ,\pi,
T'')$ of it and $h,h'\in V_{\tau^\circ}$, with $\tau^\circ_k=\tilde
q_{lk}(t)$, $\upsilon^\circ_k=u_k$, $q_k(h)=q_{lk}(g)$ and
$q_k(h\pi)=q_{lk}(g\kappa)=q_k(c)$, and the same for $h'$. So
$((\upsilon^\circ,T''),h\pi,h'\pi)$ and $(y,c,c')$ have the same
image in $\Pi_k$; as $\mathcal U'_m$ is of level $k$,
$h\pi\mathrel{(\mathcal U'_m)_{(\upsilon^\circ,T'')}}h'\pi$, and
$h\mathrel{(\mathcal U_m)_{x^\circ}}h'$ by \eqref{eq:S}. As
$q_k(x^\circ,h,h')=q_{lk}(z,g,g')$, this is $g\mathrel{\Omega_z}g'$.

(c) \emph{$\mathcal H$ is an $E$-solution.} Let $(\upsilon,\gamma,
T'')$ be a child of $x=(\tau,T')$ and $g\gamma\mathrel{\mathcal
H_{(\upsilon,T'')}}g'\gamma$; let $m$ be the number of objects of
$t_l$. If $\gamma_l$ is one-step, $u_l$ has $m+1\le N$ objects, so
$g\gamma\mathrel{(\mathcal U_{m+1})_{(\upsilon,T'')}}g'\gamma$ by
(b), hence $g\mathrel{(\mathcal U'_m)_x}g'$ by \eqref{eq:S} for
$(\mathcal U'_m,\mathcal U_{m+1})$, and $g\mathrel{\mathcal H_x}g'$
by \eqref{eq:gen} with $\kappa=\id$. If $\gamma_l$ is an isomorphism
of strings, then $T''=\gamma_i^*T'$, and there is a chain from
$q_l(g)\gamma_l$ to $q_l(g')\gamma_l$ in $V_{(u_l,T'')}$ whose
consecutive members are pairs \eqref{eq:gen} (or their reverses),
with isomorphisms $\kappa_a$ onto $u_l$ and nodes $(\upsilon_a,
T''_a)$, $T''_a=q_{li}(\kappa_a)^*T''$. Composing the members with
$\gamma_l^{-1}$ gives a chain from $q_l(g)$ to $q_l(g')$ whose
consecutive members are the pairs \eqref{eq:gen} for the
isomorphisms $\gamma_l\kappa_a$ onto $t_l$, and $T''_a=q_{li}(\gamma_l
\kappa_a)^*T'$. So $g\mathrel{\mathcal H_x}g'$.

By Steps 1 and 3, $\mathcal F$ is the intersection of the
$E$-solutions of finite level containing it. By
Theorem~\ref{thm:level} each of them contains $\bar E$, and $\bar E$
is an $E$-solution with (i) and (ii). For the $F$-solutions the proof
is the same, with the sets $K_\tau$ and $\cS_n(\cX_k)(t,s_k)$.
\end{proof}

For a set $(\rho_\lambda)$ of relations on a profinite set $V$ let
$\widehat\bigvee_\lambda\rho_\lambda$ be the intersection of the open
equivalence relations on $V$ containing all $\rho_\lambda$, as in
\cite[\S8]{Proftrees}; for the empty set of relations it is the
equality. Under the hypothesis of Theorem~\ref{thm:closed} all
threads are bounded, so the nodes are the pairs of a string of $\cM$
and a tree; the \emph{height} of a node $(\tau,T')$ is $N$ minus the
number of objects of the string $\tau$ of $\cM$, and by
Lemma~\ref{lem:tarrows}(b) a child $(\upsilon,\gamma,T'')$ for which
$\gamma$ is not an isomorphism has smaller height. Define relations
$\widehat E_x$ on $V_\tau$, for the nodes $x=(\tau,T')$, by recursion
on the height:
\[
\widehat E_x=\widehat\bigvee\gamma^*\widehat E_{(\upsilon,T'')}
\]
over the children $(\upsilon,\gamma,T'')$ of $x$ for which $\gamma$
is not an isomorphism, where $g\mathrel{\gamma^*\rho}g'$ iff
$g\gamma\mathrel\rho g'\gamma$; and relations $\widehat F_x$ on
$K_\tau$ in the same way.

\begin{thm}\label{thm:rec}
Under the hypothesis of Theorem~\ref{thm:closed}, $\bar E_x=\widehat
E_x$ and $\bar F_x=\widehat F_x$ for every node $x$. So
\[
S_{(i,T)}=\{\Theta\in G_\sigma\mid\id\mathrel{\widehat E_{(\sigma,T)}}
\Theta\},\qquad R_{(i,T)}=\{\Theta\in H_\sigma\mid\id\mathrel{
\widehat F_{(\sigma,T)}}\Theta\}.
\]
\end{thm}
\begin{proof}
We treat $E$; the proof for $F$ is the same. Let $\widehat E$ be the
family of the $\widehat E_x$.

(1) \emph{$\widehat E_x\subseteq\bar E_x$}, by induction on the
height: for every $k$, $\varepsilon^k_x$ is an open equivalence
relation which contains $\gamma^*\varepsilon^k_{(\upsilon,T'')}
\supseteq\gamma^*\bar E_{(\upsilon,T'')}\supseteq\gamma^*\widehat
E_{(\upsilon,T'')}$ for the children which are not isomorphisms.

(2) \emph{Transport.} Let $\iota\colon\upsilon\to\tau$ be an
isomorphism and $T'\in\mathcal T$ at $t_i$. Then for $h,h'\in
V_\tau$: $h\iota\mathrel{\widehat E_{(\upsilon,\iota_i^*T')}}h'\iota$
iff $h\mathrel{\widehat E_{(\tau,T')}}h'$. By induction on the
height, the two nodes having the same height: $(\omega,\gamma,T'')
\mapsto(\omega,\iota\gamma,T'')$ is a bijection from the children of
$(\upsilon,\iota_i^*T')$ to the children of $(\tau,T')$, since
$\Delta(\iota_i^*T',\gamma_i)=\Delta(T',\iota_i\gamma_i)$; it
preserves the children which are not isomorphisms; $h\mapsto h\iota$
is a homeomorphism $V_\tau\to V_\upsilon$; and
$(h\iota)\gamma=h(\iota\gamma)$.

(3) \emph{$\widehat E$ is an $E$-solution.} For a child with $\gamma$
not an isomorphism \eqref{eq:sol} holds by the definition of
$\widehat E_x$; for an isomorphism $\gamma$ we have
$T''=\gamma_i^*T'$ and (2).

(4) \emph{$\widehat E$ is closed in $\Pi$.} We show by induction on
$d$: if a net $(x_a,g_a,g'_a)$ in $\widehat E$ with nodes $x_a$ of
height at most $d$ converges to $(x,g,g')$ in $\Pi$, then
$g\mathrel{\widehat E_x}g'$. Suppose not, $x=(\tau,T')$. There is an
open equivalence relation $\mathcal E\supseteq\widehat E_x$ on
$V_\tau$ with $(g,g')\notin\mathcal E$; let $k$ and $\mathcal E^+$ be
as in Step~1 of the proof of Theorem~\ref{thm:closed}. For $a$ large,
$x_a=(\tau_a,T')$ with $(\tau_a)_k=t_k$; let $\mathcal E^a$ be the
inverse image of $\mathcal E^+$ under $q_k$ on $V_{\tau_a}$, an open
equivalence relation, and $(g_a,g'_a)\notin\mathcal E^a$ for $a$
large. As $g_a\mathrel{\widehat E_{x_a}}g'_a$, the relation $\mathcal
E^a$ does not contain $\widehat E_{x_a}$; so $x_a$ has a child
$(\upsilon_a,\gamma_a,T''_a)$ with $\gamma_a$ not an isomorphism (in
particular $d\ge1$) and $h_a,h'_a\in V_{\tau_a}$ with
\[
h_a\gamma_a\mathrel{\widehat E_{(\upsilon_a,T''_a)}}h'_a\gamma_a,
\qquad(h_a,h'_a)\notin\mathcal E^a.
\]
The points $(x_a,(\upsilon_a,\gamma_a,T''_a),h_a,h'_a)$ of the
compact space $\Pi'$ have a subnet converging to a point
$(x,(\upsilon,\gamma,T''),h,h')$; then $(h,h')\notin\mathcal E$, and
the points $((\upsilon_a,T''_a),h_a\gamma_a,h'_a\gamma_a)$ of
$\widehat E$, at nodes of height at most $d-1$, converge to
$((\upsilon,T''),h\gamma,h'\gamma)$. By the induction hypothesis
$h\gamma\mathrel{\widehat E_{(\upsilon,T'')}}h'\gamma$, and by (3)
$h\mathrel{\widehat E_x}h'$, so $(h,h')\in\mathcal E$: a
contradiction. Note that the limit $\gamma$ of the $\gamma_a$ can be
an isomorphism, and that the height of the limit node is not
restricted in the statement proved by induction.

(5) By (3), (4) and the definition, $\widehat E$ is an $E$-solution
with (i) and (ii) of Theorem~\ref{thm:closed}; so $\bar
E\subseteq\widehat E$, and with (1) $\bar E=\widehat E$.
\end{proof}

For the towers of \cite[\S7]{Proftrees} (functors which are the
identity on objects and full, so that they reflect isomorphisms and
the strings are bounded by the number of objects) this is
\cite[Prop.~8.3]{Proftrees}; there a child which is not an
isomorphism is not an isomorphism at any stage, and the children over
a label form a compact space. In general the children which are not
isomorphisms can converge to isomorphisms, as in
Example~\ref{ex:interval}.

\subsection*{Two examples}
In both examples the inverse system consists of cone categories
$\cX(P,\Phi)$ (Section~\ref{sec:depth}), and so does the limit, in
the following sense: the objects of $\cM$ other than $c$ form a poset
$P_\infty$, the arrows of $\cM$ are the automorphisms of $c$, the
inequalities of $P_\infty$ and the elements of the profinite sets
$\Phi(o)=\cM(o,c)$, on which $Z=\{1,\omega\}$ acts freely, and
composition with an inequality $o'<o$ is a restriction map
$\Phi(o)\to\Phi(o')$. Let $\sigma=(c)$.

\begin{lem}\label{lem:conesol}
For such $\cM$ let $(\approx_o)_{o\in P_\infty}$ be equivalence
relations on the sets $\Phi(o)$ such that $x|_{o'}\approx_{o'}
x'|_{o'}$ implies $x\approx_ox'$ for $o'<o$, and no $x$ is related to
$\omega x$. Then the family $\mathcal G$ with $\mathcal G_{(\tau,T')}$
the relation $\approx_{\tau_0}$ if $\tau_0\ne c$, and the equality on
$Z$ if $\tau_0=c$, is an $E$-solution, and
$(\id,\omega)\notin\mathcal G_{((c),T)}$.
\end{lem}
\begin{proof}
Let $\gamma\colon\upsilon\to\tau$ be a thread arrow. If $\tau_0\ne
c$, then $\gamma$ is an inequality $\upsilon_0\le\tau_0$, and
\eqref{eq:sol} is the first hypothesis. If $\tau_0=c$, then
$g,g'\in Z$ and $\gamma$ is in $Z$ or an element $x$ of some
$\Phi(o)$; and $gx\approx_og'x$ implies $g=g'$ by the second
hypothesis.
\end{proof}

\begin{ex}[the least solution]\label{ex:towersol}
Let $n=\infty$ and $(\cX_d)$ the tower of Section~\ref{sec:depth},
$i=1$ and $T$ the leg tree of the proof of Theorem~\ref{thm:tower};
so $\omega\in S_{(1,T)}$. The limit $\cM$ has the poset $P_\infty$ of
the $b_k$ ($k<\infty$), $b_\infty$, $a'$ and $z^k_j$
(Remark~\ref{rem:tower}(b)), with $\Phi(b_k)=Z\times\Gamma_{\le k}$,
the sets $\Phi(z^k_j)$, $\Phi(a')$ defined as in $\cX_d$, and
$\Phi(b_\infty)=Z\times\Gamma_\infty$, where $\Gamma_\infty=\{\beta\}
\cup\{e_k,e'_k\mid k<\infty\}\cup\{e_\infty\}$ is the one-point
compactification. Let $\approx_o$ be the equality for $o=a'$,
$z^k_j$, $b_\infty$; for $o=b_k$, $k\ge1$, the relation with the
classes $\{\theta e_{k-1},\theta e_k\}$, $\{\theta e'_{k-1},\theta
e'_k\}$ and singletons; and for $o=b_0$ the relation with the classes
$\{\theta\beta,\theta e_0,\omega\theta e'_0\}$. The hypotheses of
Lemma~\ref{lem:conesol} hold: for $z^k_j<b_k$ and $a'<b_0$ by the
definition of $\Phi(z^k_j)$, $\Phi(a')$; for $b_l<b_k$, $l>k$,
because no class of $\approx_{b_l}$ with more than one element is
contained in $Z\times\Gamma_{\le k}$; and for $z^l_j<b_k$, $l>k$,
because the restriction $\Phi(b_k)\to\Phi(z^l_j)$ is injective. So
there is an $E$-solution, by closed equivalence relations, which does
not contain $(\id,\omega)$ at the node $((c),T)$; it is not of finite
level. In particular the least $E$-solution does not give
$S_{(1,T)}$. A solution of level $d$ contains
$(e_d,e'_d)$ at the nodes at $b_d$, since $e_d$ and $e'_d$ have the
same image in $\cX_d$.
\end{ex}

\begin{ex}[closed equivalence relations]\label{ex:interval}
Let $n\ge3$. For $d\ge1$ let $N_d=2^d-1$ and $P_d$ the poset with the
elements $y$, $a'$, $z_1,\dots,z_{N_d}$ and the order $a'<y$,
$z_m<y$. Let $\Gamma_d=\{\beta,w_0,\dots,w_{N_d}\}$ and
\begin{align*}
\Phi_d(y)&=Z\times\Gamma_d,\qquad
\Phi_d(z_m)=Z\times\Gamma_d/\bigl((\theta,w_{m-1})\sim(\theta,w_m)
\bigr),\\
\Phi_d(a')&=Z\times\Gamma_d/\bigl((\theta,\beta)\sim(\theta,w_0),\
(\theta,w_{N_d})\sim(\omega\theta,\beta)\bigr),
\end{align*}
with the quotient maps as restrictions; the actions of $Z$ are free.
Let $\mathcal Y_d=\cX(P_d,\Phi_d)$, and $q_d\colon\mathcal Y_{d+1}\to
\mathcal Y_d$ the functor which is the identity on $c$, $y$, $a'$,
sends $z_{2m}$ to $z_m$ and $z_m$ to $y$ for $m$ odd, and sends
$\beta$ to $\beta$ and $w_m$ to $w_{\lfloor m/2\rfloor}$: the pair of
$z_{2m}$ is sent to the pair of $z_m$, the pair of $z_m$, $m$ odd, to
a pair of equal elements, and $w_0$, $w_{2N_d+1}$ to $w_0$,
$w_{N_d}$, so the maps $\Phi_{d+1}(o)\to\Phi_d(q_do)$ are defined,
equivariant and compatible with the restrictions. We think of $w_m$
as the word of length $d$ which is the binary expansion of $m$; $q_d$
deletes the last letter, and $z_m$ joins two consecutive words.
\[
\begin{tikzcd}[column sep=small,row sep=small]
z_1\ar[drr]&z_2\ar[dr]&\cdots&z_{N_d}\ar[dl]&a'\ar[dll]&\\
&&y\ar[rrr,"{\beta,w_0,\dots,w_{N_d}}"]&&&c\ar[loop right,"\omega"]
\end{tikzcd}
\]
\emph{The limit.} The poset $P_\infty$ has the elements $a'<y$ and
$z_r<y$ for the dyadic rationals $0<r<1$, where $z_r$ is $z_{r2^d}$
at the stages $d$ with $r2^d$ an integer and $y$ at the others. With
$\Gamma_\infty=\{\beta\}\sqcup\{0,1\}^{\mathbb N}$,
\[
\begin{gathered}
\Phi(y)=Z\times\Gamma_\infty,\qquad\Phi(z_r)=Z\times\Gamma_\infty/
(\theta r^-\sim\theta r^+),\\
\Phi(a')=Z\times\Gamma_\infty/(\theta\beta\sim\theta\bar0,\
\theta\bar1\sim\omega\theta\beta),
\end{gathered}
\]
where $r^-$, $r^+$ are the two binary expansions of $r$ and $\bar0$,
$\bar1$ the constant sequences. The arrows $z_r\to y$ are not
invertible, and they converge to the identity of $y$ when the
denominator of $r$ in lowest terms tends to infinity.

\emph{The groups.} Let $i=1$ and $T$ the leg tree of $\mathcal Y_1$
at $(c)$ whose root has the one child $\beta$, from $(y;\beta)$, with
the subtree $T_0$ whose root has the two children $a'<y$ and $z_1<y$,
without children. In the pullback of $T$ along $q_{d1}$ the node
$\beta$ has the subtree $T_0^{(y;\beta)}$, in which the arrow $a'<y$
lies over a child, and so does $z_m<y$ for $m=2^{d-1}$; for the other
$m$ the arrow $z_m<y$ is mapped to the identity of $(y;\beta)$, so it
is a collapsing child. Each $z_m$ gives $w_{m-1}\sim w_m$ and $a'$
gives $\beta\sim w_0$, $w_{N_d}\sim\omega\beta$; so
$\beta\sim\omega\beta$ at $(y;\beta)$ and $\omega\in S_{T^{(c)}}$,
for every $d$. Hence $S_{(1,T)}=Z$, and $R_{(1,T)}=1$ by the proof of
Lemma~\ref{lem:cone}: $(1,T)$ is an obstruction at $(c)$. (Already
$\HH_L\Fin^{\mathcal Y_d}$ is not a pretopos, by
Lemma~\ref{lem:cone}; the point of the example is the form of the
description.)

\emph{A solution by closed equivalence relations.} Let $v\colon
Z\times\Gamma_\infty\to\mathbb R/2\mathbb Z$ be the continuous map
\[
v(\theta,\beta)=[\theta],\qquad v(\theta,w)=[\theta]+
\textstyle\sum_{k\ge1}w_k2^{-k},\qquad[1]=0,\ [\omega]=1,
\]
$\approx_y$ its kernel pair, a closed equivalence relation whose
quotient is a circle, and $\approx_{a'}$, $\approx_{z_r}$ the
equality. The hypotheses of Lemma~\ref{lem:conesol} hold: $v(\theta
r^-)=v(\theta r^+)$, $v(\theta\beta)=v(\theta\bar0)$,
$v(\theta\bar1)=[\theta]+1=v(\omega\theta\beta)$, and $v(\omega
x)=v(x)+1$. So there is an $E$-solution by closed equivalence
relations which does not contain $(\id,\omega)$ at the node
$((c),T)$, while every $E$-solution of finite level does, by
Theorem~\ref{thm:level}. In the recursion of Theorem~\ref{thm:rec},
at the node of $(y;\beta)$ the relations of the children are the
equality, and $\widehat\bigvee$ of the pairs $(\theta r^-,\theta
r^+)$, $(\theta\beta,\theta\bar0)$, $(\theta\bar1,\omega\theta\beta)$
is the total relation: an open equivalence relation containing the
pairs $(\theta r^-,\theta r^+)$ identifies all $\theta w$, the
quotient of the Cantor set by these pairs being an interval, which is
connected. So $\widehat\bigvee$ cannot be replaced by the join in the
lattice of closed equivalence relations.
\end{ex}

\subsection*{Theories with finite sorts}
Let $\cC$ be a pretopos with finite sorts such that $\Mod(\cC)$ is an
EI-category, and $(\cX_\Lambda)_{\Lambda\in I}$ the inverse system of
Section~\ref{sec:theories}, whose limit $\cM$ is the category of the
standard models. For a string $\tilde s$ of standard models, a finite
language $\Lambda\in I$ and a leg tree $T$ of $\cX_\Lambda$ at
$\tilde s|_\Lambda$, Theorem~\ref{thm:level} describes the groups
$S_{(\Lambda,T)}$ and $R_{(\Lambda,T)}$ of Theorem~\ref{thm:A} by the
solutions of finite level, where the level is now a finite language
$\Lambda'\supseteq\Lambda$. If $n_L$ is finite, the nodes are the
pairs of a string of standard models and a tree of $\mathcal T$, the
children which are not isomorphisms are the homomorphisms of strings
$\tilde u\to\tilde t$ from a string with one model more whose
restriction to $\Lambda$ is the label of a child or an isomorphism of
strings, and by Theorem~\ref{thm:rec} the groups are given by the
recursion with the joins $\widehat\bigvee$ over this tree of strings
of models. This is Theorem~\ref{thm:F}(a),(b).

\subsection*{The criterion without leg trees}
By the proof of Theorem~\ref{thm:general} it suffices to consider the
trees $\Max(s_i,a)$ of the functors $W\colon\cS_n(\cX_i)\to\Fin$, and
for these the tree in $\cM$ is the tree of the elements of $W$. We
return to an inverse system $(\cX_k)_{k\in I}$ satisfying the
hypothesis of Theorem~\ref{thm:closed}, with limit $\cM$ and
$\cC=\varinjlim\Fin^{\cX_k}$; the threads are the strings of $\cM$.
Fix $i$, $W\colon\cS_n(\cX_i)\to\Fin$, a string $\tilde s$ of $\cM$
with thread $\sigma$ as the base, and $a\in W(s_i)$. For a string
$\tilde t$ of $\cM$ with thread $(t_k)$ put $W(\tilde t)=W(t_i)$, and
for a thread arrow $\gamma\colon\tilde u\to\tilde t$ let
$\gamma_*=W(\gamma_i)\colon W(\tilde u)\to W(\tilde t)$. Define
equivalence relations $E^W_{(\tilde t,y)}$ on $V_{\tilde
t}=\cM(\tilde t_0,\tilde s_0)$, for the strings $\tilde t$ of $\cM$
and $y\in W(\tilde t)$, by recursion on $N$ minus the number of
objects of $\tilde t$:
\begin{equation}\label{eq:EW}
E^W_{(\tilde t,y)}=\widehat\bigvee\gamma^*E^W_{(\tilde u,y')}
\end{equation}
over the pairs $(\gamma,y')$ of a one-step arrow $\gamma\colon\tilde
u\to\tilde t$ of $\cS_n(\cM)$ and $y'\in W(\tilde u)$ with
$\gamma_*y'=y$.

\begin{thm}\label{thm:elements}
Under the hypothesis of Theorem~\ref{thm:closed}:
\begin{enumerate}
\item[(a)] For $T=\Max(s_i,a)$ and every string $\tilde t$ of $\cM$
and $y\in W(\tilde t)$ such that $\Max(t_i,y)\in\mathcal T$:
$\bar E_{(\tilde t,\Max(t_i,y))}=E^W_{(\tilde t,y)}$. In particular
$S_{(i,\Max(s_i,a))}=\{\Theta\in G_{\tilde s}\mid\id\mathrel{
E^W_{(\tilde s,a)}}\Theta\}$.
\item[(b)] $\HH_L\cC$ is a pretopos if and only if for every $i$,
every $W\colon\cS_n(\cX_i)\to\Fin$, every string $\tilde s$ of $\cM$,
$a\in W(\tilde s)$ and $\Theta\in H_{\tilde s}$:
\[
\id\mathrel{E^W_{(\tilde s,a)}}\Theta\quad\Longrightarrow\quad
\Theta_*a=a .
\]
\end{enumerate}
\end{thm}
\begin{proof}
(a) Let $T=\Max(s_i,a)$. For $y\in W(t)$, $t$ a string of $\cX_i$,
the children of the root of $\Max(t,y)$ with a one-step label
$\delta\colon u\to t$ are the $(\delta,y')$ with $W(\delta)y'=y$,
with the subtrees $\Max(u,y')$; and $\delta^*\Max(t,y)=\Max(u,
W(\delta)^{-1}y)$ for an isomorphism of strings $\delta\colon u\to
t$, by the proof of (P1). So the trees of $\mathcal T$ are of the
form $\Max(t,y)$, and the children $(\tilde u,\gamma,T'')$ of a node
$(\tilde t,\Max(t_i,y))$ are those with $T''=\Max(u_i,y')$ for some
$y'\in W(\tilde u)$ with $\gamma_*y'=y$. By
Lemma~\ref{lem:tarrows}(b) the short thread arrows which are not
isomorphisms are the one-step arrows of $\cS_n(\cM)$. Hence the
relations $\widehat E_{(\tilde t,\Max(t_i,y))}$ of
Theorem~\ref{thm:rec} satisfy the recursion \eqref{eq:EW}, by
induction on the height, and they are the $\bar E$.

(b) By the proof of Theorem~\ref{thm:general}, $\HH_L\cC$ is a
pretopos iff \eqref{eq:W} holds for all $i$, $W$, threads $\sigma$ and
$a\in W(s_i)$: every element of $\rho_i(H_\sigma\cap
S_{(i,\Max(s_i,a))})$ fixes $a$. The threads are the strings of
$\cM$, $H_\sigma=H_{\tilde s}$ by Lemma~\ref{lem:threads}(c),
$\Theta_*=W(\rho_i(\Theta))$, and (a) applies.
\end{proof}

So for bounded strings the leg trees, their pullbacks and the groups
$R$ can be removed from the criterion: what remains is one family of
equivalence relations for each functor $W$ at a stage, defined by a
recursion over the elements of $W$ at the strings of $\cM$. In the
situation of Theorem~\ref{thm:A} with $n_L$ finite, the strings of
$\cM$ are the strings of standard models, the functors $W$ at the
stage $\Lambda$ are the objects of $\PP\HH_L\cC_\Lambda$ by (R9), the
set $W(\tilde t)$ is the value at $W$ of the model $\hat G_{\tilde t}$
(extended to the pretopos completion), and by (R8) at every stage the
thread arrows are the homomorphisms of these models. This is
Theorem~\ref{thm:F}(c).

\section{Questions}\label{sec:questions}

The two questions of the first version of this paper on the role of
the threads are answered by Theorem~\ref{thm:bounded} (an obstruction
never occurs at unbounded threads only) and by
Corollary~\ref{cor:tower} (the criterion for G\"odel--Dummett logic is
not equivalent to the criteria for all the logics $L_m$).
Section~\ref{sec:intrinsic} describes the groups of a leg tree in the
category of models, and for $n_L$ finite the criterion without leg
trees.

\begin{que}
(a) In Theorem~\ref{thm:F}(c) the finite languages enter only
through the objects $W$. Is it sufficient to test a smaller class of
objects, defined in terms of $\cC$?

(b) For $n_L=\infty$: do the nodes at strings of models suffice in
Theorem~\ref{thm:level}, and is there a description of the groups
which does not refer to the levels? The recursion of
Theorem~\ref{thm:rec} is then not well-founded in general, and its
least solution can be wrong (Example~\ref{ex:towersol}).
\end{que}

\section*{Declarations}
The author has no competing interests to declare.

\begin{thebibliography}{9}
\bibitem{AGM08} S.~Aguzzoli, B.~Gerla, V.~Marra, \emph{G\"odel algebras
free over finite distributive lattices}, Ann. Pure Appl. Logic 155
(2008), 183--193.
\bibitem{Carai} L.~Carai, \emph{Free algebras and coproducts in
varieties of G\"odel algebras}, J. Symb. Log., published online 2026,
doi:10.1017/jsl.2026.10194; arXiv:2406.05480.
\bibitem{Interm} R.~Shibusawa, \emph{Intermediate logics and
imaginaries: the translation into an intermediate logic commutes with
the pretopos completion of all finite propositional theories exactly
for the extensions of G\"odel--Dummett logic}, preprint (2026),
archived at DOI 10.5281/zenodo.24150610.
\bibitem{Strings} R.~Shibusawa, \emph{G\"odel--Dummett logic over a
finite category of models: the models of the translated theory are
strings of arrows, and a subgroup obstruction to commuting with the
imaginaries}, preprint (2026), archived at DOI
10.5281/zenodo.24327653.
\bibitem{Finsorts} R.~Shibusawa, \emph{Intermediate logics and the
imaginaries of a theory whose models are finite on every sort},
preprint (2026), archived at DOI 10.5281/zenodo.24356860.
\bibitem{Proftrees} R.~Shibusawa, \emph{Intermediate logics and the
imaginaries of a theory with finite sorts and finitely many sentences:
profinite leg trees}, preprint (2026), archived at DOI
10.5281/zenodo.24362902.
\end{thebibliography}

\end{document}
